% Leçon 1 — Vocabulaire et compréhension de texte : amour de la famille et la solidarité familiale — Chinois 1ère A — Cours signé OnBuch+
% Programme officiel MINESEC. Compiler avec : tectonic ZHA01-vocabulaire-amour-de-la-famille-et-la-solidarite-f.tex
\newcommand{\DOCMATIERE}{Chinois}
\newcommand{\DOCNIVEAU}{1ère A}
\newcommand{\DOCMODULE}{Module 1 : Vie familiale et intégration sociale}
\newcommand{\DOCLECON}{ZHA01}
\newcommand{\DOCTITRE}{Leçon 1 — Vocabulaire et compréhension de texte : amour de la famille et la solidarité familiale}
% OnBuch+ — préambule « cours signés » ENRICHI (compilé avec tectonic / XeLaTeX)
% Système visuel « Pop » d'OnBuch+ : fond crème (jamais blanc), bordures encre
% épaisses, ombres « dures » décalées, titres Archivo Black en capitales.
% Version MATHÉMATIQUES : graphes (pgfplots/tikz), boîtes pédagogiques
% (définition, propriété, méthode, attention, le savais-tu, exercice résolu,
% exercice, TP, bilan), page de garde et signature OnBuch+ ancrée en bas.
%
% Macros attendues dans main.tex AVANT \input{preamble} :
%   \DOCMATIERE  \DOCNIVEAU  \DOCMODULE  \DOCLECON (numéro)  \DOCTITRE
\documentclass[11pt]{article}

\usepackage{fontspec}
\usepackage[french]{babel} % français = langue principale ; le chinois via \zh{..} / textebox (xeCJK)
\usepackage{xeCJK}
\usepackage{pifont} % \ding{51} \ding{55} utilisés par les modèles
\usepackage[a4paper,margin=2cm,top=3.4cm,bottom=2.8cm,headheight=1.4cm,footskip=1.3cm]{geometry}
\usepackage{amsmath,amssymb,amsfonts}
\usepackage{mathtools} % \xmapsto, \coloneqq... souvent utilisés par les modèles
\usepackage{cancel} % simplifications barrées (bilans de forces)
% \abs{z} (module, valeur absolue) et \norm{u} (norme), utilisés par les modèles
\providecommand{\abs}{}\let\abs\relax\DeclarePairedDelimiter{\abs}{\lvert}{\rvert}
\providecommand{\norm}{}\let\norm\relax\DeclarePairedDelimiter{\norm}{\lVert}{\rVert}
\usepackage[table]{xcolor} % [table] : fournit aussi \rowcolors (alternance de lignes)
\usepackage{pagecolor}
\usepackage{tikz}
\usetikzlibrary{arrows.meta,positioning,calc,shapes.geometric,shapes.misc,shapes.arrows,decorations.pathmorphing,decorations.pathreplacing,decorations.markings,patterns,shadows,mindmap,trees,fit,backgrounds,matrix,intersections,angles,quotes}
\usepackage{pgfplots}
\usepackage[version=4]{mhchem} % équations nucléaires \ce{^{235}_{92}U}
\usepackage[european,siunitx]{circuitikz} % schémas électriques
\pgfplotsset{compat=1.18}
\usepgfplotslibrary{fillbetween,groupplots} % aires entre courbes (calcul intégral)
\usepackage{enumitem}
\usepackage{fancyhdr}
% [most] charge toutes les bibliothèques tcolorbox (listings, minted, raster,
% documentation...) et gonfle inutilement la mémoire TeX sur un document long
% (~300 boîtes) : on ne charge que ce qui est réellement utilisé.
\usepackage[skins,breakable]{tcolorbox}
\usepackage{graphicx}
\usepackage{colortbl}
\usepackage{booktabs}
\usepackage{tabularx}
\usepackage{diagbox} % case d'en-tête diagonale des tableaux à double entrée
% Colonnes extensibles centrée / gauche / droite (C, L, R), souvent utilisées par les modèles
\newcolumntype{C}{>{\centering\arraybackslash}X}
\newcolumntype{L}{>{\raggedright\arraybackslash}X}
\newcolumntype{R}{>{\raggedleft\arraybackslash}X}
\usepackage{array}
\usepackage{multicol}
\usepackage{multirow}
\usepackage{siunitx}
% parse-numbers=false : accepte aussi \SI{2,5 \times 10^{-3}}{...}
\sisetup{locale=FR,output-decimal-marker={,},group-minimum-digits=4,parse-numbers=false}
\DeclareSIUnit\ohm{\ensuremath{\symup{\Omega}}} % U+2126 absent de la police
\DeclareSIUnit\division{div} % réglages d'oscilloscope (ms/div, V/div)
\DeclareSIUnit\tr{tr} % tours (vitesse de rotation, ex. \SI{1500}{\tr\per\minute})
\DeclareSIUnit\molar{M} % concentration molaire (ex. \SI{3}{\molar}, \SI{1}{\micro\molar})
\DeclareSIUnit\an{an} % année (le modèle confond parfois avec \year, un compteur TeX) — ex. \SI{5}{\kilo\meter\per\an}
\DeclareSIUnit\FCFA{FCFA} % franc CFA (ex. \SI{500}{\FCFA\per\kilo\gram})
\DeclareSIUnit\degreeBrix{\textdegree Brix} % teneur en sucre d'un sirop/jus (réfractométrie)
\DeclareSIUnit\month{mois} % le modèle utilise parfois \month, un compteur TeX, comme unité de durée
\usepackage{qrcode}
% \degree : commande fréquemment utilisée par les modèles pour un angle ou une
% température en dehors de \SI{}{} (non fournie par les paquets chargés).
\providecommand{\degree}{^{\circ}}
% \qed (fin de démonstration), utilisé par les modèles sans amsthm
\providecommand{\qed}{\hfill$\square$}
% \barre : double barre des valeurs interdites dans un tableau de variations (tkz-tab)
\providecommand{\barre}{\ensuremath{\|}}
% environnement proof (amsthm non chargé) utilisé par les modèles
\ifdefined\proof\else\newenvironment{proof}[1][Démonstration]{\par\noindent\textit{#1.}\ }{\qed\par}\fi
\usepackage{lastpage}
\usepackage{hyperref}
\hypersetup{hidelinks}

% ============================================================
% Palette « Pop » OnBuch+ — mêmes hex que la classe P (plus_ui.dart)
% ============================================================
\definecolor{poporange}{HTML}{F59321}
\definecolor{popdark}{HTML}{E07A0C}
\definecolor{popcream}{HTML}{FFF6E8}
\definecolor{popink}{HTML}{1C1714}
\definecolor{poppurple}{HTML}{7A5AE0}
\definecolor{popgreen}{HTML}{2E9E6A}
\definecolor{poppink}{HTML}{E0457B}
\definecolor{popblue}{HTML}{2F7DE1}
\definecolor{popgold}{HTML}{C98A12}
\definecolor{popmuted}{HTML}{8A7F76}
\definecolor{popline}{HTML}{EADFD2}
\definecolor{popgray}{HTML}{8A7F76}
\colorlet{popgrey}{popgray}
% Alias tolérants (le modèle nomme parfois les couleurs différemment)
\colorlet{popwhite}{white}
\colorlet{popblack}{popink}
\colorlet{popred}{poppink}
\colorlet{popyellow}{popgold}
% Teintes claires (fonds de tableaux/encadrés) — le modèle nomme parfois
% les couleurs avec un suffixe L (ex. popblueL) sans les avoir définies.
\colorlet{poporangeL}{poporange!20}
\colorlet{popdarkL}{popdark!20}
\colorlet{poppurpleL}{poppurple!20}
\colorlet{popgreenL}{popgreen!20}
\colorlet{poppinkL}{poppink!20}
\colorlet{popblueL}{popblue!20}
\colorlet{popgoldL}{popgold!20}
\colorlet{popredL}{popred!20}
\colorlet{popgrayL}{popgray!20}
% Teintes claires pour fonds de tableaux / zones
\colorlet{popblueL}{popblue!10!white}
\colorlet{poporangeL}{poporange!14!white}
\colorlet{popgreenL}{popgreen!12!white}
\colorlet{poppurpleL}{poppurple!12!white}
\colorlet{poppinkL}{poppink!12!white}
\colorlet{popgoldL}{popgold!14!white}

\pagecolor{popcream}

% ============================================================
% Typographie
% ============================================================
\defaultfontfeatures{Ligatures=TeX}
\setmainfont{PlusJakartaSans-Medium.ttf}[Path=../../../fonts/,BoldFont=PlusJakartaSans-Bold.ttf]
% Symboles, API et emojis absents de la police principale : repli sur DejaVu Sans (caractères actifs)
\newfontface\symfont{DejaVuSans.ttf}[Path=../../../fonts/]
\makeatletter
\newcommand\symactive[1]{\catcode#1=13 \begingroup\lccode`\~=#1 \lowercase{\endgroup\def~}{{\symfont\char#1}}}
\newcommand\symblank[1]{\catcode#1=13 \begingroup\lccode`\~=#1 \lowercase{\endgroup\def~}{}}
\newcommand\symhyph[1]{\catcode#1=13 \begingroup\lccode`\~=#1 \lowercase{\endgroup\def~}{\mbox{-}}}
\newcount\symc
\newcommand\symrange[2]{\symc=#1 \@whilenum\symc<#2 \do{\expandafter\symactive\expandafter{\the\symc}\advance\symc by 1 }}
\newcommand\symblankrange[2]{\symc=#1 \@whilenum\symc<#2 \do{\expandafter\symblank\expandafter{\the\symc}\advance\symc by 1 }}
\makeatother
\symrange{"0250}{"02FF}\symactive{"2713}\symactive{"2714}\symactive{"2717}\symactive{"2718}\symactive{"2610}\symactive{"2611}\symactive{"2612}
\symrange{"2600}{"27BF}
\catcode"2705=13 \begingroup\lccode`\~="2705 \lowercase{\endgroup\def~}{{\symfont\char"2713}}\symblank{"26BD}\symblank{"26A1}
\symrange{"2460}{"257F}\symactive{"223C}\symactive{"2248}\symactive{"2260}
\symhyph{"2011}\symhyph{"2E1A}\symblankrange{"1F300}{"1FAFF}\symblank{"FE0F}
% Caractères chinois : WenQuanYi Zen Hei (sous-ensemble GB2312 + pinyin), gras/italique simulés
\setCJKmainfont{WenQuanYiZenHei-subset.ttf}[Path=../../../fonts/,AutoFakeBold=3,AutoFakeSlant=0.2]
\xeCJKsetup{CJKspace=false}
% Pinyin : voyelles à caron (ǎ ǐ ǒ ǔ ǖ ǘ ǚ ǜ) absentes de la police principale -> caractères actifs
\newfontface\pyfont{WenQuanYiZenHei-subset.ttf}[Path=../../../fonts/]
\newcommand\pyactive[1]{\catcode#1=13 \begingroup\lccode`\~=#1 \lowercase{\endgroup\def~}{{\pyfont\char#1}}}
\pyactive{"01CD}\pyactive{"01CE}\pyactive{"01CF}\pyactive{"01D0}\pyactive{"01D1}\pyactive{"01D2}\pyactive{"01D3}\pyactive{"01D4}\pyactive{"01D5}\pyactive{"01D6}\pyactive{"01D7}\pyactive{"01D8}\pyactive{"01D9}\pyactive{"01DA}\pyactive{"01DB}\pyactive{"01DC}
% Maths en sans-serif (Fira Math) avec chiffres et lettres droites de Plus
% Jakarta, pour que les formules se fondent dans le texte.
\usepackage[math-style=ISO,bold-style=ISO]{unicode-math}
\setmathfont{FiraMath-Regular.otf}[Path=../../../fonts/]
\setmathfont{PlusJakartaSans-Medium.ttf}[Path=../../../fonts/,range={up/num,up/{Latin,latin},\mathup}]
\usepackage{icomma} % virgule décimale sans espace en mode math
% Caractères mathématiques Unicode absents des polices de texte (Plus Jakarta,
% Archivo Black) : rendus par leur équivalent mathématique, en texte comme en
% formule — sinon ils s'affichent en carré vide (ex. « Arithmétique dans ℤ »).
\usepackage{newunicodechar}
\newunicodechar{ℝ}{\ensuremath{\mathbb{R}}}
\newunicodechar{ℤ}{\ensuremath{\mathbb{Z}}}
\newunicodechar{ℂ}{\ensuremath{\mathbb{C}}}
\newunicodechar{ℕ}{\ensuremath{\mathbb{N}}}
\newunicodechar{ℚ}{\ensuremath{\mathbb{Q}}}
\newunicodechar{𝔻}{\ensuremath{\mathbb{D}}}
\newunicodechar{α}{\ensuremath{\alpha}}
\newunicodechar{β}{\ensuremath{\beta}}
\newunicodechar{γ}{\ensuremath{\gamma}}
\newunicodechar{δ}{\ensuremath{\delta}}
\newunicodechar{ε}{\ensuremath{\varepsilon}}
\newunicodechar{θ}{\ensuremath{\theta}}
\newunicodechar{λ}{\ensuremath{\lambda}}
\newunicodechar{μ}{\ensuremath{\mu}}
\newunicodechar{σ}{\ensuremath{\sigma}}
\newunicodechar{τ}{\ensuremath{\tau}}
\newunicodechar{φ}{\ensuremath{\varphi}}
\newunicodechar{ψ}{\ensuremath{\psi}}
\newunicodechar{ω}{\ensuremath{\omega}}
\newunicodechar{Σ}{\ensuremath{\Sigma}}
% majuscules produites par \MakeUppercase dans les titres (α-aminés -> Α)
\newunicodechar{Α}{\ensuremath{\alpha}}
\newunicodechar{Β}{\ensuremath{\beta}}
\newunicodechar{Γ}{\ensuremath{\Gamma}}
\newunicodechar{Δ}{\ensuremath{\Delta}}
\newunicodechar{Ε}{\ensuremath{\varepsilon}}
\newunicodechar{Θ}{\ensuremath{\Theta}}
\newunicodechar{Λ}{\ensuremath{\Lambda}}
\newunicodechar{Μ}{\ensuremath{\mu}}
\newunicodechar{Τ}{\ensuremath{\tau}}
\newunicodechar{Φ}{\ensuremath{\Phi}}
\newunicodechar{Ψ}{\ensuremath{\Psi}}
\newunicodechar{Ω}{\ensuremath{\Omega}}
\newunicodechar{↦}{\ensuremath{\mapsto}}
\newunicodechar{∈}{\ensuremath{\in}}
\newunicodechar{∉}{\ensuremath{\notin}}
\newunicodechar{∧}{\ensuremath{\wedge}}
\newunicodechar{⇔}{\ensuremath{\Leftrightarrow}}
\newunicodechar{⟺}{\ensuremath{\Longleftrightarrow}}
\newunicodechar{⇒}{\ensuremath{\Rightarrow}}
\newunicodechar{⟹}{\ensuremath{\Longrightarrow}}
\newunicodechar{∘}{\ensuremath{\circ}}
\newunicodechar{∪}{\ensuremath{\cup}}
\newunicodechar{∩}{\ensuremath{\cap}}
\newunicodechar{≡}{\ensuremath{\equiv}}
\newunicodechar{⊂}{\ensuremath{\subset}}
\newunicodechar{⊥}{\ensuremath{\perp}}
\newunicodechar{∃}{\ensuremath{\exists}}
\newunicodechar{∀}{\ensuremath{\forall}}
\newunicodechar{∛}{\ensuremath{\sqrt[3]{\,}}}
\newunicodechar{∜}{\ensuremath{\sqrt[4]{\,}}}
\newunicodechar{⃗}{\ensuremath{{}^{\to}}}
\newunicodechar{ⁿ}{\ifmmode^{n}\else\textsuperscript{\ensuremath{n}}\fi}
\newunicodechar{ˣ}{\ifmmode^{x}\else\textsuperscript{\ensuremath{x}}\fi}
\newunicodechar{ᵇ}{\ifmmode^{b}\else\textsuperscript{\ensuremath{b}}\fi}
\newunicodechar{ʳ}{\ifmmode^{r}\else\textsuperscript{\ensuremath{r}}\fi}
\newunicodechar{ᵘ}{\ifmmode^{u}\else\textsuperscript{\ensuremath{u}}\fi}
\newunicodechar{ʸ}{\ifmmode^{y}\else\textsuperscript{\ensuremath{y}}\fi}
\newunicodechar{ᵃ}{\ifmmode^{a}\else\textsuperscript{\ensuremath{a}}\fi}
\newunicodechar{ᵉ}{\ifmmode^{e}\else\textsuperscript{\ensuremath{e}}\fi}
\newunicodechar{ᵏ}{\ifmmode^{k}\else\textsuperscript{\ensuremath{k}}\fi}
\newunicodechar{ᵖ}{\ifmmode^{p}\else\textsuperscript{\ensuremath{p}}\fi}
\newunicodechar{ᴺ}{\ifmmode^{N}\else\textsuperscript{\ensuremath{N}}\fi}
\newunicodechar{ᶜ}{\ifmmode^{c}\else\textsuperscript{\ensuremath{c}}\fi}
\newunicodechar{ʰ}{\ifmmode^{h}\else\textsuperscript{\ensuremath{h}}\fi}
\newunicodechar{ᵠ}{\ifmmode^{q}\else\textsuperscript{\ensuremath{q}}\fi}
\newunicodechar{ₙ}{\ifmmode_{n}\else\textsubscript{\ensuremath{n}}\fi}
\newunicodechar{ₐ}{\ifmmode_{a}\else\textsubscript{\ensuremath{a}}\fi}
\newunicodechar{ᵢ}{\ifmmode_{i}\else\textsubscript{\ensuremath{i}}\fi}
\newunicodechar{ᵣ}{\ifmmode_{r}\else\textsubscript{\ensuremath{r}}\fi}
\newunicodechar{ₖ}{\ifmmode_{k}\else\textsubscript{\ensuremath{k}}\fi}
\newunicodechar{ᵦ}{\ifmmode_{\beta}\else\textsubscript{\ensuremath{\beta}}\fi}
\newunicodechar{ₓ}{\ifmmode_{x}\else\textsubscript{\ensuremath{x}}\fi}
\newfontfamily\popdisplay{ArchivoBlack-Regular.ttf}[Path=../../../fonts/]
\newfontfamily\poplogo{SpaceGrotesk-Bold.ttf}[Path=../../../fonts/]
\newfontfamily\popsemi{PlusJakartaSans-SemiBold.ttf}[Path=../../../fonts/]
\newfontfamily\popxbold{PlusJakartaSans-ExtraBold.ttf}[Path=../../../fonts/]
\newfontfamily\popxboldls{PlusJakartaSans-ExtraBold.ttf}[Path=../../../fonts/,LetterSpace=3.0]

\setlist[enumerate]{leftmargin=1.7em,itemsep=5pt,topsep=4pt}
\setlist[itemize]{leftmargin=1.7em,itemsep=4pt,topsep=4pt,label={\color{poporange}\textbullet}}
\setlength{\parindent}{0pt}
\setlength{\parskip}{5pt}
\renewcommand{\arraystretch}{1.25}

% pgfplots : style Pop par défaut
\pgfplotsset{
  every axis/.append style={axis line style={popink,line width=1pt},tick style={popink},
    grid style={popline,line width=0.5pt},label style={font=\small},tick label style={font=\footnotesize}},
  popcurve/.style={poporange,line width=1.8pt,no markers,smooth},
}
% Même style disponible dans un \draw TikZ simple (hors pgfplots)
\tikzset{popcurve/.style={poporange,line width=1.8pt}}
\tikzset{popgrid/.style={popline,very thin}}
\colorlet{popcurve}{poporange} % le nom du style est parfois utilisé comme couleur

% ============================================================
% Pill / squiggle
% ============================================================
\newcommand{\poppill}[3]{%
  \tikz[baseline=-0.55ex]\node[draw=popink,line width=1.2pt,rounded corners=2.6pt,%
    fill=#2,inner xsep=6.5pt,inner ysep=3pt]{%
    \popxboldls\fontsize{7}{7}\selectfont\color{#3}\MakeUppercase{#1}};%
}
\newcommand{\popsquiggle}[1]{%
  \tikz[baseline=-0.4ex]{
    \draw[#1,line width=2.6pt,line cap=round]
      plot [smooth] coordinates {(0,0.09) (0.34,0.01) (0.68,0.21) (1.02,0.01) (1.36,0.10)};
  }%
}
% Mot-clé surligné (marqueur orange sous le texte)
\newcommand{\cle}[1]{{\popxbold\color{popdark}#1}}

% ============================================================
% En-tête / pied
% ============================================================
\pagestyle{fancy}
\fancyhf{}
\renewcommand{\headrulewidth}{0pt}
\renewcommand{\footrulewidth}{0pt}
\lhead{\raisebox{-0.15ex}{\poplogo\fontsize{13}{13}\selectfont\color{popink}OnBuch\color{poporange}+}}
\rhead{\poppill{\DOCMATIERE\ \textbullet\ \DOCNIVEAU\ \textbullet\ Leçon \DOCLECON}{white}{popink}}
\lfoot{\popxbold\fontsize{7.6}{7.6}\selectfont\color{popmuted}\MakeUppercase{Cours signé OnBuch+}\ \textbullet\ {\popsemi\DOCTITRE}}
\rfoot{\tikz[baseline=-0.6ex]\node[rounded corners=8.5pt,fill=popink,minimum height=17pt,minimum width=17pt,inner xsep=5pt,inner sep=0]{\popxbold\fontsize{8}{8}\selectfont\color{popcream}\thepage\,/\,\pageref*{LastPage}};}

% ============================================================
% Titres
% ============================================================
\newcommand{\chaptitre}[2]{%
  \begin{tcolorbox}[enhanced,colback=poporange,colframe=popink,boxrule=2.6pt,arc=11pt,
      drop shadow=popink,left=15pt,right=15pt,top=12pt,bottom=11pt,before skip=3pt,after skip=18pt]
    \poppill{#2}{popink}{popcream}\par\vspace{9pt}
    {\raggedright\hyphenpenalty=10000\exhyphenpenalty=10000\popdisplay\fontsize{22}{24}\selectfont\color{popcream}\MakeUppercase{#1}\par}
  \end{tcolorbox}
}
% Section numérotée : pastille orange avec le numéro + titre Archivo
\newcounter{popsec}
\newcommand{\coursec}[1]{%
  \stepcounter{popsec}\setcounter{popsub}{0}%
  \par\vspace{16pt}\noindent
  \begin{minipage}{\linewidth}
  \tikz[baseline=-0.7ex]\node[draw=popink,line width=1.6pt,fill=poporange,rounded corners=4pt,minimum size=22pt,inner sep=0]{\popdisplay\fontsize{12}{12}\selectfont\color{popcream}\thepopsec};%
  \hspace{8pt}{\popdisplay\fontsize{15}{17}\selectfont\color{popink}\MakeUppercase{#1}}\par
  \vspace{4pt}\noindent{\color{poporange}\rule{36pt}{3.6pt}}
  \end{minipage}\par\nopagebreak\vspace{8pt}
}
\newcounter{popsub}
\newcommand{\courssub}[1]{%
  \stepcounter{popsub}%
  \par\vspace{9pt}\noindent{\popxbold\fontsize{11.5}{13}\selectfont\color{popink}{\color{poporange}\thepopsec.\thepopsub}\hspace{6pt}#1}\par\nopagebreak\vspace{4pt}
}

% ============================================================
% Boîtes pédagogiques — même recette : fond blanc, bordure épaisse colorée,
% ombre dure encre, badge plein. Titre optionnel [#1].
% ============================================================
\tcbset{popbase/.style={breakable,enhanced,colback=white,boxrule=2.3pt,arc=9pt,
  shadow={3pt}{-3pt}{0pt}{popink},left=14pt,right=14pt,top=11pt,bottom=11pt,before skip=11pt,after skip=15pt,
  fontupper=\popsemi\fontsize{10.6}{14.5}\selectfont\color{popink},
  fontlower=\popsemi\fontsize{10.6}{14.5}\selectfont\color{popink}}}
% Coche dessinée (le glyphe ✓ n'existe pas dans les polices du document)
\newcommand{\popcheck}[1]{\tikz[baseline=-0.5ex]\draw[#1,line width=1.6pt,line cap=round,line join=round](-0.13,0.0)--(-0.04,-0.09)--(0.13,0.1);}
\providecommand{\coche}{\popcheck{popgreen}}
\providecommand{\resultat}[1]{\par\smallskip\noindent\fcolorbox{popgreen}{popgreenL}{\parbox{\dimexpr\linewidth-2\fboxsep-2\fboxrule\relax}{\textbf{Résultat.} #1}}\par\smallskip}
\newcommand{\popboxhead}[3]{\poppill{#1}{#2}{white}\ifx\relax#3\relax\else\hspace{6pt}{\popxbold\fontsize{10.6}{12}\selectfont\color{popink}#3}\fi\par\vspace{7pt}}

\newtcolorbox{aretenir}[1][]{popbase,colframe=popgreen,before upper={\popboxhead{À retenir}{popgreen}{#1}}}
% Texte en chinois (document de lecture, dialogue, extrait) : cadre
% d'étude. Un titre optionnel (source, auteur) s'affiche dans l'en-tête.
\newtcolorbox{textebox}[1][]{popbase,colframe=popblue,before upper={\popboxhead{文本 · Texte}{popblue}{#1}},after upper={}}
% Marques ✓ / ✗ (forme correcte / fautive) souvent utilisées par les modèles
\providecommand{\cross}{\textcolor{poppink}{\ensuremath{\times}}}
\providecommand{\xmark}{\cross}\providecommand{\crossmark}{\cross}\providecommand{\cmark}{\ensuremath{\checkmark}}
% Ligne de réponse / justification à compléter (inventée par les modèles)
\providecommand{\justification}[1]{\par\noindent\textit{Justification~:} \dotfill\par}
% \textipa (package tipa non chargé) : la police ne couvre pas l'API -> simple passage
\providecommand{\textipa}[1]{#1}
% Soulignement (ulem non chargé) : \uline, \ul
\providecommand{\uline}[1]{\underline{#1}}\providecommand{\ul}[1]{\underline{#1}}
% \glossaire{\item ...} inventée par les modèles : liste de glossaire
\providecommand{\glossaire}[1]{\begin{itemize}#1\end{itemize}}
% \alert (beamer) inventée par les modèles : texte mis en évidence
\providecommand{\alert}[1]{\textbf{\textcolor{poppink}{#1}}}
% variantes ASCII d'ordinaux écrites en macro
\providecommand{\iere}{\textsuperscript{re}}\providecommand{\ieme}{\textsuperscript{e}}\providecommand{\ere}{\textsuperscript{re}}\providecommand{\eme}{\textsuperscript{e}}\providecommand{\er}{\textsuperscript{er}}\providecommand{\ie}{\textsuperscript{e}}
% Ordinaux français écrits en macro par les modèles : 1\ière, 3\ième
\newcommand{\ière}{\textsuperscript{re}}\newcommand{\ième}{\textsuperscript{e}}
% \exemple (inventée par les modèles) : étiquette « Exemple : »
\providecommand{\exemple}{\textbf{Exemple~:}\ }
% \square utilisable aussi hors mode math (cases à cocher)
\AtBeginDocument{\let\squareMath\square\renewcommand{\square}{\ensuremath{\squareMath}}}
% Mot ou phrase en chinois à l'intérieur d'un paragraphe français
\newcommand{\zh}[1]{\ifmmode\text{#1}\else{#1}\fi}
\let\eng\zh \let\esp\zh \let\all\zh % alias tolérés
% flèches d'intonation inventées par les modèles
\providecommand{\textsearrow}{}\renewcommand{\textsearrow}{\ensuremath{\searrow}}\providecommand{\textnearrow}{}\renewcommand{\textnearrow}{\ensuremath{\nearrow}}
% \fr{..} : mot français (inventé par les modèles)
\providecommand{\fr}[1]{\textit{#1}}
% zhbox : encadré de texte chinois inventé par les modèles
\newenvironment{zhbox}{\begin{quote}}{\end{quote}}
% \vs : « versus » inventé par les modèles
\providecommand{\vs}{\textit{vs}\ }
% \blank : case à compléter inventée par les modèles
\providecommand{\blank}[1][]{\underline{\hspace{1.6em}}}
\newtcolorbox{exemplebox}[1][]{popbase,colframe=poporange,before upper={\popboxhead{Exemple}{poporange}{#1}}}
\newtcolorbox{definition}[1][]{popbase,colframe=popblue,before upper={\popboxhead{Définition}{popblue}{#1}}}
\newtcolorbox{propriete}[1][]{popbase,colframe=poppurple,before upper={\popboxhead{Propriété}{poppurple}{#1}}}
\newtcolorbox{methode}[1][]{popbase,colframe=popgold,before upper={\popboxhead{Méthode}{popgold}{#1}}}
\newtcolorbox{attention}[1][]{popbase,colframe=poppink,before upper={\popboxhead{Attention}{poppink}{#1}}}
\newtcolorbox{experience}[1][]{popbase,colframe=popblue,colback=popblueL,before upper={\popboxhead{Expérience}{popblue}{#1}}}
% « Le savais-tu ? » : carte violette pleine (contexte, culture, Cameroun)
\newtcolorbox{savaistu}[1][]{popbase,colback=poppurple,colframe=popink,
  fontupper=\popsemi\fontsize{10.4}{14.2}\selectfont\color{white},
  before upper={\poppill{Le savais-tu\,?}{white}{poppurple}\ifx\relax#1\relax\else\hspace{6pt}{\popxbold\color{white}#1}\fi\par\vspace{7pt}}}
% Exercice résolu : énoncé en haut, \tcblower puis la solution sur fond crème
\newcounter{popexo}\newcounter{popexr}
\newtcolorbox{exoresolu}[1][]{popbase,colframe=poporange,colbacklower=popcream,
  segmentation style={popink,line width=1.2pt,dashed},
  before upper={\stepcounter{popexr}\popboxhead{Exercice résolu \thepopexr}{poporange}{#1}},
  before lower={\poppill{Solution}{popink}{popcream}\par\vspace{6pt}}}
% Exercice (non résolu en place) — la correction est dans la partie Corrigés
\newtcolorbox{exercice}[1][]{popbase,colframe=popink,
  before upper={\stepcounter{popexo}\popboxhead{Exercice \thepopexo}{popink}{#1}}}
% Corrigé
\newtcolorbox{corrige}[1][]{popbase,colframe=popgreen,colback=popgreenL,
  before upper={\popboxhead{Corrigé}{popgreen}{#1}}}
% Objectifs / prérequis / bilan
\newtcolorbox{objectifs}{popbase,colframe=popink,colback=poporangeL,
  before upper={\popboxhead{Objectifs de la leçon}{popink}{}}}
\newtcolorbox{prerequis}{popbase,colframe=popink,colback=popblueL,
  before upper={\popboxhead{Prérequis}{popblue}{}}}
\newtcolorbox{bilan}{popbase,colframe=popink,colback=white,boxrule=2.8pt,
  before upper={\popboxhead{Fiche bilan}{poporange}{L'essentiel à connaître pour l'examen}}}
% Figure encadrée avec légende
\newtcolorbox{popfigure}[1][]{enhanced,breakable=false,colback=white,colframe=popline,boxrule=1.4pt,arc=8pt,
  left=10pt,right=10pt,top=10pt,bottom=8pt,before skip=10pt,after skip=12pt,halign=center,
  lower separated=false,halign lower=center,
  fontlower=\popsemi\fontsize{9}{11.5}\selectfont\color{popmuted},
  #1}
\newcommand{\legende}[1]{\par\vspace{6pt}{\popsemi\fontsize{9}{11.5}\selectfont\color{popmuted}#1}}

% ============================================================
% Signature OnBuch+
% ============================================================
\newcommand{\poplogoinline}[1]{{\poplogo\fontsize{#1}{#1}\selectfont\color{popink}OnBuch\color{poporange}+}}
\newcommand{\signatureonbuch}{%
  \begin{tcolorbox}[enhanced,colback=popink,colframe=popink,boxrule=2.2pt,arc=10pt,
      drop shadow=poporange,left=14pt,right=14pt,top=10pt,bottom=10pt,before skip=0pt,after skip=0pt]
    \tikz[baseline=-0.7ex]\node[circle,fill=popgreen,draw=popcream,line width=1.3pt,minimum size=22pt,inner sep=0]{\popcheck{white}};%
    \hspace{9pt}\begin{minipage}[c]{0.62\linewidth}
      {\popxbold\fontsize{12}{13}\selectfont\color{popcream}Signé\ \textbullet\ {\poplogo OnBuch\color{poporange}+}}\par\vspace{2pt}
      {\raggedright\popsemi\fontsize{8}{10.5}\selectfont\color{popline}\MakeUppercase{Équipe pédagogique OnBuch+}\\\MakeUppercase{Conforme au programme officiel MINESEC}\par}
    \end{minipage}\hfill
    {\poplogo\fontsize{22}{22}\selectfont\color{popcream}OnBuch\color{poporange}+}
  \end{tcolorbox}%
}

% ============================================================
% Page de garde
% #1 titre · #2 matière · #3 niveau · #4 module · #5 n° leçon · #6 durée ·
% #7 description / objectifs (texte ou itemize)
% ============================================================
\newcommand{\pagegarde}[7]{%
  \thispagestyle{empty}
  \newgeometry{a4paper,margin=1.7cm,top=1.6cm,bottom=1.4cm}%
  \begin{tikzpicture}[remember picture,overlay]
    \fill[poporange!18] ([xshift=-3.2cm,yshift=-3.6cm]current page.north east) circle (4.6cm);
    \fill[poppurple!14] ([xshift=2.4cm,yshift=7.6cm]current page.south west) circle (3.8cm);
  \end{tikzpicture}%
  {\poplogo\fontsize{30}{30}\selectfont\color{popink}OnBuch\color{poporange}+}\hfill
  \raisebox{0.6ex}{\poppill{Cours signé \textbullet\ Premium}{popink}{popcream}}\par
  \vspace{1pt}\popsquiggle{poppurple}\par
  \vspace{8pt}
  \begin{tcolorbox}[enhanced,colback=poporange,colframe=popink,boxrule=2.8pt,arc=13pt,
      drop shadow=popink,left=18pt,right=18pt,top=12pt,bottom=13pt,after skip=10pt]
    \poppill{#2 \textbullet\ #3}{popink}{popcream}\hspace{5pt}\poppill{#4}{white}{popink}\par\vspace{8pt}
    {\popdisplay\fontsize{14}{14}\selectfont\color{popink}LEÇON #5}\par\vspace{4pt}
    {\raggedright\hyphenpenalty=10000\exhyphenpenalty=10000\popdisplay\fontsize{25}{27}\selectfont\color{popcream}\MakeUppercase{#1}\par}
    \vspace{8pt}
    {\popxbold\fontsize{9.5}{9.5}\selectfont\color{popink}\MakeUppercase{Durée conseillée : #6}}
  \end{tcolorbox}
  \begin{tcolorbox}[enhanced,colback=white,colframe=popink,boxrule=2.2pt,arc=10pt,
      drop shadow=popink,left=16pt,right=16pt,top=10pt,bottom=10pt,after skip=8pt]
    \poppill{Dans cette leçon}{poppurple}{white}\par\vspace{5pt}
    \popsemi\fontsize{9.3}{12.6}\selectfont\color{popink}#7
  \end{tcolorbox}
  \noindent\begin{minipage}[t]{0.487\linewidth}
    \begin{tcolorbox}[enhanced,colback=popgreen,colframe=popink,boxrule=2.2pt,arc=10pt,
        drop shadow=popink,left=11pt,right=11pt,top=10pt,bottom=10pt,height=3.1cm,valign=center]
      \tcbox[colback=white,colframe=popink,boxrule=1.3pt,arc=5pt,boxsep=0pt,nobeforeafter,
        left=3pt,right=3pt,top=3pt,bottom=3pt,valign=center]{\qrcode[height=1.9cm]{https://play.google.com/store/apps/details?id=com.luvvix.onbuch&hl=fr}}%
      \hspace{8pt}\begin{minipage}[c]{0.5\linewidth}
        {\popdisplay\fontsize{11}{12.5}\selectfont\color{white}\MakeUppercase{Télécharge OnBuch}}\par\vspace{4pt}
        {\raggedright\popsemi\fontsize{8.4}{11}\selectfont\color{white}Gratuit \textbullet\ Google Play\\Tuteur IA, annales, concours, résultats\par}
      \end{minipage}
    \end{tcolorbox}
  \end{minipage}\hfill
  \begin{minipage}[t]{0.487\linewidth}
    \begin{tcolorbox}[enhanced,colback=poppurple,colframe=popink,boxrule=2.2pt,arc=10pt,
        drop shadow=popink,left=11pt,right=11pt,top=10pt,bottom=10pt,height=3.1cm,valign=center]
      \tikz[baseline=-0.7ex]\node[circle,fill=white,draw=popink,line width=1.3pt,minimum size=1.6cm,inner sep=0]{\popdisplay\fontsize{24}{24}\selectfont\color{poppurple}+};%
      \hspace{8pt}\begin{minipage}[c]{0.6\linewidth}
        {\popdisplay\fontsize{11}{12.5}\selectfont\color{white}\MakeUppercase{Passe à OnBuch+}}\par\vspace{4pt}
        {\raggedright\popsemi\fontsize{8.4}{11}\selectfont\color{white}Tous les cours signés, Tuteur IA illimité, fiches PDF — onglet OnBuch+ de l'app\par}
      \end{minipage}
    \end{tcolorbox}
  \end{minipage}
  \vspace{8pt}
  \popcontenu
  \vfill
  \signatureonbuch
  \restoregeometry
  \newpage
}

% Rangée « Ce cours contient » (page de garde)
\newcommand{\poptile}[3]{% #1 couleur #2 gros texte #3 légende
  \begin{tcolorbox}[enhanced,colback=#1,colframe=popink,boxrule=2pt,arc=9pt,shadow={2.5pt}{-2.5pt}{0pt}{popink},
      left=6pt,right=6pt,top=9pt,bottom=9pt,halign=center,valign=center,height=1.75cm,width=\linewidth,nobeforeafter]
    {\popdisplay\fontsize{12.5}{13}\selectfont\color{white}\MakeUppercase{#2}}\par\vspace{3pt}
    {\centering\popsemi\fontsize{7.8}{9.5}\selectfont\color{white}#3\par}
  \end{tcolorbox}}
\newcommand{\popcontenu}{%
  \par\noindent{\popxboldls\fontsize{7.5}{7.5}\selectfont\color{popmuted}CE COURS SIGNÉ CONTIENT}\par\vspace{7pt}
  \noindent\begin{minipage}[t]{0.235\linewidth}\poptile{popblue}{Cours}{complet, illustré, pas à pas}\end{minipage}\hfill
  \begin{minipage}[t]{0.235\linewidth}\poptile{poporange}{Méthodes}{et exercices résolus}\end{minipage}\hfill
  \begin{minipage}[t]{0.235\linewidth}\poptile{poppink}{Exercices}{type examen + corrigés}\end{minipage}\hfill
  \begin{minipage}[t]{0.235\linewidth}\poptile{popgreen}{Fiche bilan}{l'essentiel à retenir}\end{minipage}\par}

% Sommaire
\newtcolorbox{sommaire}{enhanced,colback=white,colframe=popink,boxrule=2.2pt,arc=10pt,
  drop shadow=popink,left=18pt,right=18pt,top=13pt,bottom=13pt,before skip=6pt,after skip=20pt,
  before upper={\poppill{Sommaire}{poppurple}{white}\par\vspace{9pt}\popsemi\fontsize{10.6}{14.5}\selectfont\color{popink}}}

% Signature de fin de cours — poussée tout en bas de la dernière page
\newcommand{\findecours}{%
  \par\vfill\nopagebreak
  \begin{center}\popsquiggle{poporange}\end{center}\vspace{4pt}
  \signatureonbuch
}
\AtBeginDocument{\def\llbracket{\mathopen{[\mkern-3.5mu[}}\def\rrbracket{\mathclose{]\mkern-3.5mu]}}} % intervalles d'entiers (glyphe ⟦ absent de la police)
% Glyphes absents de Fira Math : redéfinis à partir de glyphes disponibles
\AtBeginDocument{%
  \def\setminus{\mathbin{\backslash}}\let\smallsetminus\setminus
  \def\vdots{\mathord{\vbox{\baselineskip3.2pt\lineskiplimit0pt\kern4pt\hbox{.}\hbox{.}\hbox{.}}}}%
  \def\ddots{\mathinner{\mkern1mu\raise6.5pt\hbox{.}\mkern2mu\raise3.6pt\hbox{.}\mkern2mu\raise0.7pt\hbox{.}\mkern1mu}}%
  \def\triangle{\mathord{\scalebox{0.9}{$\symup{\Delta}$}}}%
  \def\nabla{\mathord{\raisebox{\depth}{\scalebox{1}[-1]{$\symup{\Delta}$}}}}%
  \def\checkmark{\popcheck{popdark}}%
  \def\star{\mathbin{\ast}}\def\bigstar{\mathord{\ast}}%
  \def\diamond{\mathbin{\scalebox{0.6}{$\Diamond$}}}%
  \def\bigcup{\mathop{\mathchoice{\raisebox{-0.3em}{\scalebox{1.7}{$\cup$}}}{\scalebox{1.25}{$\cup$}}{\cup}{\cup}}\displaylimits}%
  \def\bigcap{\mathop{\mathchoice{\raisebox{-0.3em}{\scalebox{1.7}{$\cap$}}}{\scalebox{1.25}{$\cap$}}{\cap}{\cap}}\displaylimits}%
  \def\lesseqqgtr{\mathrel{\lessgtr}}\def\bigoplus{\mathop{\oplus}\displaylimits}%
}
\providecommand{\ssi}{si et seulement si} % abréviation fréquente
\AtBeginDocument{\providecommand{\wideparen}{\overparen}} % arc de cercle

%% FIGFIX-BEGIN v4 — correctifs globaux des figures (docs/onbuch-generation/tools/figfix)
\usepackage{adjustbox}\usepackage{etoolbox}
% toute figure TikZ/pgfplots plus large que la ligne est réduite à la largeur disponible
\BeforeBeginEnvironment{tikzpicture}{\begin{adjustbox}{max width=\linewidth}}
\AfterEndEnvironment{tikzpicture}{\end{adjustbox}}
% légendes pgfplots lisibles par-dessus les courbes
\pgfplotsset{every axis/.append style={legend style={fill=white,fill opacity=0.92,text opacity=1,draw=black!15,inner sep=3pt}}}
% étiquettes d'arêtes (syntaxe edge["..."]) avec fond blanc
\tikzset{every edge quotes/.append style={fill=white,inner sep=1.2pt}}
%% FIGFIX-END
\begin{document}

\pagegarde{Leçon 1 — Vocabulaire et compréhension de texte : amour de la famille et la solidarité familiale}{Chinois}{1ère A}{Module 1 : Vie familiale et intégration sociale}{ZHA01}{2 h}{Cette leçon fait découvrir aux élèves de 1ère A le vocabulaire chinois de l'amour familial et de la solidarité à travers un texte support original. Elle développe la lecture, la prononciation, l'écriture des caractères clés et la compréhension de texte.
\vspace{4pt}
\begin{multicols}{2}\small
\begin{itemize}
\item À la fin de cette leçon, je sais lire à voix haute un court texte chinois sur la famille en respectant les tons.
\item À la fin de cette leçon, je sais reconnaître et traduire au moins 15 mots de vocabulaire liés à l'amour familial et à la solidarité.
\item À la fin de cette leçon, je sais identifier les radicaux (clés) et tracer l'ordre des traits de caractères comme 爱, 家, 帮.
\item À la fin de cette leçon, je sais répondre à des questions de compréhension sur un texte chinois simple.
\item À la fin de cette leçon, je sais employer les collocations usuelles : 爱家人, 互相帮助, 关心父母.
\item À la fin de cette leçon, je sais comparer brièvement la solidarité familiale au Cameroun et en Chine.
\end{itemize}\end{multicols}}

\begin{sommaire}
\begin{multicols}{2}
\begin{enumerate}
\item Mise en route et découverte du thème
\item Texte support : lecture et traduction
\item Glossaire du thème
\item Clés (radicaux) et ordre des traits
\item Compréhension du texte
\item Collocations et production guidée
\item Activité d'intégration
\item Méthodes et exercices résolus
\item Exercices
\item Corrigés des exercices
\item Fiche bilan
\end{enumerate}
\end{multicols}
\end{sommaire}

\begin{objectifs}
\begin{itemize}
\item À la fin de cette leçon, je sais lire à voix haute un court texte chinois sur la famille en respectant les tons.
\item À la fin de cette leçon, je sais reconnaître et traduire au moins 15 mots de vocabulaire liés à l'amour familial et à la solidarité.
\item À la fin de cette leçon, je sais identifier les radicaux (clés) et tracer l'ordre des traits de caractères comme 爱, 家, 帮.
\item À la fin de cette leçon, je sais répondre à des questions de compréhension sur un texte chinois simple.
\item À la fin de cette leçon, je sais employer les collocations usuelles : 爱家人, 互相帮助, 关心父母.
\item À la fin de cette leçon, je sais comparer brièvement la solidarité familiale au Cameroun et en Chine.
\item À la fin de cette leçon, je sais rédiger 3 à 4 phrases guidées sur ma propre famille en utilisant le vocabulaire de la leçon.
\end{itemize}
\end{objectifs}

\begin{prerequis}
\begin{itemize}
\item Les membres de la famille (爸爸, 妈妈, 哥哥, 姐姐...) vus en Seconde
\item La lecture du pinyin et les quatre tons
\item Les structures de base : sujet + verbe + complément (我爱我家)
\item Les radicaux élémentaires (人, 女, 子, 口) et l'ordre général des traits
\item Les pronoms personnels et possessifs (我, 你, 他, 我的, 我们的)
\end{itemize}
\end{prerequis}

\newpage

\coursec{Mise en route et découverte du thème}

\courssub{Situation d'accroche : une valeur partagée entre le Cameroun et la Chine}

Imaginez la scène : à Yaoundé comme à Douala, quand un membre de la famille tombe malade, toute la famille se mobilise. Les oncles envoient de l'argent, les tantes préparent à manger, les cousins rendent visite, les grands-parents prient et conseillent. Personne ne reste seul face à la difficulté. Cette solidarité familiale, si naturelle au Cameroun, est aussi une valeur fondamentale en Chine, où l'on dit : \zh{家和万事兴} (\emph{jiā hé wàn shì xīng}) : « quand la famille est en harmonie, tout prospère ».

Aujourd'hui, nous allons lire le récit d'une famille chinoise et découvrir le vocabulaire qui permet de parler de l'amour et de la solidarité en famille. À la fin de cette leçon, vous serez capables de présenter votre propre famille en chinois et de raconter un geste de solidarité que vous avez vécu ou observé.

\textbf{Problématique de la leçon :} comment le chinois exprime-t-il l'amour et l'entraide au sein de la famille, et quels mots faut-il connaître pour en parler ?

\courssub{Brainstorming : les deux mots-clés du thème}

Avant toute lecture, posons-nous une question simple : comment dit-on « famille » et « amour » en chinois ?

\begin{definition}[Les deux mots-clés du thème]
\zh{家} (\emph{jiā}) : la maison, le foyer, la famille. Ce caractère représente à l'origine un toit (\zh{宀}) sous lequel se trouve un cochon (\zh{豕}) : autrefois, avoir du bétail sous son toit, c'était avoir un foyer.\par
\zh{爱} (\emph{ài}) : aimer, l'amour. C'est un verbe (\zh{我爱我家}, \emph{wǒ ài wǒ jiā}, « j'aime ma famille ») et un nom (\zh{爱}, « l'amour »). Remarquez le cœur (\zh{心}) caché dans sa forme traditionnelle \emph{ài} : l'amour vient du cœur.
\end{definition}

En combinant ces deux mots, on obtient le titre de notre thème :

\begin{aretenir}[Le thème du module]
\zh{家庭的爱与团结}\par
\emph{jiātíng de ài yǔ tuánjié}\par
« L'amour et la solidarité (l'union) de la famille »
\begin{itemize}
\item \zh{家庭} (\emph{jiātíng}) : la famille, le foyer (nom) ;
\item \zh{的} (\emph{de}) : mot-outil de possession et de qualification (« de ») ;
\item \zh{与} (\emph{yǔ}) : et (conjonction, plus soutenue et plus écrite que \zh{和}, \emph{hé}) ;
\item \zh{团结} (\emph{tuánjié}) : l'union, la solidarité ; s'unir.
\end{itemize}
\end{aretenir}

\begin{attention}[Piège des francophones : 爱 ne se conjugue pas]
En français, « aimer » se conjugue (j'aime, tu aimes, il aime). En chinois, \zh{爱} (\emph{ài}) est invariable : \zh{我爱} (j'aime), \zh{他爱} (il aime), \zh{我们爱} (nous aimons). N'ajoutez jamais de terminaison au verbe chinois. De même, attention au ton : \zh{爱} est au quatrième ton (\emph{ài}, ton descendant), à ne pas confondre avec \zh{矮} (\emph{ǎi}, « petit de taille », troisième ton). Dire \zh{我很矮} (\emph{Wǒ hěn ǎi}, « je suis petit ») au lieu de \zh{我很爱} (\emph{Wǒ hěn ài}, « j'aime beaucoup ») change complètement le sens !
\end{attention}

\courssub{Carte mentale : autour de 家}

Observons la carte mentale ci-dessous. Elle organise les quatre grandes idées que nous allons rencontrer dans le texte : l'amour, l'aide, la sollicitude et l'union.

\begin{popfigure}
\begin{tikzpicture}[mindmap, grow cyclic, every node/.style={concept, align=center}, concept color=popblue!40, level 1/.append style={level distance=3.2cm, sibling angle=90}, text=popdark]
\node[concept, concept color=poporange!50, font=\Large, align=center]{家\\ \emph{jiā}\\ famille}
  child[concept color=popgreen!40] { node[concept, align=center]{爱\\ \emph{ài}\\ amour} }
  child[concept color=popblue!40] { node[concept, align=center]{帮助\\ \emph{bāngzhù}\\ aide} }
  child[concept color=poppink!40] { node[concept, align=center]{关心\\ \emph{guānxīn}\\ sollicitude} }
  child[concept color=popgold!50] { node[concept, align=center]{团结\\ \emph{tuánjié}\\ union} };
\end{tikzpicture}
\legende{Carte mentale du thème : autour de 家 (jiā), la famille.}
\end{popfigure}

Lisez chaque branche à voix haute en respectant les tons : \emph{ài} (4\up{e} ton, descendant), \emph{bāngzhù} (1\up{er} ton + 4\up{e} ton), \emph{guānxīn} (deux 1\up{ers} tons, hauts et plans), \emph{tuánjié} (2\up{e} ton + 2\up{e} ton, montants).

\courssub{Un proverbe à mémoriser}

\begin{textebox}[Proverbe chinois : 家和万事兴]
家和万事兴。\\
\emph{jiā hé wàn shì xīng}\\
« Famille en harmonie, dix mille affaires prospèrent », c'est-à-dire : « quand la famille est en harmonie, tout réussit ».\\[4pt]
Découpage : 家 (jiā, famille) + 和 (hé, harmonie) + 万事 (wàn shì, toutes les affaires) + 兴 (xīng, prospérer).\\[4pt]
Ce proverbe est affiché dans de nombreuses maisons chinoises, souvent calligraphié sur un rouleau rouge accroché au mur du salon. Il exprime l'idée que l'harmonie familiale est la base de toute réussite.
\end{textebox}

\begin{savaistu}[La famille élargie : un point commun Cameroun--Chine]
Au Cameroun comme en Chine, la famille ne se limite pas au père, à la mère et aux enfants : c'est la \emph{famille élargie} qui compte. Les oncles, les tantes, les grands-parents et les cousins jouent un rôle central dans l'éducation des enfants, l'entraide financière et les grandes décisions. En Chine, les grands-parents participent très souvent à la garde des petits-enfants pendant que les parents travaillent ; au Cameroun, un enfant peut être confié à un oncle en ville pour faire ses études. Cette conception diffère de la famille nucléaire occidentale, centrée sur le couple et ses enfants. C'est pourquoi le vocabulaire chinois de la parenté est très précis : il distingue par exemple l'oncle paternel (\zh{伯伯}, \emph{bóbo}) de l'oncle maternel (\zh{舅舅}, \emph{jiùjiu}), et le grand-père paternel (\zh{爷爷}, \emph{yéye}) du grand-père maternel (\zh{外公}, \emph{wàigōng}).
\end{savaistu}

\courssub{Révision des acquis de Seconde : les membres de la famille}

En Seconde, vous avez appris à nommer les membres de la famille proche. Avant de lire le texte, réactivons ce vocabulaire indispensable.

\begin{center}
\begin{tabularx}{\linewidth}{|X|X|X|X|}
\hline
\rowcolor{popblueL} Caractères & Pinyin & Nature & Traduction \\
\hline
家 & jiā & nom & famille, maison \\
\hline
家庭 & jiātíng & nom & famille, foyer \\
\hline
家人 & jiārén & nom & membres de la famille, les siens \\
\hline
爸爸 & bàba & nom & papa, père \\
\hline
妈妈 & māma & nom & maman, mère \\
\hline
哥哥 & gēge & nom & frère aîné \\
\hline
姐姐 & jiějie & nom & sœur aînée \\
\hline
弟弟 & dìdi & nom & frère cadet \\
\hline
妹妹 & mèimei & nom & sœur cadette \\
\hline
爷爷 & yéye & nom & grand-père (paternel) \\
\hline
奶奶 & nǎinai & nom & grand-mère (paternelle) \\
\hline
爱 & ài & verbe / nom & aimer ; amour \\
\hline
帮助 & bāngzhù & verbe / nom & aider ; aide \\
\hline
关心 & guānxīn & verbe & se soucier de, entourer de soins \\
\hline
团结 & tuánjié & verbe / nom & s'unir ; union, solidarité \\
\hline
互相 & hùxiāng & adverbe & mutuellement, l'un l'autre \\
\hline
\end{tabularx}
\end{center}

\begin{attention}[Le chinois distingue l'aîné du cadet]
Contrairement au français qui dit simplement « frère » et « sœur », le chinois impose de préciser l'ordre de naissance : \zh{哥哥} (frère aîné) ≠ \zh{弟弟} (frère cadet) ; \zh{姐姐} (sœur aînée) ≠ \zh{妹妹} (sœur cadette). Erreur fréquente des francophones : traduire « mon frère » par \zh{我的哥哥} alors qu'il est plus jeune que soi. Demandez-vous toujours : aîné ou cadet ? Autre piège : dans les noms de parenté redoublés (\zh{爸爸}, \zh{妈妈}, \zh{哥哥}…), la deuxième syllabe se prononce au \textbf{ton neutre}, sans accent : \emph{bàba}, \emph{māma}, \emph{gēge}.
\end{attention}

\begin{exoresolu}[Révision : compléter avec le bon terme de parenté]
Énoncé : complétez la phrase avec le mot qui convient. « Ma sœur cadette aime beaucoup ma grand-mère. » → \zh{我的\_\_\_\_很爱我的奶奶。}
\tcblower
Solution détaillée : « sœur cadette » se dit \zh{妹妹} (\emph{mèimei}) ; on ne peut pas utiliser \zh{姐姐} qui désigne la sœur aînée. La phrase correcte est donc : \zh{我的妹妹很爱我的奶奶。} (\emph{Wǒ de mèimei hěn ài wǒ de nǎinai.}) « Ma petite sœur aime beaucoup ma grand-mère. » Remarquez : \zh{很} (\emph{hěn}, « très ») précède le verbe \zh{爱}, et \zh{的} marque la possession (\zh{我的}, « mon/ma »). Structure : Sujet (\zh{我的妹妹}) + adverbe (\zh{很}) + verbe (\zh{爱}) + complément (\zh{我的奶奶}).
\end{exoresolu}

\courssub{Objectifs de la leçon et tâche finale}

\begin{aretenir}[Ce que vous saurez faire à la fin de la leçon]
\begin{enumerate}
\item Lire et comprendre un court texte chinois sur l'amour et la solidarité familiale ;
\item Reconnaître, prononcer (pinyin et tons) et employer le vocabulaire du thème : \zh{爱}, \zh{帮助}, \zh{关心}, \zh{团结}, \zh{家庭}, \zh{互相}… ;
\item Identifier les clés (radicaux) et l'ordre des traits de quelques caractères du thème ;
\item Répondre à des questions de compréhension sur le texte.
\end{enumerate}
\textbf{Tâche finale :} présenter oralement et par écrit votre famille en chinois (4 à 6 phrases), en racontant un geste de solidarité familiale, sur le modèle : \zh{我很爱我的家。我的家人常常互相帮助。} (\emph{Wǒ hěn ài wǒ de jiā. Wǒ de jiārén chángcháng hùxiāng bāngzhù.}) « J'aime beaucoup ma famille. Les membres de ma famille s'entraident souvent. » Notez l'adverbe \zh{互相} (\emph{hùxiāng}, « mutuellement »), placé \textbf{avant} le verbe : \zh{互相帮助}, « s'aider mutuellement, s'entraider ».
\end{aretenir}

\courssub{Première écoute globale du texte (sans support écrit)}

L'enseignant lit maintenant le texte support à voix haute, deux fois, à vitesse normale. Vous écoutez \textbf{sans regarder le texte écrit} : l'objectif n'est pas de tout comprendre, mais de saisir le sens général et de repérer les mots que vous connaissez déjà.

\begin{methode}[Comment écouter un texte chinois]
\begin{enumerate}
\item \textbf{Première écoute : repérez les mots connus.} Notez mentalement les mots déjà appris (\zh{家}, \zh{爱}, \zh{妈妈}, \zh{爸爸}…). Ils vous donnent le thème et les personnages.
\item \textbf{Deuxième écoute : repérez le sujet et le verbe de chaque phrase.} En chinois comme en français, la phrase de base suit l'ordre Sujet + Verbe (+ Complément). Demandez-vous : qui fait quoi ?
\item \textbf{Formulez une hypothèse globale.} En une phrase en français : de qui parle le texte et que se passe-t-il ?
\item \textbf{Vérifiez ensuite avec le texte écrit.} La lecture qui suivra confirmera ou corrigera votre hypothèse.
\end{enumerate}
Ne cherchez jamais à traduire mot à mot pendant l'écoute : c'est le meilleur moyen de perdre le fil. Visez d'abord le sens global.
\end{methode}

\begin{experience}[À vous d'écouter]
Fermez votre manuel. L'enseignant lit le texte deux fois. Après la première écoute, levez la main et citez un mot que vous avez reconnu. Après la deuxième écoute, répondez à ces deux questions en français :
\begin{enumerate}
\item De qui parle le texte ?
\item Quel événement rassemble la famille ?
\end{enumerate}
Comparez vos réponses avec celles de votre voisin avant la lecture écrite de la section suivante.
\end{experience}

\begin{exercice}[Préparation à la lecture]
Sans dictionnaire, associez chaque mot chinois à sa traduction française : \zh{爱}, \zh{帮助}, \zh{关心}, \zh{团结} / a) s'unir ; b) aimer ; c) se soucier de ; d) aider. Puis composez une phrase simple avec \zh{爱} et un membre de votre famille.
\end{exercice}

\begin{corrige}[Exercice de préparation]
\zh{爱} → b) aimer ; \zh{帮助} → d) aider ; \zh{关心} → c) se soucier de ; \zh{团结} → a) s'unir. Exemple de phrase : \zh{我爱我的妈妈。} (\emph{Wǒ ài wǒ de māma.}) « J'aime ma maman. » Structure : Sujet (\zh{我}) + Verbe (\zh{爱}) + Complément (\zh{我的妈妈}). Autres réponses possibles : \zh{我爱我的哥哥。} (\emph{Wǒ ài wǒ de gēge.}) « J'aime mon grand frère » ; \zh{我很爱我的家。} (\emph{Wǒ hěn ài wǒ de jiā.}) « J'aime beaucoup ma famille. »
\end{corrige}

\coursec{Texte support : lecture et traduction}

Dans la section précédente, nous avons découvert le thème du module : la famille, lieu d'amour et de solidarité. Vous avez déjà rencontré quelques mots comme \zh{家} (\emph{jiā}, « famille, maison ») et \zh{爱} (\emph{ài}, « aimer »). Nous allons maintenant entrer dans le cœur de la leçon : la lecture d'un court texte chinois original, phrase par phrase, jusqu'à pouvoir le lire à voix haute, le comprendre sans traduction et le traduire vous-mêmes. C'est la compétence de base du programme de Première : \cle{comprendre un texte écrit} sur la vie familiale.

\courssub{Première lecture globale : écouter avant de comprendre}

Avant toute analyse, lisez le texte en silence une première fois, sans chercher à tout comprendre. Repérez seulement : de qui parle-t-on ? Que font les personnages ? Où et quand ?

\begin{textebox}[我有一个幸福的家 — « J'ai une famille heureuse »]
我有一个幸福的家。\\
Wǒ yǒu yí ge xìngfú de jiā.\\
爸爸妈妈很爱我，我也很爱他们。\\
Bàba māma hěn ài wǒ, wǒ yě hěn ài tāmen.\\
爸爸每天努力工作，妈妈在家做饭。\\
Bàba měitiān nǔlì gōngzuò, māma zài jiā zuòfàn.\\
我和妹妹帮助妈妈做家务。\\
Wǒ hé mèimei bāngzhù māma zuò jiāwù.\\
我们每天互相帮助，互相关心。\\
Wǒmen měitiān hùxiāng bāngzhù, hùxiāng guānxīn.\\
周末，全家人一起吃饭、聊天。\\
Zhōumò, quánjiārén yìqǐ chīfàn, liáotiān.\\
我觉得家是最温暖的地方。\\
Wǒ juéde jiā shì zuì wēnnuǎn de dìfang.\\
我爱我的家！\\
Wǒ ài wǒ de jiā!
\end{textebox}

\textbf{Traduction naturelle d'ensemble :} « J'ai une famille heureuse. Papa et maman m'aiment beaucoup, et moi aussi je les aime beaucoup. Papa travaille dur chaque jour, maman cuisine à la maison. Ma petite sœur et moi aidons maman à faire les tâches ménagères. Chaque jour, nous nous aidons et nous nous soucions les uns des autres. Le week-end, toute la famille mange et bavarde ensemble. Je trouve que la famille est l'endroit le plus chaleureux. J'aime ma famille ! »

\begin{savaistu}[La famille en Chine et au Cameroun]
En Chine comme au Cameroun, la famille est au cœur de la société. Le mot \zh{家} (\emph{jiā}) désigne à la fois la « maison » et la « famille » : pour un Chinois, les deux notions sont inséparables. Le repas du week-end en famille (\zh{全家人一起吃饭}) est un moment sacré, exactement comme le grand repas familial du dimanche à Yaoundé ou à Douala, où l'on partage le ndolé ou le poulet DG. La solidarité familiale — aider ses parents, s'occuper de ses frères et sœurs — est une valeur partagée par les deux cultures.
\end{savaistu}

\courssub{Lecture phrase par phrase : pinyin, traduction littérale, traduction naturelle}

Pour bien comprendre un texte chinois, on procède en trois temps pour chaque phrase : \textbf{prononciation} (pinyin et tons), \textbf{traduction littérale} (mot à mot, pour voir la structure), puis \textbf{traduction naturelle} (en bon français).

\textbf{Phrase 1.} \zh{我有一个幸福的家。}\\
\emph{Wǒ yǒu yí ge xìngfú de jiā.}\\
Littéral : « Moi + avoir + un (classificateur) + heureux + (de) + famille. »\\
Naturel : « J'ai une famille heureuse. »\\
Observation : le verbe \zh{有} (\emph{yǒu}, « avoir, il y a ») fonctionne comme le français « avoir ». Entre le numéral \zh{一} et le nom, un \cle{classificateur} est obligatoire : ici \zh{个} (\emph{ge}), le classificateur général, le plus fréquent. Remarquez aussi la \textbf{mutation de ton} : \zh{一} seul se dit \emph{yī} (1\textsuperscript{er} ton), mais devant un 4\textsuperscript{e} ton (ici \emph{ge}, ton neutre issu de \emph{gè}), il se prononce \emph{yí}. Enfin, \zh{的} (\emph{de}) relie l'adjectif \zh{幸福} au nom \zh{家} : « une famille \emph{heureuse} ».

\textbf{Phrase 2.} \zh{爸爸妈妈很爱我，我也很爱他们。}\\
\emph{Bàba māma hěn ài wǒ, wǒ yě hěn ài tāmen.}\\
Littéral : « Papa maman très aimer moi, moi aussi très aimer eux. »\\
Naturel : « Papa et maman m'aiment beaucoup, et moi aussi je les aime beaucoup. »\\
Observation : en chinois, on juxtapose souvent \zh{爸爸妈妈} sans mot « et ». L'adverbe \zh{很} (\emph{hěn}) se place \textbf{avant} le verbe \zh{爱} ; devant un verbe psychologique comme « aimer », \zh{很} garde pleinement son sens de « beaucoup ». \zh{也} (\emph{yě}, « aussi ») se place juste après le sujet, avant le verbe. Le pluriel des personnes se marque avec \zh{们} : \zh{他们} (\emph{tāmen}, « eux, ils »).

\textbf{Phrase 3.} \zh{爸爸每天努力工作，妈妈在家做饭。}\\
\emph{Bàba měitiān nǔlì gōngzuò, māma zài jiā zuòfàn.}\\
Littéral : « Papa chaque jour avec-effort travailler, maman à maison faire-cuisine. »\\
Naturel : « Papa travaille dur chaque jour, maman cuisine à la maison. »\\
Observation : le complément de temps \zh{每天} (\emph{měitiān}, « chaque jour ») se place juste après le sujet, avant le verbe — contrairement au français où « chaque jour » va souvent à la fin. De même, le complément de lieu \zh{在家} (\emph{zài jiā}, « à la maison ») se place \textbf{avant} le verbe \zh{做饭} : en chinois, on dit « maman à-la-maison cuisine », jamais *\zh{妈妈做饭在家}.

\textbf{Phrase 4.} \zh{我和妹妹帮助妈妈做家务。}\\
\emph{Wǒ hé mèimei bāngzhù māma zuò jiāwù.}\\
Littéral : « Moi et petite-sœur aider maman faire tâches-ménagères. »\\
Naturel : « Ma petite sœur et moi aidons maman à faire les tâches ménagères. »\\
Observation : \zh{和} (\emph{hé}, « et ») relie deux personnes sujets. \zh{帮助} (\emph{bāngzhù}, « aider ») est suivi d'une personne puis d'un verbe : \cle{aider quelqu'un à faire quelque chose} = \zh{帮助 + personne + verbe}.

\textbf{Phrase 5.} \zh{我们每天互相帮助，互相关心。}\\
\emph{Wǒmen měitiān hùxiāng bāngzhù, hùxiāng guānxīn.}\\
Littéral : « Nous chaque jour mutuellement aider, mutuellement soucier-de. »\\
Naturel : « Chaque jour, nous nous aidons et nous nous soucions les uns des autres. »\\
C'est la phrase centrale du texte : elle exprime la \cle{solidarité familiale} grâce au mot-clé \zh{互相} (\emph{hùxiāng}, « mutuellement, les uns les autres »).

\textbf{Phrase 6.} \zh{周末，全家人一起吃饭、聊天。}\\
\emph{Zhōumò, quánjiārén yìqǐ chīfàn, liáotiān.}\\
Littéral : « Week-end, famille-entière ensemble manger, bavarder. »\\
Naturel : « Le week-end, toute la famille mange et bavarde ensemble. »\\
Observation : le complément de temps \zh{周末} peut se placer en tête de phrase (avant le sujet), c'est un choix de style très courant ; après le sujet, il faudrait dire \zh{全家人周末一起吃饭}. \zh{一起} (\emph{yìqǐ}, « ensemble ») se place avant le verbe. La virgule de juxtaposition \zh{、} relie deux verbes coordonnés : \zh{吃饭} et \zh{聊天}.

\textbf{Phrase 7.} \zh{我觉得家是最温暖的地方。}\\
\emph{Wǒ juéde jiā shì zuì wēnnuǎn de dìfang.}\\
Littéral : « Moi trouver famille être le-plus chaleureux (de) endroit. »\\
Naturel : « Je trouve que la famille est l'endroit le plus chaleureux. »\\
Observation : \zh{觉得} (\emph{juéde}, « penser, trouver ») introduit une opinion, suivi directement d'une phrase complète, sans mot équivalent à « que ». Le superlatif chinois est très simple : \zh{最} (\emph{zuì}, « le plus ») + adjectif, sans « de » ni article.

\textbf{Phrase 8.} \zh{我爱我的家！}\\
\emph{Wǒ ài wǒ de jiā!}\\
Naturel : « J'aime ma famille ! »\\
Phrase de conclusion, courte et forte : sujet + verbe + complément, l'ordre le plus simple du chinois. \zh{我的家} : le pronom possessif se forme avec \zh{的} : \zh{我} + \zh{的} = « mon, ma ».

\begin{exemplebox}[Deux phrases modèles à réutiliser]
\zh{我很爱我的家。} \emph{Wǒ hěn ài wǒ de jiā.} « J'aime beaucoup ma famille. »\\[2pt]
\zh{全家人一起吃饭。} \emph{Quánjiārén yìqǐ chīfàn.} « Toute la famille mange ensemble. »
\end{exemplebox}

\courssub{Frise des actions : visualiser le déroulement du texte}

Le texte raconte une journée et une semaine types de cette famille. La frise ci-dessous résume l'enchaînement des actions ; servez-vous-en pour raconter le texte de mémoire.

\begin{popfigure}
\begin{tikzpicture}[
  etape/.style={rectangle, rounded corners=4pt, draw=popblue, fill=popblueL, align=center, minimum width=2.9cm, minimum height=1.5cm, font=\small},
  fleche/.style={-{Stealth[length=3mm]}, thick, poporange}
]
\node[etape, align=center] (a) at (0,0){爸爸努力工作\\ \emph{Bàba nǔlì gōngzuò}\\ le père travaille dur};
\node[etape, right=1.1cm of a, align=center] (b){妈妈做饭\\ \emph{māma zuòfàn}\\ la mère cuisine};
\node[etape, right=1.1cm of b, align=center] (c){我们帮助妈妈\\ \emph{wǒmen bāngzhù māma}\\ les enfants aident};
\node[etape, right=1.1cm of c, align=center] (d){全家一起吃饭\\ \emph{quánjiā yìqǐ chīfàn}\\ toute la famille dîne};
\draw[fleche] (a) -- (b);
\draw[fleche] (b) -- (c);
\draw[fleche] (c) -- (d);
\end{tikzpicture}
\legende{Frise des actions du texte : du travail du père au repas familial du week-end.}
\end{popfigure}

\courssub{Travail sur la prononciation : les tons difficiles du texte}

La lecture à voix haute exige de maîtriser trois difficultés de tons présentes dans notre texte.

\begin{propriete}[Trois pièges de tons à soigner]
\begin{enumerate}
\item \zh{爱} \emph{ài} : \textbf{4\textsuperscript{e} ton}, ton descendant et sec, comme un « oui ! » affirmé en français. Ne le laissez pas monter : dites \emph{ài!} et non *\emph{ái}.
\item \zh{帮} \emph{bāng} : \textbf{1\textsuperscript{er} ton}, ton haut et plat, tenu sans monter ni descendre, comme si vous chantiez une note longue.
\item \zh{团结} \emph{tuánjié} (« solidarité, unité ») : \textbf{2\textsuperscript{e} + 2\textsuperscript{e} ton}, deux tons montants à la suite. Montez bien deux fois : \emph{tuán} (montée) puis \emph{jié} (nouvelle montée), sans retomber entre les deux.
\end{enumerate}
\end{propriete}

\begin{attention}[Erreurs de prononciation typiques des francophones]
\begin{itemize}
\item Dire \emph{ài} avec une intonation de question (montante) : \zh{爱} devient alors méconnaissable. Le 4\textsuperscript{e} ton doit \textbf{descendre} franchement.
\item Avaler le 1\textsuperscript{er} ton de \zh{帮} (\emph{bāng}) : le ton haut doit être \textbf{tenu}, pas expiré.
\item Dans \zh{互相} (\emph{hùxiāng}, 4\textsuperscript{e} + 1\textsuperscript{er} ton), ne pas confondre avec \emph{húxiāng} : la première syllabe descend, la seconde reste haute et plate.
\item Attention à \zh{觉得} : on prononce ici \emph{juéde} (le second caractère est au ton neutre), et non *\emph{juédé}.
\item Attention à la mutation de ton de \zh{一} : dans \zh{一个} on dit \emph{yí ge}, dans \zh{一起} on dit \emph{yìqǐ} (devant un 4\textsuperscript{e} ton, \zh{一} monte ; devant un autre ton, il descend).
\end{itemize}
\end{attention}

\textbf{Relecture chorale puis individuelle.} La classe lit d'abord le texte entier ensemble, lentement, en marquant bien les tons (le professeur donne le rythme en frappant dans ses mains à chaque syllabe). Puis chaque élève lit deux phrases à voix haute ; les autres vérifient les tons de \zh{爱}, \zh{帮}, \zh{互相} et \zh{一起}.

\courssub{Grammaire en contexte : 互相, mot-clé de la solidarité}

Observons la phrase 5 : \zh{我们互相帮助} — « nous nous aidons mutuellement ». Où se place \zh{互相} ?

\begin{propriete}[Place de 互相 (hùxiāng, « mutuellement »)]
Structure : \textbf{Sujet (pluriel) + 互相 + Verbe}\\[3pt]
\zh{互相} est un adverbe : il se place \textbf{toujours avant le verbe}, jamais après.\\[3pt]
\zh{我们互相帮助。} \emph{Wǒmen hùxiāng bāngzhù.} « Nous nous aidons mutuellement. »\\
\zh{他们互相关心。} \emph{Tāmen hùxiāng guānxīn.} « Ils se soucient les uns des autres. »\\
\zh{同学们互相学习。} \emph{Tóngxuémen hùxiāng xuéxí.} « Les camarades apprennent les uns des autres. »
\end{propriete}

\begin{attention}[Ne placez jamais 互相 après le verbe]
Le français dit « nous aidons \emph{mutuellement} » avec l'adverbe après le verbe ; le chinois fait l'inverse. On dit \zh{我们互相帮助} et jamais *\zh{我们帮助互相}. De même : \zh{我们互相爱护} (\emph{Wǒmen hùxiāng àihù}, « nous nous chérissons les uns les autres »), jamais *\zh{我们爱护互相}.
\end{attention}

\begin{methode}[Traduire une phrase chinoise en quatre étapes]
\begin{enumerate}
\item \textbf{Identifier le sujet} : qui fait l'action ? (\zh{我们}, \zh{爸爸}, \zh{全家人}…)
\item \textbf{Identifier le verbe} : quelle action ? (\zh{帮助}, \zh{做饭}, \zh{爱}…)
\item \textbf{Identifier le complément} : qui ou quoi reçoit l'action ? (\zh{妈妈}, \zh{家务}, \zh{他们}…)
\item \textbf{Traiter les adverbes avant le verbe} : \zh{很} (« très »), \zh{也} (« aussi »), \zh{都} (« tous »), \zh{互相} (« mutuellement »), \zh{一起} (« ensemble ») se traduisent en français à la place naturelle, souvent après le verbe ou en fin de phrase.
\end{enumerate}
Exemple : \zh{我也很爱他们} → sujet \zh{我} + adverbe \zh{也} + adverbe \zh{很} + verbe \zh{爱} + complément \zh{他们} → « Moi aussi très aimer eux » → « Moi aussi, je les aime beaucoup. »
\end{methode}

\begin{center}
\begin{tabularx}{\linewidth}{|X|X|X|}
\hline
\rowcolor{popblueL} \textbf{Structure chinoise} & \textbf{Exemple (caractères + pinyin)} & \textbf{Traduction} \\
\hline
Sujet + 很 + verbe & \zh{爸爸妈妈很爱我。} \emph{Bàba māma hěn ài wǒ.} & Papa et maman m'aiment beaucoup. \\
\hline
Sujet + 也 + verbe & \zh{我也很爱他们。} \emph{Wǒ yě hěn ài tāmen.} & Moi aussi je les aime beaucoup. \\
\hline
Sujet + 互相 + verbe & \zh{我们互相帮助。} \emph{Wǒmen hùxiāng bāngzhù.} & Nous nous aidons mutuellement. \\
\hline
Sujet + 一起 + verbe & \zh{全家人一起吃饭。} \emph{Quánjiārén yìqǐ chīfàn.} & Toute la famille mange ensemble. \\
\hline
最 + adjectif + 的 + nom & \zh{最温暖的地方} \emph{zuì wēnnuǎn de dìfang} & l'endroit le plus chaleureux \\
\hline
\end{tabularx}
\end{center}

\courssub{Repérage des mots nouveaux pour le glossaire}

Avant de constituer le glossaire officiel (section suivante), soulignons dans le texte les mots nouveaux et vérifions que nous savons les prononcer et les traduire.

\begin{center}
\begin{tabularx}{\linewidth}{|X|X|X|X|}
\hline
\rowcolor{popblueL} \textbf{Caractères} & \textbf{Pinyin} & \textbf{Nature} & \textbf{Traduction} \\
\hline
\zh{幸福} & xìngfú & adjectif & heureux, épanoui \\
\hline
\zh{爱} & ài & verbe & aimer \\
\hline
\zh{努力} & nǔlì & adverbe / verbe & avec effort, s'efforcer \\
\hline
\zh{做饭} & zuòfàn & verbe & cuisiner, faire à manger \\
\hline
\zh{帮助} & bāngzhù & verbe & aider \\
\hline
\zh{家务} & jiāwù & nom & tâches ménagères \\
\hline
\zh{互相} & hùxiāng & adverbe & mutuellement, les uns les autres \\
\hline
\zh{关心} & guānxīn & verbe & se soucier de, prendre soin de \\
\hline
\zh{全家} & quánjiā & nom & toute la famille \\
\hline
\zh{一起} & yìqǐ & adverbe & ensemble \\
\hline
\zh{聊天} & liáotiān & verbe & bavarder \\
\hline
\zh{温暖} & wēnnuǎn & adjectif & chaleureux \\
\hline
\end{tabularx}
\end{center}

\begin{exoresolu}[Traduction guidée d'une phrase du texte]
Énoncé : traduisez en français naturel la phrase \zh{我和妹妹帮助妈妈做家务。} en appliquant la méthode en quatre étapes, puis repérez les adverbes éventuels.
\tcblower
Solution détaillée :
\begin{enumerate}
\item Sujet : \zh{我和妹妹} (\emph{wǒ hé mèimei}) → « ma petite sœur et moi » (en français, on se mentionne en dernier par politesse).
\item Verbe : \zh{帮助} (\emph{bāngzhù}) → « aider ».
\item Compléments : \zh{妈妈} (« maman », la personne aidée) + \zh{做家务} (\emph{zuò jiāwù}, « faire les tâches ménagères », l'action aidée).
\item Adverbes : il n'y en a pas ici ; l'ordre chinois sujet–verbe–complément correspond à l'ordre français.
\end{enumerate}
Traduction littérale : « Moi et petite-sœur aider maman faire tâches-ménagères. »\\
Traduction naturelle : « Ma petite sœur et moi aidons maman à faire les tâches ménagères. »\\
Piège évité : ne pas traduire \zh{和} par « avec » ici ; entre deux personnes sujets, \zh{和} signifie « et ».
\end{exoresolu}

\begin{exercice}[À vous de lire et de traduire]
\begin{enumerate}
\item Lisez à voix haute les phrases 2, 5 et 6 du texte en veillant aux tons de \zh{爱}, \zh{互相} et \zh{一起}.
\item Traduisez en français : \zh{我们每天互相关心。} puis \zh{我觉得家是最温暖的地方。}
\item Recopiez la phrase qui exprime le mieux la solidarité familiale dans le texte, avec son pinyin.
\end{enumerate}
\end{exercice}

\begin{corrige}[Exercice « À vous de lire et de traduire »]
\begin{enumerate}
\item Lecture contrôlée en classe : \emph{Bàba māma hěn ài wǒ, wǒ yě hěn ài tāmen.} — \emph{Wǒmen měitiān hùxiāng bāngzhù, hùxiāng guānxīn.} — \emph{Zhōumò, quánjiārén yìqǐ chīfàn, liáotiān.} Vérifier le 4\textsuperscript{e} ton de \emph{ài}, le 4\textsuperscript{e} + 1\textsuperscript{er} de \emph{hùxiāng}, la mutation de ton dans \emph{yìqǐ}.
\item \zh{我们每天互相关心。} : « Chaque jour, nous nous soucions les uns des autres. » (\zh{互相} avant le verbe \zh{关心}). \zh{我觉得家是最温暖的地方。} : « Je trouve que la famille est l'endroit le plus chaleureux. » (superlatif \zh{最} + adjectif \zh{温暖}).
\item La phrase attendue est \zh{我们每天互相帮助，互相关心。} (\emph{Wǒmen měitiān hùxiāng bāngzhù, hùxiāng guānxīn.}) : c'est elle qui résume l'entraide au sein de la famille.
\end{enumerate}
\end{corrige}

\begin{aretenir}[L'essentiel de cette section]
\begin{itemize}
\item Un texte chinois se lit en trois temps : prononciation (pinyin, tons), traduction littérale (structure), traduction naturelle (bon français).
\item Les compléments de temps (\zh{每天}, \zh{周末}), de lieu (\zh{在家}) et les adverbes (\zh{很}, \zh{也}, \zh{都}, \zh{互相}, \zh{一起}) se placent \textbf{avant le verbe} en chinois.
\item \zh{互相} (\emph{hùxiāng}) = « mutuellement » : jamais après le verbe.
\item Le superlatif se construit avec \zh{最} + adjectif : \zh{最温暖的地方}, « l'endroit le plus chaleureux ».
\item Soignez les tons : \emph{ài} (4\textsuperscript{e}, descendant), \emph{bāng} (1\textsuperscript{er}, haut et plat), \emph{tuánjié} (2\textsuperscript{e} + 2\textsuperscript{e}, deux montées), et la mutation de \zh{一} (\emph{yí ge}, \emph{yìqǐ}).
\end{itemize}
\end{aretenir}

Le texte est maintenant compris et su par cœur presque sans effort. La section suivante reprendra ces mots nouveaux pour construire le glossaire complet du thème « amour de la famille et solidarité familiale ».

\coursec{Glossaire du thème}

Dans la section précédente, nous avons lu et traduit le texte support sur l'amour de la famille et la solidarité familiale. Vous avez remarqué que certains mots revenaient souvent : \zh{家}, \zh{爱}, \zh{互相}, \zh{帮助}… C'est le signe qu'il s'agit du \cle{vocabulaire central} du thème. Dans cette section, nous allons étudier ces mots un à un : leur prononciation exacte, leur nature grammaticale, leur emploi dans une phrase, et les pièges qu'ils réservent aux francophones. Un bon glossaire n'est pas une liste à réciter : c'est une boîte à outils que vous devez savoir ouvrir au bon moment.

\courssub{Observation : les mots du thème en action}

Avant tout tableau, observons. Relisez ces phrases extraites et adaptées du texte support, et repérez les mots en gras :

\begin{exemplebox}[Observer avant d'apprendre]
\zh{我爱我的家人。} \emph{(Wǒ ài wǒ de jiārén.)} « J'aime les membres de ma famille. »\\
\zh{我们互相关心，互相帮助。} \emph{(Wǒmen hùxiāng guānxīn, hùxiāng bāngzhù.)} « Nous nous soucions les uns des autres et nous nous aidons mutuellement. »\\
\zh{一家人团结在一起，就是幸福。} \emph{(Yì jiā rén tuánjié zài yìqǐ, jiù shì xìngfú.)} « Une famille unie, c'est le bonheur. »\\
\zh{妈妈照顾生病的奶奶。} \emph{(Māma zhàogù shēngbìng de nǎinai.)} « Maman prend soin de grand-mère malade. »
\end{exemplebox}

Qu'observez-vous ? D'abord, les mots chinois sont \cle{invariables} : \zh{家人} ne prend jamais de marque de pluriel, \zh{爱} ne se conjugue pas. Ensuite, l'adverbe \zh{互相} se place \textbf{avant le verbe}, exactement comme « mutuellement » en français, mais sans pronom réfléchi : le chinois dit littéralement « nous mutuellement aidons », là où le français dit « nous \emph{nous} aidons ». Retenez ces deux faits, nous y reviendrons.

\courssub{Tableau de vocabulaire : les mots clés}

Voici le glossaire officiel de la leçon. Apprenez-le en trois passages : lecture à voix haute, vérification du sens, réemploi dans une phrase personnelle.

\begin{center}
\begin{tabularx}{\linewidth}{|X|X|X|X|}
\hline
\rowcolor{popblueL} Caractères & Pinyin & Nature & Traduction \\
\hline
爱 & ài & verbe & aimer \\
\hline
家 & jiā & nom & famille, maison, foyer \\
\hline
家人 & jiārén & nom & membres de la famille \\
\hline
幸福 & xìngfú & adjectif & heureux, comblé \\
\hline
帮助 & bāngzhù & verbe & aider \\
\hline
互相 & hùxiāng & adverbe & mutuellement, l'un l'autre \\
\hline
关心 & guānxīn & verbe & se soucier de, s'intéresser à \\
\hline
团结 & tuánjié & verbe / adjectif & s'unir ; uni, solidaire \\
\hline
温暖 & wēnnuǎn & adjectif & chaleureux, doux \\
\hline
一起 & yìqǐ & adverbe & ensemble \\
\hline
\end{tabularx}
\end{center}

\begin{center}
\begin{tabularx}{\linewidth}{|X|X|X|X|}
\hline
\rowcolor{popgreenL} Caractères & Pinyin & Nature & Traduction \\
\hline
照顾 & zhàogù & verbe & prendre soin de \\
\hline
支持 & zhīchí & verbe & soutenir, appuyer \\
\hline
分享 & fēnxiǎng & verbe & partager \\
\hline
感情 & gǎnqíng & nom & affection, sentiments \\
\hline
\end{tabularx}
\end{center}

\begin{definition}[家人 (jiārén) : un mot collectif]
\zh{家人} \emph{(jiārén)} signifie « les membres de la famille ». C'est un \cle{nom collectif invariable} : il ne prend \textbf{jamais} de marque de pluriel (le chinois n'en a pas) et il ne se combine pas directement avec un chiffre seul. Pour compter, on utilise le classificateur \zh{口} (kǒu) avec \zh{家} : \zh{我家有四口人。} \emph{(Wǒ jiā yǒu sì kǒu rén.)} « Ma famille compte quatre personnes. » On ne dit jamais \zh{四个家人}. Comparez avec le français « famille », qui reste singulier même s'il désigne plusieurs personnes.
\end{definition}

\courssub{Emplois et exemples commentés}

Étudions maintenant l'emploi de chaque mot clé dans des phrases complètes, avec la structure à retenir.

\begin{exemplebox}[Les verbes du thème en contexte]
\zh{我爱我的家。} \emph{(Wǒ ài wǒ de jiā.)} « J'aime ma famille. » — Structure : Sujet + \zh{爱} + nom.\\
\zh{家人应该互相支持。} \emph{(Jiārén yīnggāi hùxiāng zhīchí.)} « Les membres de la famille doivent se soutenir mutuellement. » — Structure : Sujet + \zh{应该} + \zh{互相} + verbe.\\
\zh{我们互相关心。} \emph{(Wǒmen hùxiāng guānxīn.)} « Nous nous soucions les uns des autres. »\\
\zh{姐姐每天晚上照顾弟弟。} \emph{(Jiějie měitiān wǎnshang zhàogù dìdi.)} « Grande sœur prend soin de son petit frère chaque soir. »\\
\zh{我们一起吃饭，一起分享快乐。} \emph{(Wǒmen yìqǐ chīfàn, yìqǐ fēnxiǎng kuàilè.)} « Nous mangeons ensemble et partageons ensemble la joie. »\\
\zh{他们的感情很深。} \emph{(Tāmen de gǎnqíng hěn shēn.)} « Leur affection est très profonde. »\\
\zh{温暖的家让我很幸福。} \emph{(Wēnnuǎn de jiā ràng wǒ hěn xìngfú.)} « Un foyer chaleureux me rend très heureux. »
\end{exemplebox}

\begin{propriete}[Place de 互相 et de 一起]
\zh{互相} \emph{(hùxiāng)} et \zh{一起} \emph{(yìqǐ)} sont des \cle{adverbes} : en chinois, l'adverbe se place \textbf{toujours avant le verbe}, jamais en fin de phrase comme souvent en français.\\
Structure : Sujet + \zh{互相}/\zh{一起} + Verbe (+ Complément).\\
\zh{他们一起学习。} \emph{(Tāmen yìqǐ xuéxí.)} « Ils étudient ensemble. » (et non \zh{他们学习一起}).\\
Remarque : le français exprime la réciprocité avec un pronom (« nous \emph{nous} aidons ») ; le chinois utilise uniquement l'adverbe \zh{互相}, sans pronom réfléchi. Ne traduisez donc jamais « nous nous » mot à mot.
\end{propriete}

\begin{attention}[爱 (ài) ou 喜欢 (xǐhuan) ?]
Piège classique du francophone ! \zh{爱} s'emploie pour les \textbf{personnes} et les \textbf{choses aimées} profondément : \zh{我爱我的家人。} « J'aime ma famille. » ; \zh{我爱中国。} « J'aime la Chine. »\\
Mais pour dire « aimer \textbf{faire} quelque chose », on utilise \zh{喜欢} + verbe : \zh{我喜欢帮助家人。} \emph{(Wǒ xǐhuan bāngzhù jiārén.)} « J'aime aider ma famille. » Dire \zh{我爱帮助家人} est peu naturel ou trop emphatique. Retenez : \zh{爱} + nom ; \zh{喜欢} + verbe (ou + nom pour les goûts : \zh{我喜欢恩多莱。} « J'aime le ndolé. »).
\end{attention}

\begin{savaistu}[Pourquoi un cochon sous un toit ?]
Le caractère \zh{家} (jiā, « famille, maison ») est composé de deux éléments : la clé du \textbf{toit} \zh{宀} (mián) en haut, et \zh{豕} (shǐ), un ancien caractère signifiant « cochon », en dessous. Dans la Chine ancienne, posséder un cochon sous son toit était le signe d'un foyer établi et prospère : la maisonnée avait de quoi vivre. Le caractère raconte donc une histoire : un toit + de quoi manger = une famille. Au Cameroun comme en Chine, la maison familiale reste le cœur de la solidarité : on y accueille les proches, on y partage les repas, on y soigne les anciens.
\end{savaistu}

\courssub{Carte mentale du lexique}

Pour mémoriser durablement, organisez les mots autour du mot central \zh{家} :

\begin{popfigure}
\begin{tikzpicture}[mindmap, grow cyclic, every node/.style={concept}, concept color=popblueL, text=popdark, root concept/.append style={concept color=popblue, text=white, font=\large\bfseries}, level 1/.append style={level distance=3.2cm, sibling angle=72}, level 2/.append style={level distance=2.2cm, sibling angle=45, concept color=popcream}]
\node[root concept, align=center]{家 jiā\\ famille}
  child { node[align=center]{爱 ài\\ aimer} }
  child { node[align=center]{家人 jiārén\\ membres} child { node[align=center]{感情 gǎnqíng\\ affection} } }
  child { node[align=center]{幸福 xìngfú\\ heureux} child { node[align=center]{温暖 wēnnuǎn\\ chaleureux} } }
  child { node[align=center]{团结 tuánjié\\ s'unir} child { node[align=center]{一起 yìqǐ\\ ensemble} } }
  child { node[align=center]{帮助 bāngzhù\\ aider} child { node[align=center]{互相关心\\ hùxiāng guānxīn} } child { node[align=center]{照顾 zhàogù\\ soigner} } };
\end{tikzpicture}
\legende{Carte mentale du vocabulaire de la solidarité familiale autour de 家.}
\end{popfigure}

\courssub{Texte de réemploi : lire le vocabulaire en contexte}

\begin{textebox}[一个温暖的家庭 (Yí ge wēnnuǎn de jiātíng) — « Une famille chaleureuse »]
\zh{我家有五口人：爸爸、妈妈、哥哥、妹妹和我。我们住在雅温得。}\\
\emph{(Wǒ jiā yǒu wǔ kǒu rén : bàba, māma, gēge, mèimei hé wǒ. Wǒmen zhù zài Yǎwēndé.)}\\
« Ma famille compte cinq personnes : papa, maman, mon grand frère, ma petite sœur et moi. Nous habitons à Yaoundé. »\\
\zh{我们一家人非常团结。爸爸妈妈每天工作很辛苦，但是他们总是关心我们。}\\
\emph{(Wǒmen yì jiā rén fēicháng tuánjié. Bàba māma měitiān gōngzuò hěn xīnkǔ, dànshì tāmen zǒngshì guānxīn wǒmen.)}\\
« Notre famille est très unie. Papa et maman travaillent dur chaque jour, mais ils se soucient toujours de nous. »\\
\zh{哥哥喜欢帮助我做作业，我喜欢照顾妹妹。我们互相关心，互相支持。}\\
\emph{(Gēge xǐhuan bāngzhù wǒ zuò zuòyè, wǒ xǐhuan zhàogù mèimei. Wǒmen hùxiāng guānxīn, hùxiāng zhīchí.)}\\
« Mon grand frère aime m'aider à faire mes devoirs, et moi j'aime prendre soin de ma petite sœur. Nous nous soucions les uns des autres et nous nous soutenons mutuellement. »\\
\zh{晚上，我们一起吃饭，一起分享一天的事情。星期天，我们一起去看奶奶。奶奶说：「家人在一起，就是最大的幸福。」}\\
\emph{(Wǎnshang, wǒmen yìqǐ chīfàn, yìqǐ fēnxiǎng yì tiān de shìqing. Xīngqītiān, wǒmen yìqǐ qù kàn nǎinai. Nǎinai shuō : « Jiārén zài yìqǐ, jiù shì zuì dà de xìngfú. »)}\\
« Le soir, nous dînons ensemble et nous partageons les événements de la journée. Le dimanche, nous allons ensemble voir grand-mère. Grand-mère dit : "La famille réunie, c'est le plus grand bonheur." »\\
\zh{我爱我的家。我们的家不大，但是很温暖，因为我们的感情很深。}\\
\emph{(Wǒ ài wǒ de jiā. Wǒmen de jiā bù dà, dànshì hěn wēnnuǎn, yīnwèi wǒmen de gǎnqíng hěn shēn.)}\\
« J'aime ma famille. Notre maison n'est pas grande, mais elle est très chaleureuse, parce que notre affection est profonde. »
\end{textebox}

Ce court texte réemploie \textbf{treize} mots du glossaire. Soulignez-les tous : c'est le meilleur test de reconnaissance.

\courssub{Exercice de prononciation : la bataille des tons}

Le chinois est une langue à tons : changer le ton change le mot. Deux paires de notre glossaire sont dangereuses pour les francophones.

\begin{propriete}[Paires minimales à distinguer]
\zh{爱} \emph{ài} (4\textsuperscript{e} ton, descendant et sec) « aimer » ≠ \emph{āi} (1\textsuperscript{er} ton, haut et plat) comme dans \zh{哎} « hé ! ».\\
\zh{家} \emph{jiā} (1\textsuperscript{er} ton, haut et plat) « famille » ≠ \zh{假} \emph{jiǎ} (3\textsuperscript{e} ton, qui descend puis remonte) « faux, vacances ».\\
Méthode : le 1\textsuperscript{er} ton se chante haut et long, comme une note tenue ; le 4\textsuperscript{e} ton tombe brutalement, comme un ordre ; le 3\textsuperscript{e} ton plonge puis rebondit. Prononcez à voix haute : \emph{jiā — jiǎ — jiā} ; \emph{ài — āi — ài}. Répétez cinq fois chaque série.
\end{propriete}

\begin{exoresolu}[Distinguer les tons et employer le bon mot]
Énoncé. (1) Transcrivez en pinyin avec tons : 爱, 家, 幸福, 互相, 照顾. (2) Complétez avec \zh{爱} ou \zh{喜欢} : \zh{我\_\_\_\_我的家人。} / \zh{我\_\_\_\_帮助妈妈。} (3) Traduisez : \zh{我们互相支持。}
\tcblower
Solution détaillée. (1) 爱 \emph{ài} (4\textsuperscript{e} ton), 家 \emph{jiā} (1\textsuperscript{er} ton), 幸福 \emph{xìngfú} (4\textsuperscript{e} + 2\textsuperscript{e} ton), 互相 \emph{hùxiāng} (4\textsuperscript{e} + 1\textsuperscript{er} ton), 照顾 \emph{zhàogù} (4\textsuperscript{e} + 4\textsuperscript{e} ton). Attention à \emph{xìng} : la voyelle est \emph{i}, le \emph{x} se prononce comme un « s » mouillé, la langue contre les dents du bas.\\
(2) \zh{我爱我的家人。} : devant un \textbf{nom} (\zh{家人}), on emploie \zh{爱}. \zh{我喜欢帮助妈妈。} : devant un \textbf{verbe} (\zh{帮助}), on emploie \zh{喜欢}. C'est exactement la règle vue dans la boîte « Attention ».\\
(3) \zh{我们互相支持。} \emph{(Wǒmen hùxiāng zhīchí.)} « Nous nous soutenons mutuellement. » Ne traduisez pas « nous nous soutenons » par deux mots : le chinois n'a pas de pronom réfléchi ; \zh{互相} suffit, placé avant le verbe.
\end{exoresolu}

\courssub{Mémorisation guidée : la méthode en trois temps}

\begin{methode}[Apprendre 14 mots en 15 minutes]
\begin{enumerate}
\item \textbf{Lecture} (5 min) : lisez le tableau à voix haute, colonne par colonne, en exagérant les tons. Regardez le caractère en prononçant : l'œil, la bouche et l'oreille travaillent ensemble.
\item \textbf{Recouvrement} (5 min) : cachez la colonne « Traduction » avec une feuille et donnez le sens de chaque mot ; puis cachez les caractères et retrouvez-les à partir de la traduction. Cochez les mots ratés et repassez uniquement ceux-là.
\item \textbf{Dictée de 5 mots} (5 min) : faites-vous dicter ces cinq mots : \zh{家人}, \zh{幸福}, \zh{互相}, \zh{帮助}, \zh{温暖}. Écrivez le caractère, le pinyin avec tons et la traduction. Vérifiez : un ton oublié compte comme une faute.
\end{enumerate}
\end{methode}

\begin{exercice}[À vous de jouer]
(1) Traduisez en chinois : « Les membres de la famille doivent s'aider mutuellement. » / « Nous mangeons ensemble le soir. » (2) Construisez une phrase avec \zh{照顾} et une avec \zh{感情}. (3) Classez les 14 mots du glossaire en trois colonnes : verbes, noms, adjectifs et adverbes.
\end{exercice}

\begin{corrige}[Exercice « À vous de jouer »]
(1) \zh{家人应该互相帮助。} \emph{(Jiārén yīnggāi hùxiāng bāngzhù.)} ; \zh{我们晚上一起吃饭。} \emph{(Wǒmen wǎnshang yìqǐ chīfàn.)} — le complément de temps \zh{晚上} se place avant l'adverbe \zh{一起}.\\
(2) Exemples attendus : \zh{我照顾生病的奶奶。} « Je prends soin de grand-mère malade. » ; \zh{我们家的感情很深。} « Dans notre famille, l'affection est profonde. »\\
(3) Verbes : \zh{爱}, \zh{帮助}, \zh{关心}, \zh{团结}, \zh{照顾}, \zh{支持}, \zh{分享}. Noms : \zh{家}, \zh{家人}, \zh{感情}. Adjectifs : \zh{幸福}, \zh{温暖} (\zh{团结} est aussi adjectif). Adverbes : \zh{互相}, \zh{一起}.
\end{corrige}

\begin{aretenir}[L'essentiel du glossaire]
\begin{itemize}
\item \zh{家人} est un nom collectif invariable ; on compte avec \zh{口} : \zh{五口人}.
\item \zh{互相} et \zh{一起} se placent toujours \textbf{avant le verbe} ; pas de pronom réfléchi en chinois.
\item \zh{爱} + nom (personnes, pays) ; \zh{喜欢} + verbe (« aimer faire »).
\item Les tons font le sens : \emph{jiā} « famille » ≠ \emph{jiǎ} « faux » ; \emph{ài} « aimer » (4\textsuperscript{e} ton sec).
\item Le caractère \zh{家} = toit \zh{宀} + cochon \zh{豕} : un foyer où l'on a de quoi vivre.
\end{itemize}
\end{aretenir}

Votre glossaire est maintenant construit, prononcé et mémorisé. Dans la prochaine section, nous ouvrirons ces mots comme des boîtes : nous étudierons leurs \cle{clés} (radicaux) et l'\cle{ordre des traits} pour apprendre à les écrire correctement.

\coursec{Clés (radicaux) et ordre des traits}

Dans la section précédente, nous avons mémorisé le vocabulaire du thème « amour de la famille et solidarité familiale » grâce au tableau du glossaire : \zh{爱} (ài, aimer), \zh{家} (jiā, famille), \zh{帮助} (bāngzhù, aider), \zh{团结} (tuánjié, solidarité), \zh{温暖} (wēnnuǎn, chaleureux). Mais reconnaître un caractère ne suffit pas : à l'examen, vous devrez aussi l'\textbf{écrire correctement}. Pour cela, il faut comprendre comment il est construit. C'est l'objet de cette section : démonter chaque caractère comme on démonte un mécanisme, identifier sa \cle{clé} (radical) et appliquer l'\cle{ordre des traits}.

\courssub{Observation : un caractère n'est pas un dessin, c'est une construction}

Regardons d'abord le mot \zh{爱} (ài). Un élève francophone qui le voit pour la première fois essaie souvent de le « dessiner » d'un seul geste, comme une image. Erreur ! En réalité, ce caractère est \textbf{assemblé} à partir de pièces plus petites :

\begin{center}
\zh{爱} = \zh{爫} (la griffe, la main) + \zh{冖} (le couvercle) + \zh{友} (l'ami)
\end{center}

L'idée ancienne est belle et facile à retenir : \emph{la main qui saisit l'ami} — c'est l'amour, l'affection. Chaque caractère chinois fonctionne ainsi : une ou plusieurs \textbf{composantes} dont l'une, la \cle{clé} (ou radical), donne une indication de \textbf{sens}, et dont une autre donne souvent une indication de \textbf{prononciation} (composant phonétique).

\begin{definition}[Clé (radical) et trait]
La \cle{clé} (部首 \zh{bùshǒu}) est la composante d'un caractère qui classe le caractère dans les dictionnaires et qui indique le plus souvent le \textbf{domaine de sens} : la clé \zh{日} (soleil) apparaît dans les mots liés au temps et à la lumière, la clé \zh{宀} (toit) dans les mots liés à la maison. Le \cle{trait} (笔画 \zh{bǐhuà}) est chaque geste élémentaire du pinceau ou du stylo : le point (丶), l'horizontal (一), le vertical (丨), le crochet, la tombante gauche (丿), la tombante droite (乀).
\end{definition}

\courssub{Les règles générales de l'ordre des traits}

Avant d'étudier nos cinq caractères, rappelons les trois règles d'or qui régissent l'écriture de \textbf{tous} les caractères chinois. Elles sont simples, mais les francophones les oublient souvent parce qu'en français on écrit « au fil de la plume ».

\begin{propriete}[Règles générales de l'ordre des traits]
\begin{enumerate}
\item \textbf{De haut en bas} : on écrit d'abord la partie supérieure du caractère, puis la partie inférieure. Exemple : \zh{爱} : d'abord \zh{爫} en haut, ensuite \zh{友} en bas.
\item \textbf{De gauche à droite} : on écrit d'abord la composante de gauche, puis celle de droite. Exemple : \zh{暖} : d'abord \zh{日} à gauche, ensuite la partie droite.
\item \textbf{L'extérieur avant l'intérieur, puis on ferme} : pour les caractères encadrés, on trace l'enceinte, on écrit l'intérieur, et on trace le trait de fermeture en dernier. Exemple : \zh{团} : l'enceinte \zh{囗} s'ouvre, on écrit \zh{才} à l'intérieur, puis le trait horizontal du bas ferme la boîte.
\end{enumerate}
Règle complémentaire : l'horizontal précède le vertical quand ils se croisent (exemple : \zh{十} : d'abord 一, puis 丨).
\end{propriete}

\courssub{Étude détaillée des cinq caractères du thème}

\textbf{1. \zh{爱} ài — aimer}\par
Radical : \zh{爫} (la griffe, variante de la main). Nombre de traits : \textbf{10}. Structure : haut–bas. On écrit d'abord la griffe \zh{爫} (quatre petits traits en haut), puis le couvercle \zh{冖}, puis \zh{友} (l'ami) en bas. Image mentale : \emph{la main (爫) protège (冖) l'ami (友)}.

\begin{popfigure}
\begin{tikzpicture}[scale=1, every node/.style={font=\normalsize}]
\node[draw=popblue, thick, rounded corners, fill=popblueL, minimum width=1.6cm, minimum height=1.2cm, align=center] (zhua) at (0,2.2) {\zh{爫}\\ griffe (main)};
\node[draw=poporange, thick, rounded corners, fill=poporangeL, minimum width=1.6cm, minimum height=1.2cm, align=center] (mi) at (3,2.2) {\zh{冖}\\ couvercle};
\node[draw=popgreen, thick, rounded corners, fill=popgreenL, minimum width=1.6cm, minimum height=1.2cm, align=center] (you) at (6,2.2) {\zh{友}\\ ami};
\node[font=\Large] at (1.5,2.2) {+};
\node[font=\Large] at (4.5,2.2) {+};
\node[draw=poppurple, very thick, rounded corners, fill=poppurpleL, minimum width=1.8cm, minimum height=1.4cm, align=center] (ai) at (9,2.2) {\zh{爱}\\ ài, aimer};
\draw[-{Stealth}, very thick, popdark] (7,2.2) -- (8,2.2);
\node[align=center, font=\small, text=popmuted] at (4.5,0.8) {Traits 1 à 4 : la griffe \zh{爫} ; traits 5 à 6 : le couvercle \zh{冖} ; traits 7 à 10 : \zh{友}};
\draw[-{Stealth}, thick, popblue] (zhua.south) -- (0,1.2);
\draw[-{Stealth}, thick, poporange] (mi.south) -- (3,1.2);
\draw[-{Stealth}, thick, popgreen] (you.south) -- (6,1.2);
\end{tikzpicture}
\legende{Décomposition de \zh{爱} : la main qui saisit l'ami (10 traits, de haut en bas).}
\end{popfigure}

\textbf{2. \zh{家} jiā — famille, maison}\par
Radical : \zh{宀} (le toit). Nombre de traits : \textbf{10}. Ordre précis : d'abord le \textbf{point} en haut, puis le \textbf{point} à gauche, puis le \textbf{crochet horizontal} — cela forme le toit \zh{宀} — ensuite on écrit \zh{豕} (le cochon) en dessous, en 7 traits. Image mentale : \emph{sous le toit, un cochon} : dans la Chine ancienne, la maison familiale était celle où l'on élevait le cochon.

\begin{popfigure}
\begin{tikzpicture}[scale=1, every node/.style={font=\normalsize}]
\node[draw=popblue, thick, rounded corners, fill=popblueL, minimum width=2cm, minimum height=1.2cm, align=center] (bao) at (0,2) {\zh{宀}\\ toit (traits 1 à 3)};
\node[font=\Large] at (2,2) {+};
\node[draw=poporange, thick, rounded corners, fill=poporangeL, minimum width=2cm, minimum height=1.2cm, align=center] (shi) at (4,2) {\zh{豕}\\ cochon (traits 4 à 10)};
\node[font=\Large] at (6,2) {=};
\node[draw=poppurple, very thick, rounded corners, fill=poppurpleL, minimum width=2cm, minimum height=1.4cm, align=center] (jia) at (8.4,2) {\zh{家}\\ jiā, famille};
\draw[-{Stealth}, thick, popdark] (bao.south) .. controls (2,0.6) .. (shi.south);
\node[align=center, font=\small, text=popmuted] at (4.2,0.2) {Ordre du toit \zh{宀} : 1. point du haut ; 2. point de gauche ; 3. crochet horizontal};
\end{tikzpicture}
\legende{Décomposition de \zh{家} : le toit \zh{宀} + \zh{豕} (10 traits).}
\end{popfigure}

\begin{propriete}[La famille du radical \zh{宀} (le toit)]
Le radical \zh{宀} (toit) apparaît dans de nombreux caractères liés à la maison et à la vie domestique :
\begin{itemize}
\item \zh{家} (jiā) : famille, maison ;
\item \zh{安} (ān) : paix, tranquillité — \emph{une femme \zh{女} sous le toit, c'est la paix} ;
\item \zh{室} (shì) : pièce, salle — comme dans \zh{教室} (jiàoshì, salle de classe).
\end{itemize}
Quand vous rencontrez un caractère inconnu coiffé de \zh{宀}, pensez d'abord : « cela concerne probablement la maison ».
\end{propriete}

\textbf{3. \zh{帮} bāng — aider}\par
Composition : \zh{邦} (pays, région) + \zh{巾} (tissu, turban). Radical : \zh{巾}. Nombre de traits : \textbf{9}. Structure : haut–bas ; on écrit d'abord \zh{邦} en haut (6 traits), puis \zh{巾} en bas (3 traits : vertical, crochet horizontal, vertical). Ce caractère est le cœur du mot \zh{帮助} (bāngzhù, aider) — le verbe de la solidarité.

\textbf{4. \zh{团} tuán — groupe, réunion}\par
Radical : \zh{囗} (l'enceinte, l'enceinte fermée — attention, ce n'est pas \zh{口} la bouche !). Nombre de traits : \textbf{6}. C'est le modèle parfait de la règle « l'extérieur avant l'intérieur, puis on ferme » : 1. vertical gauche de l'enceinte ; 2. crochet horizontal-vertical (haut et droite de l'enceinte) ; 3 à 5. on écrit \zh{才} à l'intérieur ; 6. trait horizontal du bas qui \textbf{ferme} la boîte. On le retrouve dans \zh{团结} (tuánjié, solidarité, union).

\textbf{5. \zh{暖} nuǎn — chaud, chaleureux}\par
Radical : \zh{日} (le soleil), placé \textbf{à gauche}. Nombre de traits : \textbf{13}. Règle « de gauche à droite » : on écrit d'abord le soleil \zh{日} (4 traits), puis la partie droite (9 traits). Logique du sens : \emph{c'est le soleil qui donne la chaleur}. On le trouve dans \zh{温暖} (wēnnuǎn, chaleureux) — la chaleur de la famille.

\begin{center}
\begin{tabularx}{\linewidth}{|X|X|X|X|}
\hline
\rowcolor{popblueL} Caractère & Radical (clé) & Nombre de traits & Structure et règle \\
\hline
\zh{爱} ài, aimer & \zh{爫} griffe (main) & 10 & haut $\rightarrow$ bas \\
\hline
\zh{家} jiā, famille & \zh{宀} toit & 10 & haut $\rightarrow$ bas ; point, point, crochet, puis \zh{豕} \\
\hline
\zh{帮} bāng, aider & \zh{巾} tissu & 9 & haut $\rightarrow$ bas : \zh{邦} puis \zh{巾} \\
\hline
\zh{团} tuán, groupe & \zh{囗} enceinte & 6 & extérieur $\rightarrow$ intérieur $\rightarrow$ fermeture \\
\hline
\zh{暖} nuǎn, chaud & \zh{日} soleil & 13 & gauche $\rightarrow$ droite \\
\hline
\end{tabularx}
\end{center}

\courssub{Méthode et pièges}

\begin{methode}[Comment mémoriser un caractère en trois étapes]
\begin{enumerate}
\item \textbf{Identifier le radical} : il donne le sens général. Pour \zh{暖}, le soleil \zh{日} annonce la chaleur ; pour \zh{家}, le toit \zh{宀} annonce la maison.
\item \textbf{Identifier le reste} (souvent le composant phonétique ou une image) : \zh{友} dans \zh{爱}, \zh{豕} dans \zh{家}. Fabriquez une petite histoire : « la main saisit l'ami ».
\item \textbf{Écrire en comptant les traits} : tracez le caractère lentement, à voix haute, en comptant « un, deux, trois… » jusqu'au nombre exact (10 pour \zh{爱}, 13 pour \zh{暖}). Si le compte n'y est pas, un trait a été oublié ou inventé.
\end{enumerate}
\end{methode}

\begin{attention}[Ne confondez pas \zh{爱} et \zh{受} !]
Les francophones confondent souvent \zh{爱} (ài, aimer) et \zh{受} (shòu, recevoir) : même silhouette générale, même griffe \zh{爫} en haut, même couvercle \zh{冖}. La différence est en bas : \zh{爱} contient \zh{友} (l'ami) tandis que \zh{受} contient \zh{又} (la main droite). Comparez : \zh{我爱我家。} (Wǒ ài wǒ jiā. « J'aime ma famille. ») et \zh{我受到帮助。} (Wǒ shòudào bāngzhù. « Je reçois de l'aide. »). À l'examen, un trait de trop ou de moins change le sens de toute la phrase !
\end{attention}

\begin{exemplebox}[Un proverbe à retenir pour l'examen]
\zh{团结就是力量。}\\
\emph{Tuánjié jiùshì lìliang.}\\
« L'union fait la force. »\\
Ce proverbe illustre parfaitement notre thème de solidarité : \zh{团} (enceinte qui réunit) + \zh{结} (nouer) = s'unir. Employez-le dans vos productions écrites sur la famille et l'entraide.
\end{exemplebox}

\begin{exoresolu}[Compter les traits et identifier les radicaux]
Énoncé. Pour chaque caractère, donne le radical, le nombre de traits et la règle d'ordre appliquée : \zh{爱}, \zh{团}, \zh{暖}.
\tcblower
Solution détaillée.
\begin{enumerate}
\item \zh{爱} : radical \zh{爫} (griffe) ; 10 traits ; règle de haut en bas (\zh{爫} puis \zh{冖} puis \zh{友}).
\item \zh{团} : radical \zh{囗} (enceinte) ; 6 traits ; règle extérieur $\rightarrow$ intérieur $\rightarrow$ fermeture : l'enceinte s'ouvre, \zh{才} s'écrit dedans, le trait du bas ferme.
\item \zh{暖} : radical \zh{日} (soleil) à gauche ; 13 traits ; règle de gauche à droite : \zh{日} d'abord, partie droite ensuite.
\end{enumerate}
\end{exoresolu}

\begin{exercice}[Atelier d'écriture]
Sur ton cahier, trace chaque caractère \textbf{5 fois} en comptant les traits à voix haute : \zh{爱} (10), \zh{家} (10), \zh{帮} (9), \zh{团} (6), \zh{暖} (13). Puis recopie la phrase \zh{我爱我家，团结就是力量。} (Wǒ ài wǒ jiā, tuánjié jiùshì lìliang. « J'aime ma famille, l'union fait la force. ») en veillant à l'ordre des traits.
\end{exercice}

\begin{corrige}[Atelier d'écriture]
Vérification par le compte : si tu arrives à 10 pour \zh{爱} et \zh{家}, 9 pour \zh{帮}, 6 pour \zh{团} et 13 pour \zh{暖}, ton écriture est complète. Contrôle visuel : \zh{爱} finit par \zh{友} (et non \zh{又}) ; \zh{家} commence par le point du toit ; \zh{团} se termine par le trait horizontal qui ferme l'enceinte ; \zh{暖} commence par le soleil \zh{日} bien étroit à gauche.
\end{corrige}

\begin{aretenir}[L'essentiel de la section]
Trois règles d'ordre : haut $\rightarrow$ bas, gauche $\rightarrow$ droite, extérieur $\rightarrow$ intérieur puis fermeture. Cinq caractères du thème : \zh{爱} (10 traits, radical \zh{爫}), \zh{家} (10 traits, radical \zh{宀}), \zh{帮} (9 traits, radical \zh{巾}), \zh{团} (6 traits, radical \zh{囗}), \zh{暖} (13 traits, radical \zh{日}). Le radical donne le sens ; compter les traits garantit une écriture exacte.
\end{aretenir}

\coursec{Compréhension du texte}

Dans la section précédente, nous avons appris à écrire les caractères du thème en respectant les clés (radicaux) et l'ordre des traits : \zh{家} (clé du toit 宀), \zh{爱} (clé de la griffe 爫 et du cœur 心), \zh{帮} (clé de la serviette 巾). Nous savons donc maintenant \emph{lire} et \emph{écrire} les mots. Il reste l'étape la plus importante : \cle{comprendre} le texte dans son ensemble, c'est-à-dire être capable de répondre à des questions précises en retrouvant l'information dans le texte et en la reformulant. C'est exactement la compétence évaluée à l'examen : « Lire et comprendre un court texte chinois, puis répondre à des questions de compréhension ».

\courssub{Le texte support : rappel de lecture}

Avant de répondre aux questions, relisons attentivement le texte support de la leçon. Lis-le d'abord en silence, puis à voix haute en suivant le pinyin, puis vérifie ta compréhension avec la traduction.

\begin{textebox}[我的家 — Wǒ de jiā — « Ma famille »]
我叫李天，是喀麦隆人，住在雅温得。我家有五口人：爸爸、妈妈、哥哥、妹妹和我。\\
Wǒ jiào Lǐ Tiān, shì Kāmàilóng rén, zhù zài Yǎwēndé. Wǒ jiā yǒu wǔ kǒu rén : bàba, māma, gēge, mèimei hé wǒ.\\
« Je m'appelle Li Tian, je suis camerounais et j'habite à Yaoundé. Ma famille compte cinq personnes : papa, maman, mon grand frère, ma petite sœur et moi. »\\[4pt]
爸爸妈妈很爱我们。爸爸每天很早就去工作，晚上回家以后，常常帮妈妈做饭。妈妈很爱干净，也把家照顾得很好。\\
Bàba māma hěn ài wǒmen. Bàba měitiān hěn zǎo jiù qù gōngzuò, wǎnshang huí jiā yǐhòu, chángcháng bāng māma zuò fàn. Māma hěn ài gānjìng, yě bǎ jiā zhàogù de hěn hǎo.\\
« Papa et maman nous aiment beaucoup. Papa part travailler très tôt chaque jour ; le soir, quand il rentre à la maison, il aide souvent maman à faire la cuisine. Maman aime beaucoup la propreté et prend aussi très bien soin de la maison. »\\[4pt]
我们孩子也爱爸爸妈妈。哥哥帮妹妹学习中文，我帮妈妈买菜，妹妹帮爸爸拿东西。周末，全家人一起吃饭、聊天、看电视。星期天，我们有时候去看爷爷奶奶。\\
Wǒmen háizi yě ài bàba māma. Gēge bāng mèimei xuéxí Zhōngwén, wǒ bāng māma mǎi cài, mèimei bāng bàba ná dōngxi. Zhōumò, quán jiā rén yìqǐ chī fàn, liáo tiān, kàn diànshì. Xīngqītiān, wǒmen yǒu shíhou qù kàn yéye nǎinai.\\
« Nous, les enfants, nous aimons aussi papa et maman. Mon grand frère aide ma petite sœur à apprendre le chinois, j'aide maman à acheter les légumes, et ma petite sœur aide papa à porter les affaires. Le week-end, toute la famille mange, bavarde et regarde la télévision ensemble. Le dimanche, nous allons parfois voir grand-père et grand-mère. »\\[4pt]
我觉得家是最温暖的地方，因为家人互相关心、互相帮助。\\
Wǒ juéde jiā shì zuì wēnnuǎn de dìfang, yīnwèi jiārén hùxiāng guānxīn, hùxiāng bāngzhù.\\
« Je pense que la maison est l'endroit le plus chaleureux, parce que les membres de la famille prennent soin les uns des autres et s'entraident. »
\end{textebox}

\begin{center}
\begin{tabularx}{\linewidth}{|X|X|X|X|}
\hline
\rowcolor{popblueL} Caractères & Pinyin & Nature & Traduction \\
\hline
家人 & jiārén & nom & membres de la famille \\
\hline
互相 & hùxiāng & adverbe & mutuellement, l'un l'autre \\
\hline
关心 & guānxīn & verbe & se soucier de, prendre soin de \\
\hline
帮助 & bāngzhù & verbe & aider \\
\hline
温暖 & wēnnuǎn & adjectif & chaleureux \\
\hline
照顾 & zhàogù & verbe & s'occuper de \\
\hline
觉得 & juéde & verbe & penser, trouver (opinion) \\
\hline
地方 & dìfang & nom & endroit, lieu \\
\hline
\end{tabularx}
\end{center}

\begin{attention}[Deux points de langue à remarquer dans le texte]
\begin{itemize}
\item \zh{我家有五口人} : pour compter les membres d'une famille, le chinois utilise le classificateur spécial \zh{口} (kǒu, litt. « bouche ») et non le classificateur général \zh{个}. On dit \zh{五口人} « cinq personnes (dans la famille) », jamais \zh{五个人} dans ce contexte familial.
\item \zh{最温暖的地方} : \zh{最} (zuì) marque le superlatif, « le plus ». Structure : \textbf{最 + adjectif}. Exemple : \zh{雅温得是喀麦隆最大的城市之一。} \emph{Yǎwēndé shì Kāmàilóng zuì dà de chéngshì zhīyī.} « Yaoundé est l'une des plus grandes villes du Cameroun. »
\end{itemize}
\end{attention}

\courssub{Méthode : comment répondre à une question de compréhension}

Répondre à une question de compréhension en chinois n'est pas une affaire d'inspiration : c'est une \cle{technique} en trois étapes. Le texte contient toujours la réponse ; ton travail est de la \emph{repérer}, de la \emph{citer}, puis de la \emph{reformuler}.

\begin{methode}[Répondre à une question de compréhension en trois étapes]
\begin{enumerate}
\item \textbf{Repérer le mot-clé de la question.} Souligne le mot important de la question (un verbe, un nom, un mot interrogatif comme \zh{谁} « qui », \zh{什么} « quoi », \zh{什么时候} « quand », \zh{为什么} « pourquoi »), puis cherche ce mot (ou un synonyme) dans le texte.
\item \textbf{Recopier la phrase concernée.} Une fois le mot-clé trouvé, recopie la phrase complète du texte qui contient la réponse. C'est ta \emph{justification} : à l'examen, une réponse non justifiée par le texte perd des points.
\item \textbf{Reformuler en phrase complète.} Transforme la phrase du texte en réponse directe à la question, en adaptant les pronoms (\zh{我} « je » devient \zh{他} « il » si la question porte sur le narrateur, etc.). Une réponse en un seul mot est acceptée seulement pour les questions fermées (voir la boîte Attention plus bas).
\end{enumerate}
\end{methode}

\begin{exemplebox}[Application de la méthode, étape par étape]
Question : \zh{他们周末做什么？} \emph{Tāmen zhōumò zuò shénme ?} « Que font-ils le week-end ? »\\[4pt]
\textbf{Étape 1 — mot-clé :} \zh{周末} « week-end ». On le repère dans le texte : \zh{周末，全家人一起吃饭、聊天、看电视。}\\
\textbf{Étape 2 — citation :} on recopie cette phrase.\\
\textbf{Étape 3 — reformulation :} \zh{他们周末一起吃饭、聊天、看电视。} \emph{Tāmen zhōumò yìqǐ chī fàn, liáo tiān, kàn diànshì.} « Le week-end, ils mangent, bavardent et regardent la télévision ensemble. »\\[4pt]
Autre exemple : \zh{李天觉得家怎么样？} \emph{Lǐ Tiān juéde jiā zěnmeyàng ?} « Que pense Li Tian de la maison ? »\\
Réponse : \zh{他觉得家是最温暖的地方。} \emph{Tā juéde jiā shì zuì wēnnuǎn de dìfang.} « Il pense que la maison est l'endroit le plus chaleureux. » (Le pronom \zh{我} du texte devient \zh{他} dans la réponse.)
\end{exemplebox}

\courssub{Compréhension globale : de qui parle-t-on ? De quoi parle-t-on ?}

Avant les questions de détail, l'examen commence toujours par deux questions \emph{globales}. Pour y répondre, on ne cherche pas une phrase précise : on regarde le texte entier.

\begin{definition}[Compréhension globale]
La \cle{compréhension globale} consiste à identifier le \textbf{sujet} du texte (de qui parle-t-on ?) et son \textbf{thème principal} (de quoi parle-t-on ?) en une ou deux phrases, sans entrer dans les détails. On s'appuie sur le titre, la première phrase et la dernière phrase, qui résument souvent l'idée principale.
\end{definition}

\begin{enumerate}
\item \textbf{De qui parle le texte ?} Le texte parle de \zh{李天} (Lǐ Tiān), un élève camerounais de Yaoundé, et de sa famille : ses parents, son grand frère et sa petite sœur.
\item \textbf{Quel est le thème principal ?} Le thème est \cle{l'amour de la famille et la solidarité familiale} : \zh{家人互相关心、互相帮助} « les membres de la famille prennent soin les uns des autres et s'entraident ». La dernière phrase du texte donne directement cette idée.
\end{enumerate}

\begin{popfigure}
\begin{tikzpicture}[
node distance=0.55cm and 0.9cm,
q/.style={rectangle, rounded corners=3pt, draw=popblue, fill=popblueL, align=center, font=\small, inner sep=5pt},
a/.style={rectangle, rounded corners=3pt, draw=popgreen, fill=popgreenL, align=center, font=\small, inner sep=5pt},
fleche/.style={-{Stealth[length=2.5mm]}, thick, popdark}]
\node[q, align=center] (qui){谁 ?\\ Qui ?};
\node[q, right=of qui, align=center] (quoi){做什么 ?\\ Quoi ?};
\node[q, right=of quoi, align=center] (quand){什么时候 ?\\ Quand ?};
\node[q, right=of quand, align=center] (pourquoi){为什么 ?\\ Pourquoi ?};
\node[a, below=of qui, align=center] (r1){李天一家\\ Lǐ Tiān et sa famille};
\node[a, below=of quoi, align=center] (r2){互相关心、互相帮助\\ s'entraider};
\node[a, below=of quand, align=center] (r3){周末、星期天\\ le week-end};
\node[a, below=of pourquoi, align=center] (r4){家是最温暖的地方\\ la maison est chaleureuse};
\draw[fleche] (qui) -- (r1);
\draw[fleche] (quoi) -- (r2);
\draw[fleche] (quand) -- (r3);
\draw[fleche] (pourquoi) -- (r4);
\end{tikzpicture}
\legende{Organigramme de compréhension : quatre mots interrogatifs pour interroger le texte.}
\end{popfigure}

\courssub{Compréhension fine : retrouver les détails}

La \cle{compréhension fine} porte sur les détails : les actions de chaque personne, les moments, les lieux. Applique la méthode en trois étapes.

\textbf{Que font les parents ?}
\begin{itemize}
\item \zh{爸爸每天很早就去工作。} \emph{Bàba měitiān hěn zǎo jiù qù gōngzuò.} « Papa part travailler très tôt chaque jour. »
\item \zh{晚上回家以后，常常帮妈妈做饭。} \emph{Wǎnshang huí jiā yǐhòu, chángcháng bāng māma zuò fàn.} « Le soir, en rentrant, il aide souvent maman à faire la cuisine. »
\item \zh{妈妈把家照顾得很好。} \emph{Māma bǎ jiā zhàogù de hěn hǎo.} « Maman s'occupe très bien de la maison. »
\end{itemize}

\textbf{Que font les enfants ?}
\begin{itemize}
\item \zh{哥哥帮妹妹学习中文。} \emph{Gēge bāng mèimei xuéxí Zhōngwén.} « Le grand frère aide la petite sœur à apprendre le chinois. »
\item \zh{我帮妈妈买菜。} \emph{Wǒ bāng māma mǎi cài.} « Le narrateur aide maman à acheter les légumes. »
\item \zh{妹妹帮爸爸拿东西。} \emph{Mèimei bāng bàba ná dōngxi.} « La petite sœur aide papa à porter les affaires. »
\end{itemize}

\textbf{Quand la famille se réunit-elle ?}
\begin{itemize}
\item \zh{周末，全家人一起吃饭、聊天、看电视。} \emph{Zhōumò, quán jiā rén yìqǐ chī fàn, liáo tiān, kàn diànshì.} « Le week-end, toute la famille se réunit pour manger, bavarder et regarder la télévision. »
\item \zh{星期天，我们有时候去看爷爷奶奶。} \emph{Xīngqītiān, wǒmen yǒu shíhou qù kàn yéye nǎinai.} « Le dimanche, ils vont parfois voir les grands-parents paternels. »
\end{itemize}

\begin{center}
\begin{tabularx}{\linewidth}{|X|X|X|}
\hline
\rowcolor{popblueL} Mot interrogatif & Question & Réponse dans le texte \\
\hline
谁 shéi « qui » & 谁很爱孩子们？ & 爸爸妈妈很爱他们。 \\
\hline
什么 shénme « quoi » & 哥哥帮妹妹做什么？ & 帮妹妹学习中文。 \\
\hline
什么时候 shénme shíhou « quand » & 全家人什么时候一起吃饭？ & 周末。 \\
\hline
怎么样 zěnmeyàng « comment » & 李天觉得家怎么样？ & 家是最温暖的地方。 \\
\hline
为什么 wèishénme « pourquoi » & 为什么家很温暖？ & 因为家人互相关心、互相帮助。 \\
\hline
\end{tabularx}
\end{center}

\begin{attention}[Le mot 谁 : deux prononciations acceptées]
Le mot interrogatif \zh{谁} « qui » se prononce couramment \emph{shéi} dans la langue parlée ; la prononciation \emph{shuí} est la forme littéraire et reste correcte. À l'oral comme à l'examen, les deux sont admises, mais retiens \emph{shéi} pour la conversation. Attention aussi à \zh{什么} : le second caractère \zh{么} porte un \emph{ton léger} (me), pas un ton plein.
\end{attention}

\courssub{Répondre en chinois : les réponses courtes}

\begin{propriete}[La réponse courte chinoise : on reprend le verbe de la question]
En français, on répond « oui » ou « non ». En chinois, il n'existe pas de mot unique pour « oui » : on \cle{reprend le verbe (ou l'adjectif) de la question}, éventuellement précédé de \zh{很} ou suivi d'un complément. Pour répondre négativement, on place \zh{不} devant le verbe repris.\\[4pt]
\zh{你爱家人吗？} \emph{Nǐ ài jiārén ma ?} « Aimes-tu ta famille ? »\\
→ \zh{爱。} \emph{Ài.} « Oui. » (litt. « J'aime. »)\\
→ \zh{我很爱他们。} \emph{Wǒ hěn ài tāmen.} « Oui, je les aime beaucoup. »\\
→ \zh{不爱。} \emph{Bú ài.} « Non. »\\[4pt]
\zh{爸爸妈妈很爱你吗？} \emph{Bàba māma hěn ài nǐ ma ?} « Tes parents t'aiment-ils beaucoup ? »\\
→ \zh{爸爸妈妈很爱我。} \emph{Bàba māma hěn ài wǒ.} « Oui, papa et maman m'aiment beaucoup. »
\end{propriete}

\begin{attention}[Pièges fréquents des francophones]
\begin{itemize}
\item Ne traduis jamais « oui » par un mot inventé : \zh{是} seul ne répond qu'aux questions contenant \zh{是} (\zh{你是学生吗？} → \zh{是。}). Pour \zh{你爱家人吗？}, la réponse est \zh{爱}, pas \zh{是}.
\item N'oublie pas d'\emph{adapter les pronoms} quand tu reformules : le texte dit \zh{我爱爸爸妈妈}, mais si la question porte sur \zh{李天}, ta réponse doit dire \zh{他爱爸爸妈妈}.
\item Dans une réponse complète, garde l'ordre chinois : Sujet + temps/lieu + verbe. On dit \zh{我们星期天去看爷爷奶奶} et non \zh{我们去看爷爷奶奶星期天}.
\item Le mot \zh{互相} « mutuellement » se place \emph{avant} le verbe : \zh{家人互相帮助}, jamais après.
\item Sandhi tonal de \zh{不} : devant un quatrième ton, \zh{不} (bù) devient \emph{bú} : \zh{不爱} se prononce \emph{bú ài}. À l'écrit, on note la prononciation réelle dans le pinyin.
\end{itemize}
\end{attention}

\courssub{Vrai ou faux justifié}

Pour les questions \emph{vrai/faux}, la justification est obligatoire : cite la phrase exacte du texte qui prouve ton jugement.

\begin{exoresolu}[Vrai ou faux ? Justifie par une citation du texte]
Énoncé. Indique si chaque phrase est vraie (对 duì) ou fausse (错 cuò), puis cite la phrase du texte qui le prouve.\\
1. 李天家有六口人。 \emph{Lǐ Tiān jiā yǒu liù kǒu rén.} « La famille de Li Tian compte six personnes. »\\
2. 爸爸晚上帮妈妈做饭。 \emph{Bàba wǎnshang bāng māma zuò fàn.} « Le soir, papa aide maman à cuisiner. »\\
3. 妹妹帮哥哥学习中文。 \emph{Mèimei bāng gēge xuéxí Zhōngwén.} « La petite sœur aide le grand frère à apprendre le chinois. »
\tcblower
\textbf{Solution détaillée.}\\
1. \textbf{错 (faux).} Le texte dit : \zh{我家有五口人} « ma famille compte cinq personnes » (papa, maman, le grand frère, la petite sœur et le narrateur), pas six.\\
2. \textbf{对 (vrai).} Le texte dit : \zh{晚上回家以后，常常帮妈妈做饭} « le soir, en rentrant à la maison, il aide souvent maman à faire la cuisine ».\\
3. \textbf{错 (faux).} Le texte dit : \zh{哥哥帮妹妹学习中文} « c'est le grand frère qui aide la petite sœur à apprendre le chinois », et non l'inverse. Attention à l'ordre des rôles : \zh{哥哥} est l'aîné, \zh{妹妹} la cadette.
\end{exoresolu}

\courssub{Question d'opinion : à toi de parler}

La dernière étape de la compréhension est la \cle{question d'opinion} : elle te fait quitter le texte pour donner ton avis, avec le verbe \zh{觉得} « penser, trouver » et la justification introduite par \zh{因为} « parce que ».

\begin{exemplebox}[Modèle de réponse d'opinion]
Question : \zh{你觉得家是最温暖的地方吗？为什么？} \emph{Nǐ juéde jiā shì zuì wēnnuǎn de dìfang ma ? Wèishénme ?} « Penses-tu que la maison soit l'endroit le plus chaleureux ? Pourquoi ? »\\[4pt]
Réponse modèle 1 (d'accord) : \zh{我觉得家是最温暖的地方，因为家人互相关心、互相帮助。} \emph{Wǒ juéde jiā shì zuì wēnnuǎn de dìfang, yīnwèi jiārén hùxiāng guānxīn, hùxiāng bāngzhù.} « Je pense que la maison est l'endroit le plus chaleureux, parce que les membres de la famille prennent soin les uns des autres et s'entraident. »\\[4pt]
Réponse modèle 2 (nuancée) : \zh{我觉得家很温暖，可是学校也很好，因为学校里有朋友。} \emph{Wǒ juéde jiā hěn wēnnuǎn, kěshì xuéxiào yě hěn hǎo, yīnwèi xuéxiào li yǒu péngyou.} « Je trouve la maison très chaleureuse, mais l'école est bien aussi, parce qu'à l'école il y a des amis. »
\end{exemplebox}

Structure à retenir : \textbf{Sujet + 觉得 + opinion，因为 + justification.} Cette structure est réutilisable dans toutes les questions d'opinion de l'année.

\courssub{Entraînement et corrigé collectif}

\begin{exercice}[Questions de compréhension sur le texte « 我的家 »]
Réponds en chinois, par des phrases complètes, en citant le texte.\\
1. 李天是哪儿人？他住在哪儿？ \emph{Lǐ Tiān shì nǎr rén ? Tā zhù zài nǎr ?}\\
2. 谁很爱孩子们？ \emph{Shéi hěn ài háizimen ?}\\
3. 爸爸每天什么时候去工作？ \emph{Bàba měitiān shénme shíhou qù gōngzuò ?}\\
4. 孩子们帮家人做什么？（说两个例子） \emph{Háizimen bāng jiārén zuò shénme ?}\\
5. 全家人周末做什么？ \emph{Quán jiā rén zhōumò zuò shénme ?}\\
6. Vrai ou faux : 李天星期天常常去看爷爷奶奶。 Justifie par une citation.\\
7. 你觉得家是最温暖的地方吗？为什么？
\end{exercice}

\begin{corrige}[Exercice « Questions de compréhension »]
1. \zh{李天是喀麦隆人，住在雅温得。} \emph{Lǐ Tiān shì Kāmàilóng rén, zhù zài Yǎwēndé.} « Li Tian est camerounais et habite à Yaoundé. » (Citation : \zh{是喀麦隆人，住在雅温得}.)\\
2. \zh{爸爸妈妈很爱孩子们。} \emph{Bàba māma hěn ài háizimen.} « Papa et maman aiment beaucoup les enfants. » (Citation : \zh{爸爸妈妈很爱我们} ; on adapte \zh{我们} en \zh{孩子们}.)\\
3. \zh{爸爸每天很早就去工作。} \emph{Bàba měitiān hěn zǎo jiù qù gōngzuò.} « Papa part travailler très tôt chaque jour. »\\
4. Deux exemples possibles : \zh{哥哥帮妹妹学习中文，我帮妈妈买菜。} \emph{Gēge bāng mèimei xuéxí Zhōngwén, wǒ bāng māma mǎi cài.} « Le grand frère aide la petite sœur à apprendre le chinois, et le narrateur aide maman à acheter les légumes. » (On accepte aussi \zh{妹妹帮爸爸拿东西}.)\\
5. \zh{周末，全家人一起吃饭、聊天、看电视。} \emph{Zhōumò, quán jiā rén yìqǐ chī fàn, liáo tiān, kàn diànshì.} « Le week-end, toute la famille mange, bavarde et regarde la télévision ensemble. »\\
6. \textbf{错 (faux).} Le texte dit \zh{星期天，我们有时候去看爷爷奶奶} : ils y vont \emph{parfois} (\zh{有时候}), pas \emph{souvent} (\zh{常常}). Piège classique : un seul adverbe change tout le sens.\\
7. Réponse libre, avec la structure \zh{我觉得……因为……}. Exemple : \zh{我觉得家是最温暖的地方，因为爸爸妈妈很爱我，我也很爱他们。} \emph{Wǒ juéde jiā shì zuì wēnnuǎn de dìfang, yīnwèi bàba māma hěn ài wǒ, wǒ yě hěn ài tāmen.} « Je pense que la maison est l'endroit le plus chaleureux, parce que papa et maman m'aiment beaucoup, et moi aussi je les aime beaucoup. »
\end{corrige}

\begin{aretenir}[L'essentiel de la section]
\begin{itemize}
\item Méthode en trois étapes : \cle{repérer} le mot-clé → \cle{citer} la phrase du texte → \cle{reformuler} en adaptant les pronoms.
\item Les mots interrogatifs guident la recherche : \zh{谁} qui, \zh{什么} quoi, \zh{什么时候} quand, \zh{怎么样} comment, \zh{为什么} pourquoi.
\item Réponse courte : on \cle{reprend le verbe} de la question (\zh{你爱家人吗？} → \zh{爱。}) ; la négation se fait avec \zh{不}.
\item Vrai/faux : la justification par une \cle{citation exacte} du texte est obligatoire ; attention aux adverbes (\zh{常常} ≠ \zh{有时候}).
\item Question d'opinion : \zh{我觉得 + opinion，因为 + justification}.
\item Points de langue du texte : classificateur \zh{口} pour les membres de la famille ; superlatif avec \zh{最} + adjectif.
\end{itemize}
\end{aretenir}

La compréhension du texte est maintenant assurée : tu sais identifier le thème, retrouver les détails, répondre en chinois et justifier tes réponses. La section suivante te fera passer de la compréhension à la production : tu utiliseras les collocations du texte (\zh{帮妈妈做饭}, \zh{互相关心}…) pour parler et écrire à ton tour sur ta propre famille.

\coursec{Collocations et production guidée}

Dans la section précédente, vous avez vérifié votre compréhension du texte sur l'amour de la famille et la solidarité familiale. Vous savez maintenant lire et comprendre ce thème. Il reste une étape décisive : passer de la \emph{compréhension} à la \emph{production}. C'est l'objectif de cette dernière partie : mémoriser les combinaisons de mots naturelles du chinois (les \cle{collocations}), puis rédiger et présenter oralement un court paragraphe sur votre propre famille.

\courssub{Observer d'abord : les collocations du thème}

\begin{definition}[Qu'est-ce qu'une collocation ?]
Une \cle{collocation} est une combinaison habituelle et stable de deux ou plusieurs mots. En chinois comme en français, on ne combine pas les mots au hasard : on dit « prendre soin de » et non « faire soin de ». En chinois, on dit \zh{照顾弟妹} (zhàogù dìmèi) et non une autre combinaison. Apprendre les collocations par blocs, c'est parler un chinois naturel et gagner du temps à l'examen.
\end{definition}

Observez les phrases suivantes, tirées du texte étudié et de la vie quotidienne. Repérez les blocs verbaux qui reviennent.

\begin{exemplebox}[Les collocations en contexte]
\zh{我很爱我的家人。} \emph{Wǒ hěn ài wǒ de jiārén.} « J'aime beaucoup ma famille. »\\
\zh{我们一家人互相帮助。} \emph{Wǒmen yījiārén hùxiāng bāngzhù.} « Toute notre famille s'entraide. »\\
\zh{爸爸妈妈很关心我们。} \emph{Bàba māma hěn guānxīn wǒmen.} « Papa et maman se soucient beaucoup de nous. »\\
\zh{姐姐每天照顾弟妹。} \emph{Jiějie měitiān zhàogù dìmèi.} « La grande sœur prend soin des cadets chaque jour. »\\
\zh{晚上，我们一起吃饭。} \emph{Wǎnshang, wǒmen yíqǐ chīfàn.} « Le soir, nous mangeons ensemble. »\\
\zh{我有一个幸福的家庭。} \emph{Wǒ yǒu yí ge xìngfú de jiātíng.} « J'ai une famille heureuse. »
\end{exemplebox}

\begin{center}
\begin{tabularx}{\linewidth}{|X|X|X|X|}
\hline
\rowcolor{popblueL} Caractères & Pinyin & Nature & Traduction \\
\hline
爱家人 & ài jiārén & verbe + nom & aimer sa famille \\
\hline
互相帮助 & hùxiāng bāngzhù & adverbe + verbe & s'entraider, s'aider mutuellement \\
\hline
关心父母 & guānxīn fùmǔ & verbe + nom & se soucier de ses parents \\
\hline
照顾弟妹 & zhàogù dìmèi & verbe + nom & prendre soin de ses cadets \\
\hline
一起吃饭 & yíqǐ chīfàn & adverbe + verbe & manger ensemble \\
\hline
幸福的家庭 & xìngfú de jiātíng & adjectif + 的 + nom & une famille heureuse \\
\hline
\end{tabularx}
\end{center}

\begin{propriete}[Deux mots-outils à bien maîtriser dans ces collocations]
\begin{itemize}
\item \cle{互相} (hùxiāng, « mutuellement ») se place \textbf{avant le verbe} : Sujet + 互相 + Verbe. Exemple : \zh{我们互相帮助。} \emph{Wǒmen hùxiāng bāngzhù.} « Nous nous entraidons. » Il remplace le « se... l'un l'autre » du français.
\item \cle{一起} (yíqǐ, « ensemble ») se place aussi \textbf{avant le verbe} : Sujet + 一起 + Verbe. Exemple : \zh{我们一起吃饭。} \emph{Wǒmen yíqǐ chīfàn.} « Nous mangeons ensemble. »
\end{itemize}
En français, « mutuellement » et « ensemble » peuvent se placer après le verbe (« nous mangeons ensemble ») ; en chinois, jamais : l'adverbe précède toujours le verbe.
\end{propriete}

\begin{attention}[Erreurs fréquentes des francophones]
\begin{itemize}
\item Ne dites pas \zh{我们吃饭一起} : \zh{一起} doit précéder le verbe. Dites \zh{我们一起吃饭。}
\item \zh{互相} ne s'emploie qu'avec un sujet pluriel (au moins deux personnes). On ne peut pas dire \zh{我互相帮助} tout seul.
\item Devant un nom, l'adjectif à deux syllabes exige \zh{的} : \zh{幸福的家庭} (xìngfú de jiātíng), et non \zh{幸福家庭} à ce niveau.
\end{itemize}
\end{attention}

\courssub{Phrases modèles à trous : s'entraîner avant de rédiger}

Avant de produire librement, entraînez-vous avec des phrases à trous. Complétez mentalement, puis vérifiez.

\begin{exemplebox}[Modèles à trous]
1. \zh{我的家很\_\_\_。} \emph{Wǒ de jiā hěn \_\_\_.} (Ma famille est très \ldots) → réponses possibles : \zh{幸福} (xìngfú, heureuse), \zh{大} (dà, grande), \zh{温暖} (wēnnuǎn, chaleureuse).\\
2. \zh{我们互相\_\_\_。} \emph{Wǒmen hùxiāng \_\_\_.} (Nous \ldots mutuellement) → \zh{帮助} (bāngzhù), \zh{关心} (guānxīn), \zh{爱} (ài).\\
3. \zh{爸爸妈妈很\_\_\_我们。} \emph{Bàba māma hěn \_\_\_ wǒmen.} → \zh{关心} (guānxīn), \zh{爱} (ài).\\
4. \zh{晚上，我们一起\_\_\_。} \emph{Wǎnshang, wǒmen yíqǐ \_\_\_.} → \zh{吃饭} (chīfàn), \zh{学习} (xuéxí).
\end{exemplebox}

\begin{exoresolu}[Compléter une phrase à trous]
Énoncé : complétez la phrase \zh{姐姐每天\_\_\_弟妹。} avec la collocation correcte, puis traduisez.
\tcblower
Solution détaillée : le sens attendu est « prendre soin de ». La collocation apprise est \zh{照顾弟妹}. On obtient : \zh{姐姐每天照顾弟妹。} \emph{Jiějie měitiān zhàogù dìmèi.} « La grande sœur prend soin de ses cadets chaque jour. » Vérification : le verbe \zh{照顾} est bien placé entre le sujet \zh{姐姐} et le complément \zh{弟妹} ; l'adverbe de temps \zh{每天} précède le verbe, conformément à l'ordre des mots chinois.
\end{exoresolu}

\courssub{Production guidée : rédiger sur sa propre famille}

\begin{methode}[Rédiger un court paragraphe en quatre étapes]
\begin{enumerate}
\item \textbf{Présenter la famille} : donnez le nombre de personnes avec le classificateur \zh{口} : \zh{我家有五口人。} (Wǒ jiā yǒu wǔ kǒu rén.) « Ma famille compte cinq personnes. »
\item \textbf{Décrire les sentiments} : utilisez \zh{爱} et \zh{关心} : \zh{爸爸妈妈很关心我们。}
\item \textbf{Décrire les actions de solidarité} : utilisez \zh{互相帮助}, \zh{照顾}, \zh{一起} : \zh{我们互相帮助，一起吃饭。}
\item \textbf{Conclure par son ressenti} : utilisez \zh{我觉得} (wǒ juéde, « je trouve que ») : \zh{我觉得我的家很幸福。} (Wǒ juéde wǒ de jiā hěn xìngfú.) « Je trouve que ma famille est très heureuse. »
\end{enumerate}
\end{methode}

\begin{popfigure}
\begin{tikzpicture}[mindmap, grow cyclic, every node/.style={concept, align=center, font=\small},
root concept/.append style={concept color=popblue, font=\bfseries},
level 1/.append style={level distance=3.4cm, sibling angle=90},
level 2/.append style={level distance=2.6cm, sibling angle=50}]
\node[root concept, align=center]{我的家\\ \emph{Wǒ de jiā}}
  child[concept color=poporange] { node[align=center]{Population\\ combien ?\\ \zh{五口人}} }
  child[concept color=poppink] { node[align=center]{Amour / Souci\\ sentiments\\ \zh{很爱家人}} }
  child[concept color=popgreen] { node[align=center]{Aide / Unité\\ actions\\ \zh{互相帮助}} }
  child[concept color=popgold] { node[align=center]{Ressenti\\ conclusion\\ \zh{很幸福}} };
\end{tikzpicture}
\legende{Carte mentale de production : quatre branches pour construire son paragraphe sur la famille.}
\end{popfigure}

\begin{textebox}[Modèle à imiter]
我家有五口人：爸爸、妈妈、两个弟弟和我。\\
\emph{Wǒ jiā yǒu wǔ kǒu rén : bàba, māma, liǎng ge dìdi hé wǒ.}\\
« Ma famille compte cinq personnes : papa, maman, deux petits frères et moi. »\\
爸爸妈妈很关心我们，我们也很爱他们。\\
\emph{Bàba māma hěn guānxīn wǒmen, wǒmen yě hěn ài tāmen.}\\
« Papa et maman se soucient beaucoup de nous, et nous les aimons beaucoup aussi. »\\
我们互相关心，互相帮助。晚上，我们一起吃饭。\\
\emph{Wǒmen hùxiāng guānxīn, hùxiāng bāngzhù. Wǎnshang, wǒmen yíqǐ chīfàn.}\\
« Nous nous soucions les uns des autres et nous nous entraidons. Le soir, nous mangeons ensemble. »\\
我觉得我的家很幸福。\\
\emph{Wǒ juéde wǒ de jiā hěn xìngfú.}\\
« Je trouve que ma famille est très heureuse. »
\end{textebox}

\begin{attention}[Le classificateur des personnes de la famille : 口]
Pour compter les membres d'une famille, le chinois utilise le classificateur spécial \cle{口} (kǒu, littéralement « bouche » : autant de bouches à nourrir !) : \zh{三口人} (sān kǒu rén) « trois personnes », \zh{五口人} « cinq personnes ». Ne dites jamais \zh{我家有五个家人} dans ce modèle : la structure correcte et naturelle est \zh{我家有} + nombre + \zh{口} + \zh{人}. Notez aussi : devant un classificateur, « deux » se dit \zh{两} (liǎng) et non \zh{二} : \zh{两个弟弟} « deux petits frères ».
\end{attention}

\begin{center}
\begin{tabularx}{\linewidth}{|X|X|X|}
\hline
\rowcolor{popblueL} Structure & Exemple (caractères + pinyin) & Traduction \\
\hline
我家有 + nombre + 口人 & 我家有四口人。Wǒ jiā yǒu sì kǒu rén. & Ma famille compte quatre personnes. \\
\hline
Sujet + 很 + 关心/爱 + complément & 妈妈很爱我。Māma hěn ài wǒ. & Maman m'aime beaucoup. \\
\hline
Sujet pluriel + 互相 + verbe & 我们互相帮助。Wǒmen hùxiāng bāngzhù. & Nous nous entraidons. \\
\hline
Sujet + 一起 + verbe & 我们一起吃饭。Wǒmen yíqǐ chīfàn. & Nous mangeons ensemble. \\
\hline
我觉得 + phrase & 我觉得我的家很幸福。Wǒ juéde wǒ de jiā hěn xìngfú. & Je trouve que ma famille est heureuse. \\
\hline
\end{tabularx}
\end{center}

\courssub{Présentation orale et correction bienveillante}

\begin{experience}[À vous de parler !]
Par groupes, deux ou trois élèves passent au tableau et présentent leur paragraphe (3 à 4 phrases) sans le lire mot à mot. La classe écoute et relève : (1) les collocations correctement employées ; (2) les tons à retravailler. Le professeur corrige avec bienveillance, en priorité les tons qui changent le sens : \zh{爱} (ài, 4\textsuperscript{e} ton), \zh{关心} (guānxīn, 1\textsuperscript{er} ton), \zh{口} (kǒu, 3\textsuperscript{e} ton), \zh{幸福} (xìngfú, 4\textsuperscript{e} puis 2\textsuperscript{e} ton). Astuce : frappez doucement le rythme des tons avec la main en parlant.
\end{experience}

\begin{savaistu}[Un mot qui valorise votre copie : 亲情]
Pour l'examen, enrichissez votre conclusion avec \cle{亲情} (qīnqíng), « l'affection familiale, les liens du cœur ». Une phrase comme \zh{我爱我的家，因为家里有亲情。} \emph{Wǒ ài wǒ de jiā, yīnwèi jiā lǐ yǒu qīnqíng.} « J'aime ma famille, parce qu'il y a de l'affection familiale à la maison » montre un vocabulaire riche et est très appréciée des correcteurs. En Chine comme au Cameroun, la solidarité familiale est une valeur centrale : savoir l'exprimer en chinois, c'est aussi savoir parler de votre propre culture.
\end{savaistu}

\begin{exercice}[Production écrite guidée]
Rédigez 3 à 4 phrases sur votre famille en suivant la méthode en quatre étapes. Utilisez au moins trois collocations de la leçon (\zh{爱家人}, \zh{互相帮助}, \zh{关心父母}, \zh{照顾弟妹}, \zh{一起吃饭}, \zh{幸福的家庭}) et le classificateur \zh{口}.
\end{exercice}

\begin{corrige}[Exercice : production écrite guidée]
Exemple de production attendue (à adapter à votre famille) :\\
\zh{我家有六口人：爸爸、妈妈、三个妹妹和我。} \emph{Wǒ jiā yǒu liù kǒu rén : bàba, māma, sān ge mèimei hé wǒ.} « Ma famille compte six personnes : papa, maman, trois petites sœurs et moi. »\\
\zh{父母很关心我们，我每天帮助妈妈照顾妹妹。} \emph{Fùmǔ hěn guānxīn wǒmen, wǒ měitiān bāngzhù māma zhàogù mèimei.} « Mes parents se soucient beaucoup de nous ; chaque jour j'aide maman à prendre soin de mes sœurs. »\\
\zh{我们一家人互相帮助，晚上一起吃饭。} \emph{Wǒmen yījiārén hùxiāng bāngzhù, wǎnshang yíqǐ chīfàn.} « Toute notre famille s'entraide, et le soir nous mangeons ensemble. »\\
\zh{我觉得我的家很幸福，我很爱我的家人。} \emph{Wǒ juéde wǒ de jiā hěn xìngfú, wǒ hěn ài wǒ de jiārén.} « Je trouve que ma famille est très heureuse, et j'aime beaucoup ma famille. »\\
Barème indicatif : classificateur \zh{口} correct (1 pt), trois collocations (3 pts), ordre des mots et tons respectés à l'oral (2 pts).
\end{corrige}

\begin{aretenir}[L'essentiel de la section]
\begin{itemize}
\item Les collocations s'apprennent par blocs : \zh{爱家人}, \zh{互相帮助}, \zh{关心父母}, \zh{照顾弟妹}, \zh{一起吃饭}, \zh{幸福的家庭}.
\item \zh{互相} et \zh{一起} se placent toujours \textbf{avant} le verbe.
\item Les membres de la famille se comptent avec \zh{口} : \zh{我家有五口人。}
\item Un bon paragraphe suit quatre étapes : présenter → sentiments → actions de solidarité → ressenti avec \zh{我觉得}.
\item Le mot \zh{亲情} (qīnqíng) enrichit une conclusion à l'examen.
\end{itemize}
\end{aretenir}

\newpage

\coursec{Activité d'intégration}

\coursec{Leçon 1 — Vocabulaire et compréhension de texte : Amour de la famille et solidarité familiale}

\courssub{Mise en route et découverte du thème}

\begin{experience}[Brainstorming et entrée en matière]
\textbf{Étape 1 — Image et proverbe.} Le professeur projette ou affiche une photo d'une famille camerounaise (parents, grands-parents, enfants) partageant un repas de \zh{ndolé} ou de \zh{poisson braisé} autour d'un grand plat commun, et une photo d'une famille chinoise réunie pour le \zh{年夜饭} (repas du réveillon du Nouvel An). 

\textbf{Étape 2 — Discussion guidée (oral).} En classe entière, le professeur pose les questions suivantes (les réponses attendues sont en français, le vocabulaire chinois est introduit au tableau) :
\begin{itemize}
\item Que voyez-vous sur ces photos ? (Famille, repas, partage, sourires)
\item Quel mot chinois connais-tu déjà pour « famille » ? (\zh{家} \emph{jiā}, \zh{家人} \emph{jiārén})
\item Comment dit-on « Je t'aime » à ses parents en chinois ? (\zh{我爱你} \emph{Wǒ ài nǐ})
\end{itemize}

\textbf{Étape 3 — Proverbe chinois du jour.} Le professeur écrit au tableau : \zh{家和万事兴} (\emph{Jiā hé wàn shì xīng}). 
\begin{itemize}
\item \textbf{Analyse} : \zh{家} (famille) + \zh{和} (harmonieux) + \zh{万事} (toutes les affaires / tout) + \zh{兴} (prospérer).
\item \textbf{Sens} : « Quand la famille vit en harmonie, tout prospère. »
\item \textbf{Parallèle culturel} : Au Cameroun, on dit souvent « La famille, c'est sacré » ou « Le sang ne devient pas eau ». En Chine, l'harmonie familiale (\zh{家和}) est considérée comme la racine (\zh{根本}) de la réussite sociale.
\end{itemize}
\end{experience}

\courssub{Texte support : lecture et traduction}

\begin{textebox}[Texte support : Mon portrait de famille — \zh{我的家庭画像}]
\zh{我的家不大，但是很温暖。}\\
\emph{Wǒ de jiā bù dà, dànshì hěn wēnnuǎn.}\\
« Ma famille n'est pas grande, mais elle est très chaleureuse. »\\[4pt]
\zh{我家有五口人：爸爸、妈妈、哥哥、妹妹和我。}\\
\emph{Wǒ jiā yǒu wǔ kǒu rén : bàba, māma, gēge, mèimei hé wǒ.}\\
« Il y a cinq personnes dans ma famille : papa, maman, mon grand frère, ma petite sœur et moi. »\\[4pt]
\zh{爸爸妈妈很爱我们，我们也爱他们。}\\
\emph{Bàba māma hěn ài wǒmen, wǒmen yě ài tāmen.}\\
« Papa et maman nous aiment beaucoup, et nous les aimons aussi. »\\[4pt]
\zh{在家里，我们互相帮助，互相关心。}\\
\emph{Zài jiā lǐ, wǒmen hùxiāng bāngzhù, hùxiāng guānxīn.}\\
« À la maison, nous nous aidons et nous prenons soin les uns des autres. »\\[4pt]
\zh{哥哥帮我学习中文，我帮妹妹做作业。}\\
\emph{Gēge bāng wǒ xuéxí Zhōngwén, wǒ bāng mèimei zuò zuòyè.}\\
« Mon grand frère m'aide à apprendre le chinois, et j'aide ma petite sœur à faire ses devoirs. »\\[4pt]
\zh{我爱我的家，家人是我的力量！}\\
\emph{Wǒ ài wǒ de jiā, jiārén shì wǒ de lìliang!}\\
« J'aime ma famille, ma famille est ma force ! »
\end{textebox}

\begin{methode}[Comment travailler ce texte en autonomie]
\begin{enumerate}
\item \textbf{Écoute globale} : Écoute l'audio (ou le professeur) sans regarder le texte. Combien de personnes dans la famille ? L'ambiance est-elle heureuse ?
\item \textbf{Lecture à voix haute} : Lis le texte en te concentrant sur les \textbf{tons} (surtout \zh{爱} ài 4\up{e}, \zh{家} jiā 1\up{er}, \zh{互相} hùxiāng 4-1). Enregistre-toi si possible.
\item \textbf{Découpage en unités de sens} : Repère les virgules chinoises \zh{，} et les points \zh{。}. Traduis phrase par phrase.
\item \textbf{Repérage grammatical} : Entoure le classificateur \zh{口}, souligne les deux \zh{互相}, cadre la structure \zh{帮 + qqn + verbe}.
\item \textbf{Récitation} : Essaie de redire le texte en ne gardant que les mots-clés au tableau.
\end{enumerate}
\end{methode}

\courssub{Glossaire du thème}

\begin{aretenir}[Vocabulaire de base — Famille et solidarité (HSK 1-3)]
\begin{center}
\begin{tabularx}{\linewidth}{|X|X|X|X|}
\hline
\rowcolor{popblueL} Caractères & Pinyin & Nature & Traduction \\
\hline
家 & jiā & n. & famille, maison, foyer \\
\hline
家人 & jiārén & n. & membres de la famille, proches \\
\hline
口 & kǒu & m. & classificateur pour les membres de la famille \\
\hline
爸爸 & bàba & n. & papa \\
\hline
妈妈 & māma & n. & maman \\
\hline
哥哥 & gēge & n. & grand frère \\
\hline
姐姐 & jiějie & n. & grande sœur \\
\hline
弟弟 & dìdi & n. & petit frère \\
\hline
妹妹 & mèimei & n. & petite sœur \\
\hline
爷爷 & yéye & n. & grand-père paternel \\
\hline
奶奶 & nǎinai & n. & grand-mère paternelle \\
\hline
爱 & ài & v. & aimer, amour \\
\hline
关心 & guānxīn & v. & se soucier de, prendre soin de \\
\hline
帮助 & bāngzhù & v./n. & aider ; aide, assistance \\
\hline
支持 & zhīchí & v. & soutenir, appuyer \\
\hline
团结 & tuánjié & adj./v. & uni, solidaire ; s'unir \\
\hline
温暖 & wēnnuǎn & adj. & chaleureux (climat, ambiance, cœur) \\
\hline
和睦 & hémù & adj. & harmonieux (relations familiales) \\
\hline
孝顺 & xiàoshun & adj. & filial, respectueux envers les parents \\
\hline
力量 & lìliang & n. & force, énergie, soutien moral \\
\hline
互相 & hùxiāng & adv. & mutuellement, l'un l'autre, réciproquement \\
\hline
相亲相爱 & xiāngqīnxiāng'ài & id. & s'aimer tendrement les uns les autres (famille) \\
\hline
\end{tabularx}
\end{center}
\end{aretenir}

\begin{aretenir}[Vocabulaire complémentaire — Liens de parenté élargis]
\begin{center}
\begin{tabularx}{\linewidth}{|X|X|X|X|}
\hline
\rowcolor{poporangeL} Caractères & Pinyin & Nature & Traduction \\
\hline
外公 & wàigōng & n. & grand-père maternel \\
\hline
外婆 & wàipó & n. & grand-mère maternelle \\
\hline
叔叔 & shūshu & n. & oncle (frère cadet du père) \\
\hline
阿姨 & āyí & n. & tante (sœur du père/mère) / nounou \\
\hline
堂哥 & tánggē & n. & cousin paternel aîné \\
\hline
表妹 & biǎomèi & n. & cousine maternelle cadette \\
\hline
\end{tabularx}
\end{center}
\end{aretenir}

\courssub{Clés (radicaux) et ordre des traits}

\begin{definition}[Rappel : Règles de base de l'ordre des traits]
\begin{enumerate}
\item De haut en haut (\zh{三} \emph{sān}) ; de gauche à droite (\zh{川} \emph{chuān}).
\item Horizontal avant vertical (\zh{十} \emph{shí}).
\item D'abord le trait central, puis les côtés symétriques (\zh{小} \emph{xiǎo}, \zh{水} \emph{shuǐ}).
\item Enveloppe avant contenu (\zh{国} \emph{guó}) ; fermer l'enveloppe en dernier (\zh{回} \emph{huí}).
\item Le trait qui traverse tout le caractère s'écrit en dernier (\zh{中} \emph{zhōng}, \zh{王} \emph{wáng}).
\end{enumerate}
\end{definition}

\begin{aretenir}[Analyse de 6 caractères clés de la leçon]
\begin{center}
\begin{tabularx}{\linewidth}{|X|X|X|X|X|}
\hline
\rowcolor{popgreenL} Caractère & Pinyin & Radical (Clé) & Nb traits & Décomposition / Mnémotechnique \\
\hline
\zh{家} & jiā & \zh{宀} (mián, toit) & 10 & \zh{宀} + \zh{豕} (shǐ, cochon). \emph{Sous le toit, le cochon = la richesse, la famille.} \\
\hline
\zh{爱} & ài & \zh{爫} (zhǎo, griffe) & 10 & \zh{爫} (main qui agit) + \zh{友} (ami) + \zh{心} (cœur). \emph{Aimer, c'est tendre la main vers un ami avec son cœur.} (Forme simplifiée) \\
\hline
\zh{互} & hù & \zh{彳} (chì, pas) & 4 & \zh{彳} + \zh{寸} (cùn, mesure/main). \emph{Aller vers l'autre pour échanger.} \\
\hline
\zh{助} & zhù & \zh{力} (lì, force) & 7 & \zh{且} (qiě, phonétique) + \zh{力}. \emph{Apporter sa force en plus.} \\
\hline
\zh{温} & wēn & \zh{氵} (shuǐ, eau) & 12 & \zh{氵} + \zh{皿} (mǐn, bol) + \zh{日} (rì, soleil). \emph{Le soleil réchauffe l'eau dans le bol = tiède.} \\
\hline
\zh{口} & kǒu & \zh{口} (kǒu, bouche) & 3 & Pictogramme : une bouche ouverte. \emph{Classificateur pour « bouches à nourrir » = personnes.} \\
\hline
\end{tabularx}
\end{center}
\end{aretenir}

\begin{methode}[Écriture guidée : \zh{家} (jiā) et \zh{爱} (ài)]
\textbf{\zh{家} (10 traits) :} 
1. Point \zh{丶} (haut gauche du toit) → 2. Horizontal court \zh{一} → 3. Horizontal long \zh{一} (fait le toit) → 4. Vertical \zh{丨} (descend du milieu) → 5. Horizontal \zh{一} (base du toit) → 
\textit{Intérieur (cochon \zh{豕}) :} 6. Horizontal \zh{一} → 7. Vertical \zh{丨} → 8. Horizontal \zh{一} → 9. Trait crochet \zh{乙} → 10. Point \zh{丶} (queue).

\textbf{\zh{爱} (10 traits, simplifié) :} 
1. Horizontal \zh{一} (haut \zh{爫}) → 2. Crochet vertical \zh{㇆} → 3. Point \zh{丶} → 4. Point \zh{丶} (fin \zh{爫}) → 
\textit{Milieu \zh{友} :} 5. Horizontal \zh{一} → 6. Vertical \zh{丨} → 7. Horizontal \zh{一} → 8. Crochet horizontal \zh{㇆} → 
\textit{Bas \zh{心} :} 9. Point \zh{丶} (gauche) → 10. Point \zh{丶} (milieu) → 11. Point \zh{丶} (droite) → 12. Trait oblique long \zh{㇏} (base). 
\emph{Note :} En simplifié, \zh{愛} (13 traits, trad.) perd le \zh{心} au milieu pour le mettre en bas, et \zh{友} remplace \zh{爫}+\zh{夂}.
\end{methode}

\begin{experience}[Atelier d'écriture : « Mon arbre généalogique en caractères »]
Sur une feuille A3, dessine un arbre. Sur le tronc, écris \zh{家}. Sur les grosses branches, écris les caractères des parents/grands-parents (\zh{父}, \zh{母}, \zh{爷}, \zh{奶}). Sur les petites branches, écris ton prénom chinois et ceux de tes frères/sœurs (\zh{哥}, \zh{姐}, \zh{弟}, \zh{妹}). Respecte l'ordre des traits. Affiche en classe.
\end{experience}

\courssub{Compréhension du texte}

\begin{exercice}[Questions sur le texte support « Mon portrait de famille »]
\textbf{Consigne :} Réponds en chinois (caractères + pinyin) aux questions suivantes, en te basant \textbf{uniquement} sur le texte.
\begin{enumerate}
\item \zh{作者家有几口人？} (\emph{Zuòzhě jiā yǒu jǐ kǒu rén?})
\item \zh{谁帮助作者学习中文？} (\emph{Shéi bāngzhù zuòzhě xuéxí Zhōngwén?})
\item \zh{作者帮助谁做作业？} (\emph{Zuòzhě bāngzhù shéi zuò zuòyè?})
\item \zh{作者一家人在家里做什么？} (Deux actions avec \zh{互相}) (\emph{...zài jiā lǐ zuò shénme?})
\item \zh{作者怎么形容自己的家？} (Deux adjectifs) (\emph{...zěnme xíngróng zìjǐ de jiā?})
\item \zh{翻译成法语 :} \zh{家人是我的力量。}
\end{enumerate}
\end{exercice}

\begin{corrige}[Exercice Compréhension]
\begin{enumerate}
\item \zh{作者家有五口人。} (\emph{Zuòzhě jiā yǒu wǔ kǒu rén.}) — Il y a 5 personnes dans la famille de l'auteur.
\item \zh{哥哥帮助作者学习中文。} (\emph{Gēge bāngzhù zuòzhě xuéxí Zhōngwén.}) — C'est le grand frère qui aide l'auteur à apprendre le chinois.
\item \zh{作者帮助妹妹做作业。} (\emph{Zuòzhě bāngzhù mèimei zuò zuòyè.}) — L'auteur aide sa petite sœur à faire ses devoirs.
\item \zh{互相帮助，互相关心。} (\emph{Hùxiāng bāngzhù, hùxiāng guānxīn.}) — Ils s'aident mutuellement et se soucient les uns des autres.
\item \zh{不大，但是很温暖。} (\emph{Bù dà, dànshì hěn wēnnuǎn.}) — Pas grande, mais très chaleureuse.
\item « Ma famille est ma force. » / « Les membres de ma famille sont ma force. »
\end{enumerate}
\end{corrige}

\begin{exercice}[Vrai ou Faux / 判断正误]
\textbf{Consigne :} Coche \zh{对} (juste) ou \zh{错} (faux). Corrige les phrases fausses en chinois.
\begin{enumerate}
\item \zh{作者家有六口人。} \hfill \zh{对 / 错}
\item \zh{妹妹帮助作者学习中文。} \hfill \zh{对 / 错}
\item \zh{在家里，他们互相支持。} \hfill \zh{对 / 错}
\item \zh{作者不爱他的家。} \hfill \zh{对 / 错}
\end{enumerate}
\end{exercice}

\begin{corrige}[Exercice Vrai/Faux]
\begin{enumerate}
\item \zh{错} → \zh{作者家有五口人。}
\item \zh{错} → \zh{哥哥帮助作者学习中文。}
\item \zh{错} (Le texte dit \zh{互相帮助，互相关心}, pas \zh{支持} explicitement, bien que ce soit implicite). Correction : \zh{在家里，他们互相帮助，互相关心。}
\item \zh{错} → \zh{作者很爱他的家。} / \zh{我爱我的家。}
\end{enumerate}
\end{corrige}

\courssub{Collocations et production guidée}

\begin{propriete}[Structures grammaticales essentielles du thème]
\begin{itemize}
\item \textbf{Compter les membres de la famille} \\
\cle{Sujet (我/你/他) + 家 + 有 + Nombre + 口 + 人} \\
\zh{我家有四口人。} (\emph{Wǒ jiā yǒu sì kǒu rén.}) — Ma famille compte 4 personnes. \\
\zh{你家有几口人？} (\emph{Nǐ jiā yǒu jǐ kǒu rén?}) — Combien de personnes chez toi ? \\
\emph{\textbf{Attention :}} \zh{口} (kǒu) est le classificateur \textbf{obligatoire} pour les humains dans le cadre familial. On ne dit pas \zh{五个人} pour sa propre famille.
\vspace{4pt}
\item \textbf{Exprimer l'aide : \zh{帮} (bāng) vs \zh{帮助} (bāngzhù)} \\
\cle{帮 + Quelqu'un + Verbe (action concrète, courte)} \\
\zh{妈妈帮我洗衣服。} (\emph{Māma bāng wǒ xǐ yīfu.}) — Maman m'aide à laver le linge. \\
\cle{帮助 + Quelqu'un + (做) + Verbe (plus formel, ou objet long)} \\
\zh{老师帮助我们学习汉字。} (\emph{Lǎoshī bāngzhù wǒmen xuéxí Hànzì.}) — Le prof nous aide à apprendre les caractères.
\vspace{4pt}
\item \textbf{Réciprocité : \zh{互相} (hùxiāng) + Verbe} \\
\cle{Sujet (pluriel) + 互相 + Verbe} \\
\zh{我们互相帮助。} (\emph{Wǒmen hùxiāng bāngzhù.}) — Nous nous aidons. \\
\zh{他们互相关心。} (\emph{Tāmen hùxiāng guānxīn.}) — Ils se soucient l'un de l'autre. \\
\emph{\textbf{Règle d'or :}} \zh{互相} se place \textbf{toujours avant le verbe}, jamais après. Sujet doit être pluriel (\zh{我们}, \zh{他们}, \zh{哥哥姐姐}).
\vspace{4pt}
\item \textbf{Idiome (Chengyu) figé : \zh{相亲相爱} (xiāng qīn xiāng ài)} \\
Structure : \zh{相} (mutuellement) + \zh{亲} (aimer/être proche) + \zh{相} + \zh{爱} (aimer). \\
\zh{一家人相亲相爱。} (\emph{Yì jiā rén xiāng qīn xiāng ài.}) — La famille s'aime tendrement. \\
\emph{Niveau :} HSK 4, mais incontournable pour ce thème.
\vspace{4pt}
\item \textbf{Complément de temps / lieu : Place avant le verbe} \\
\cle{Sujet + Temps / Lieu + Verbe + Objet} \\
\zh{在家里，我们互相帮助。} (\emph{Zài jiā lǐ, wǒmen hùxiāng bāngzhù.}) — À la maison, nous nous aidons. \\
\zh{困难的时候，我们互相支持。} (\emph{Kùnnan de shíhou, wǒmen hùxiāng zhīchí.}) — Dans les moments difficiles, nous nous soutenons.
\end{itemize}
\end{propriete}

\begin{attention}[Pièges fréquents des francophones — Check-list de révision]
\begin{itemize}
\item \textbf{Classificateur \zh{口}} : \zh{我家有三口人} (\checkmark) vs \zh{我家有三个人} (\textbf{✗} style parlé familier, à éviter à l'écrit/examen) vs \zh{三口人们} (\textbf{✗} \zh{们} interdit après nombre+classif.).
\item \textbf{Ordre \zh{互相}} : \zh{我们互相爱} (\checkmark) vs \zh{我们爱互相} (\textbf{✗}).
\item \textbf{\zh{帮} vs \zh{帮助}} : \zh{帮我做} (\checkmark, court) vs \zh{帮助我做作业} (\checkmark, plus long). \zh{帮助我做} (possible mais moins fluide).
\item \textbf{Tons critiques} : 
    \begin{itemize}
    \item \zh{爱} (ài, 4\up{e} ton descendant) $\neq$ \zh{艾} (ài, herbe) $\neq$ \zh{矮} (ǎi, petit).
    \item \zh{家} (jiā, 1\up{er} ton haut) $\neq$ \zh{加} (jiā, ajouter) $\neq$ \zh{夹} (jiā, presser).
    \item \zh{互} (hù, 4\up{e} ton) $\neq$ \zh{护} (hù, protéger).
    \end{itemize}
\item \textbf{\zh{和} (hé) vs \zh{跟} (gēn) vs \zh{与} (yǔ)} : Pour lister la famille : \zh{爸爸、妈妈、哥哥、妹妹和我} (\checkmark, \zh{和} lie le dernier). \zh{跟} plus oral (\zh{我跟哥哥去}).
\item \textbf{Adjectif prédicatif} : \zh{我家很温暖} (\checkmark, \zh{很} lie l'adjectif) vs \zh{我家是温暖} (\textbf{✗} \zh{是} + nom seulement).
\end{itemize}
\end{attention}

\begin{experience}[Tâche finale : Concours « Portrait de famille » — Semaine de la famille à Buea]
\textbf{Contexte :} Le club de chinois du lycée de Buea organise la « Semaine de la famille ». Concours d'affichage : « Mon portrait de famille ». Les 3 meilleurs sont affichés sous le titre \zh{家和万事兴}.

\textbf{Consignes détaillées (Travail individuel $\to$ Oral $\to$ Vote) :}
\begin{enumerate}
\item \textbf{Préparation (Brouillon - 15 min).} Remplis ta fiche :
    \begin{itemize}
    \item Nombre total : \zh{\_\_\_ 口人}
    \item Liste (ordre : aînés $\to$ cadets) : \zh{爸爸、妈妈、...、我}
    \item 2 actions de solidarité concrètes : \zh{谁帮谁？做什么？} (ex: \zh{姐姐帮我复习课文}, \zh{我帮爸爸洗车})
    \item Ton sentiment final : \zh{家人是我的...} (力量 / 依靠 / 骄傲 / 快乐)
    \end{itemize}
\item \textbf{Rédaction (15 min).} Rédige \textbf{5 à 6 phrases} en suivant ce plan :
    \begin{enumerate}
    \item Présentation : \zh{我的家...，我家有...口人：...}
    \item Ambiance/Sentiments : \zh{我们一家人... (相亲相爱 / 很和睦 / 很温暖)。}
    \item Actions solidarité (avec \zh{互相} ou \zh{帮}) : \zh{在家里，我们互相帮助：...} / \zh{哥哥帮我...，我帮妹妹...}
    \item Soutien moral (difficultés) : \zh{困难的时候，我们互相关心，互相支持。}
    \item Conclusion : \zh{我爱我的家，家人是我的力量！}
    \end{enumerate}
    \textbf{Contraintes obligatoires :} Minimum \textbf{6 mots du glossaire} (tableau p. 3) + Minimum \textbf{2 collocations} (\zh{互相帮助}, \zh{互相关心}, \zh{相亲相爱}, \zh{互相支持}).
\item \textbf{Illustration (5 min).} Dessine ou colle une photo. Légende : \zh{这是我的家人。} (\emph{Zhè shì wǒ de jiārén.}) ou \zh{这是我们全家福。} (\emph{Zhè shì wǒmen quánjiā fú.} — Photo de famille).
\item \textbf{Présentation orale (2 min/élève).} Présente sans lire mot à mot. Regarde la classe. Soigne les tons.
\item \textbf{Interaction (1 min).} Un camarade pose une question : \zh{你家有几口人？} / \zh{谁帮助你？} / \zh{你最爱谁？} Tu réponds en phrase complète.
\item \textbf{Vote \& Affichage.} Grille d'évaluation par les pairs (voir ci-dessous).
\end{enumerate}
\end{experience}

\begin{aretenir}[Grille d'auto-évaluation et d'évaluation par les pairs (Sur 10 pts)]
\begin{center}
\begin{tabularx}{\linewidth}{|X|X|X|}
\hline
\rowcolor{poppurpleL} Critère & Indicateurs de réussite & Points \\
\hline
Lexique & $\ge$ 6 mots du glossaire employés correctement & /2 \\
\hline
Collocations & $\ge$ 2 collocations (\zh{互相...}, \zh{相亲相爱}) bien placées & /2 \\
\hline
Grammaire & \zh{口} correct, ordre mots (Temps/Lieu + V), \zh{帮/帮助} structure & /2 \\
\hline
Prononciation & Tons clairs (surtout \zh{爱}, \zh{家}, \zh{互相}), fluidité & /2 \\
\hline
Interaction & Réponse correcte et complète à la question du camarade & /2 \\
\hline
\end{tabularx}
\end{center}
\textbf{Seuils :} 8-10 = \zh{优秀} (Affiché au mur) \quad 6-7 = \zh{良好} (À enrichir) \quad $<$ 6 = \zh{需加强} (Retravail avec prof).
\end{aretenir}

\begin{exoresolu}[Production modèle — Élève de Première A (Buea)]
Énoncé : Rédige ton portrait de famille (5-6 phrases, $\ge$ 6 mots glossaire, $\ge$ 2 collocations).
\tcblower
\textbf{Texte de l'élève :}\\[4pt]
\zh{我的家很大，也很温暖。}\\
\emph{Wǒ de jiā hěn dà, yě hěn wēnnuǎn.}\\
« Ma famille est grande et aussi très chaleureuse. »\\[4pt]
\zh{我家有六口人：爸爸、妈妈、两个姐姐、弟弟和我。}\\
\emph{Wǒ jiā yǒu liù kǒu rén : bàba, māma, liǎng ge jiějie, dìdi hé wǒ.}\\
« Il y a six personnes dans ma famille : papa, maman, mes deux grandes sœurs, mon petit frère et moi. »\\[4pt]
\zh{我们一家人相亲相爱，每天都很开心。}\\
\emph{Wǒmen yì jiā rén xiāng qīn xiāng ài, měitiān dōu hěn kāixīn.}\\
« Toute notre famille s'aime tendrement, et nous sommes heureux chaque jour. »\\[4pt]
\zh{在家里，我们互相帮助：姐姐帮我学习，我帮弟弟做作业。}\\
\emph{Zài jiā lǐ, wǒmen hùxiāng bāngzhù : jiějie bāng wǒ xuéxí, wǒ bāng dìdi zuò zuòyè.}\\
« À la maison, nous nous aidons mutuellement : ma grande sœur m'aide à étudier, et j'aide mon petit frère à faire ses devoirs. »\\[4pt]
\zh{困难的时候，我们互相关心，互相支持。}\\
\emph{Kùnnan de shíhou, wǒmen hùxiāng guānxīn, hùxiāng zhīchí.}\\
« Dans les moments difficiles, nous nous soucions les uns des autres et nous nous soutenons. »\\[4pt]
\zh{我爱我的家，家人是我的力量！}\\
\emph{Wǒ ài wǒ de jiā, jiārén shì wǒ de lìliang!}\\
« J'aime ma famille, ma famille est ma force ! »\\[6pt]
\textbf{Analyse pédagogique (pour l'enseignant) :}
\begin{itemize}
\item \textbf{Plan respecté} : 1. Présentation (phr. 1-2) $\to$ 2. Sentiments (phr. 3) $\to$ 3. Actions concrètes (phr. 4) $\to$ 4. Soutien moral (phr. 5) $\to$ 5. Conclusion (phr. 6).
\item \textbf{Lexique} : 8 mots du glossaire (\zh{家}, \zh{温暖}, \zh{口}, \zh{人}, \zh{爱} [dans \zh{相亲相爱}], \zh{帮助}, \zh{关心}, \zh{支持}, \zh{力量}, \zh{互相} $\times$ 3). Objectif dépassé.
\item \textbf{Collocations} : 3 utilisées (\zh{相亲相爱}, \zh{互相帮助}, \zh{互相关心}, \zh{互相支持}). Objectif dépassé.
\item \textbf{Grammaire} : \zh{六口人} (classif. ok) ; \zh{两个姐姐} (\zh{两} + classif., pas \zh{二}) ; \zh{互相} bien placé avant V ; \zh{帮 + qqn + V} correct ; Complément de temps \zh{困难的时候} en tête de phrase (ordre chinois standard).
\item \textbf{Ponctuation} : Usage correct de \zh{：} (deux points) pour introduire la liste d'actions dans la phrase 4.
\item \textbf{Conseil oral} : Marquer la pause après \zh{在家里} (sujet de phrase), après \zh{困难的时候}. Ton 4 net sur \zh{爱}, \zh{互}, \zh{助}, \zh{力}.
\end{itemize}
\end{exoresolu}

\begin{aretenir}[Ce que je retiens — Synthèse de la Leçon 1]
\begin{itemize}
\item \textbf{Vocabulaire clé} : Nommer sa famille (\zh{家人}, \zh{爸爸}...\zh{妹妹}), compter (\zh{口人}), décrire l'ambiance (\zh{温暖}, \zh{和睦}), exprimer l'amour (\zh{爱}, \zh{相亲相爱}) et la solidarité (\zh{互相帮助}, \zh{互相关心}, \zh{支持}).
\item \textbf{Structures maîtrisées} : \zh{我家有 + Nb + 口人} ; \zh{帮/帮助 + qqn + V} ; \zh{互相 + V} (position avant V) ; Complément de temps/lieu en début de phrase.
\item \textbf{Écriture} : Radicaux \zh{宀} (toit), \zh{心} (cœur), \zh{力} (force), \zh{氵} (eau), \zh{彳} (pas). Ordre des traits respecté.
\item \textbf{Culture} : \zh{家和万事兴} — L'harmonie familiale comme fondement de la réussite, valeur partagée Cameroun/Chine. Distinction des appellations (paternel/maternel, aîné/cadet) cruciale en chinois.
\item \textbf{Compétence communicative} : Je sais présenter ma famille, décrire nos interactions solidaires et exprimer mon attachement en chinois correct (HSK 2-3).
\end{itemize}
\end{aretenir}

\newpage

\coursec{Méthodes et exercices résolus}

\coursec{Leçon 1 — Vocabulaire et compréhension de texte : amour de la famille et la solidarité familiale}

\courssub{Mise en route et découverte du thème}

\begin{experience}[Brainstorming et entrée en matière]
\begin{enumerate}
\item \textbf{Observation d'image.} Projeter une photo montrant une famille camerounaise nombreuse autour d'un plat de ndolé et une famille chinoise autour d'un repas du Nouvel An (饺子 jiǎozi). Demander : « Que voyez-vous ? Quels points communs ? Quelles différences ? »
\item \textbf{Mots-clés au tableau.} En colonnes : \zh{家} (famille/foyer), \zh{爱} (amour), \zh{团结} (solidarité/unité), \zh{孝顺} (piété filiale/respect des aînés). Faire deviner le sens par les composants (clés) : \zh{宀} (toit) + \zh{豕} (cochon) = la richesse sous le toit ; \zh{爫} (griffe/main) + \zh{心} (cœur) = la main qui porte le cœur.
\item \textbf{Question déclencheur.} « Dans votre culture, comment exprime-t-on l'amour familial ? Par des mots, des actes, des cadeaux ? » Noter les réponses en français, introduire le vocabulaire chinois au fur et à mesure.
\end{enumerate}
\end{experience}

\begin{savaistu}[Famille au Cameroun et en Chine : regards croisés]
\textbf{Cameroun :} Famille élargie (oncles, tantes, cousins = « frères/sœurs »), solidarité clanique forte, respect des aînés absolu, éducation collective (« il faut tout un village pour élever un enfant »). Langues nationales : termes de parenté variables (bassa, bamileke, fulfulde, ewondo...).

\textbf{Chine :} Vocabulaire de parenté extrêmement précis (séparation lignée paternelle/maternelle, ordre de naissance). Ex. : \zh{爷爷} (père du père), \zh{外公} (père de la mère) ; \zh{哥哥} (frère aîné), \zh{弟弟} (frère cadet). Politique de l'enfant unique (1979–2015) → familles « 4-2-1 » (4 grands-parents, 2 parents, 1 enfant) ; depuis 2021 : politique des trois enfants. Valeur centrale : \zh{孝顺} (xiàoshùn, piété filiale) = respect, obéissance, soin matériel et moral aux parents âgés.

\textbf{Point commun :} La famille comme cellule de base de la société, refuge affectif et filet de sécurité sociale. L'entraide (\zh{互助} hùzhù) y est une norme, pas une option.
\end{savaistu}

\courssub{Texte support : lecture et traduction}

\begin{textebox}[Texte original : « Ma famille unie »]
\zh{我家在雅温得，有五口人：爸爸、妈妈、哥哥、妹妹和我。} \\
\zh{我们都很爱我们的家。每天妈妈做饭，我帮助她。} \\
\zh{周末，我们全家一起吃饭、看电视，哥哥帮助妹妹做作业。} \\
\zh{我觉得我们的家很温暖，很幸福。}

\emph{Wǒ jiā zài Yǎwēndé, yǒu wǔ kǒu rén : bàba, māma, gēge, mèimei hé wǒ. Wǒmen dōu hěn ài wǒmen de jiā. Měitiān māma zuò fàn, wǒ bāngzhù tā. Zhōumò, wǒmen quán jiā yìqǐ chī fàn, kàn diànshì, gēge bāngzhù mèimei zuò zuòyè. Wǒ juéde wǒmen de jiā hěn wēnnuǎn, hěn xìngfú.}

\textbf{Traduction :} « Ma famille est à Yaoundé, nous sommes cinq : papa, maman, frère aîné, sœur cadette et moi. Nous aimons tous notre famille. Chaque jour maman cuisine, je l'aide. Le week-end, toute la famille mange ensemble et regarde la télé, le frère aîné aide la sœur cadette à faire ses devoirs. Je trouve que notre foyer est très chaleureux, très heureux. »
\end{textebox}

\begin{propriete}[Structure du texte : repères syntaxiques]
\begin{itemize}
\item \textbf{Phrase 1 (Présentation) :} Sujet \zh{我家} + Verbe \zh{有} + Nombre + Classificateur \zh{口} + Nom \zh{人} + \zh{：} (énumération des membres).
\item \textbf{Phrase 2 (Sentiment) :} Sujet \zh{我们} + Adverbe \zh{都} + Adverbe de degré \zh{很} + Verbe \zh{爱} + Complément \zh{我们的家} (\zh{的} = possession/attribut).
\item \textbf{Phrase 3 (Routine) :} Temps \zh{每天} + Sujet \zh{妈妈} + Verbe \zh{做饭} ; Sujet \zh{我} + Verbe \zh{帮助} + Objet \zh{她}.
\item \textbf{Phrase 4 (Activité collective) :} Temps \zh{周末} + Sujet \zh{我们全家} + Adverbe \zh{一起} + Verbes \zh{吃饭}、\zh{看电视} (série) ; Sujet \zh{哥哥} + Verbe \zh{帮助} + Objet \zh{妹妹} + Action \zh{做作业}.
\item \textbf{Phrase 5 (Ressenti) :} Sujet \zh{我} + Verbe \zh{觉得} + Complément \zh{很温暖，很幸福} (adjectifs prédicatifs, sans \zh{是}).
\end{itemize}
\end{propriete}

\courssub{Glossaire du thème : vocabulaire de la famille et de la solidarité}

\begin{center}
\begin{tabularx}{\linewidth}{|X|X|X|X|}
\hline
\rowcolor{popblueL} \textbf{Caractères} & \textbf{Pinyin} & \textbf{Nature} & \textbf{Traduction} \\
\hline
家 & jiā & n. & famille, foyer, maison ; suffixe de spécialiste (ex. \zh{画家} peintre) \\
\hline
口 & kǒu & m. & classificateur pour les membres de la famille (bouche) \\
\hline
人 & rén & n. & personne, être humain \\
\hline
爸爸 & bàba & n. & papa, père \\
\hline
妈妈 & māma & n. & maman, mère \\
\hline
哥哥 & gēge & n. & frère aîné \\
\hline
姐姐 & jiějie & n. & sœur aînée \\
\hline
弟弟 & dìdi & n. & frère cadet \\
\hline
妹妹 & mèimei & n. & sœur cadette \\
\hline
爷爷 & yéye & n. & grand-père paternel \\
\hline
奶奶 & nǎinai & n. & grand-mère paternelle \\
\hline
外公 & wàigōng & n. & grand-père maternel \\
\hline
外婆 & wàipó & n. & grand-mère maternelle \\
\hline
爱 & ài & v. & aimer, affection \\
\hline
帮助 & bāngzhù & v. & aider, assistance \\
\hline
照顾 & zhàogù & v. & prendre soin de, s'occuper de \\
\hline
一起 & yìqǐ & adv. & ensemble \\
\hline
全家 & quán jiā & n. & toute la famille \\
\hline
每天 & měitiān & adv. & chaque jour, tous les jours \\
\hline
周末 & zhōumò & n./adv. & week-end, le week-end \\
\hline
做饭 & zuò fàn & v.o. & faire la cuisine, préparer le repas \\
\hline
吃饭 & chī fàn & v.o. & manger (un repas) \\
\hline
看电视 & kàn diànshì & v.o. & regarder la télévision \\
\hline
做作业 & zuò zuòyè & v.o. & faire ses devoirs \\
\hline
幸福 & xìngfú & adj. & heureux, bonheur \\
\hline
温暖 & wēnnuǎn & adj. & chaleureux, tiède \\
\hline
和 & hé & conj. & et (liaison noms/pronoms) \\
\hline
都 & dōu & adv. & tous, toutes (placé avant le verbe) \\
\hline
很 & hěn & adv. & très, beaucoup (devant adj. prédicatif) \\
\hline
的 & de & part. & possession (mon/ma/mes), attribut \\
\hline
觉得 & juéde & v. & penser que, avoir l'impression que, trouver que \\
\hline
比 & bǐ & prep. & (comparatif) plus que \\
\hline
大 & dà & adj. & grand, âgé (pour l'âge) \\
\hline
小 & xiǎo & adj. & petit, jeune \\
\hline
岁 & suì & m. & an (âge), classificateur d'âge \\
\hline
两 & liǎng & num. & deux (devant classificateur) \\
\hline
雅温得 & Yǎwēndé & n.pr. & Yaoundé (capitale du Cameroun) \\
\hline
\end{tabularx}
\end{center}

\courssub{Clés (radicaux) et ordre des traits}

\begin{definition}[Analyse de caractères clés du thème]
\textbf{1. 家 (jiā) — Famille/Foyer} \\
\cle{Radical (clé) :} \zh{宀} (mián, « toit »). \cle{Partie phonétique :} \zh{豕} (shǐ, « cochon »). \cle{Sens :} Sous le toit, il y a un cochon = richesse, foyer prospère. \\
\cle{Ordre des traits (10) :} 1. Point haut gauche (丶) 2. Point haut droit (丶) 3. Horizontal court (一) 4. Horizontal long (一) — forme le toit \zh{宀}. 5. Vertical (丨) 6. Horizontal (一) 7. Crochet vertical (㇆) 8. Horizontal (一) 9. Horizontal (一) 10. Horizontal (一) — forme \zh{豕}.

\textbf{2. 爱 (ài) — Aimer (simplifié)} \\
\cle{Radical :} \zh{爫} (zhuǎ, « griffe/main qui agit ») en haut. \cle{Cœur :} \zh{心} (xīn) au centre (dans la version traditionnelle \zh{愛}, le cœur est au milieu ; en simplifié \zh{爱}, le cœur est remplacé par \zh{友} « ami », mais le radical officiel reste souvent \zh{爫} ou \zh{心} selon les dictionnaires). \cle{Ordre (10 traits simplifié) :} Griffe (5 traits) → \zh{友} (7 traits : haut \zh{又} + bas \zh{一}...). \emph{Astuce :} « La main (\zh{爫}) qui tient un ami (\zh{友}) dans le cœur. »

\textbf{3. 爸 (bà) / 妈 (mā) — Père / Mère} \\
\cle{Phono-sémantiques :} Clé sémantique à gauche, phonétique à droite. \\
\zh{爸} : \zh{父} (fù, « père », clé \zh{父}) + \zh{巴} (bā, son). \zh{妈} : \zh{女} (nǚ, « femme », clé \zh{女}) + \zh{马} (mǎ, son). \\
\cle{Ordre \zh{爸} (8) :} \zh{父} (4 traits : 丿 一 丿 乀) + \zh{巴} (4 traits : 一 丨 一 丶). \\
\cle{Ordre \zh{妈} (6) :} \zh{女} (3 traits : 丿 ㇋ 一) + \zh{马} (3 traits : 丨 ㇆ 丶).

\textbf{4. 哥 (gē) / 姐 (jiě) / 弟 (dì) / 妹 (mèi) — Frères/Sœurs} \\
\cle{Clé commune :} \zh{可} (gē), \zh{女} (jiě, mèi), \zh{弓} (dì). \emph{Note :} \zh{哥} clé \zh{口} (bouche/parole), \zh{姐/妹} clé \zh{女} (femme), \zh{弟} clé \zh{弓} (arc → jeune homme). \\
\cle{Ordre \zh{哥} (10) :} \zh{可} (5) + \zh{口} (3) + trait final. \cle{Ordre \zh{姐} (8) :} \zh{女} (3) + \zh{且} (5). \cle{Ordre \zh{弟} (7) :} \zh{弓} (3) + \zh{弟} bas (4). \cle{Ordre \zh{妹} (8) :} \zh{女} (3) + \zh{未} (5).

\textbf{5. 口 (kǒu) — Bouche / Classificateur} \\
\cle{Radical :} \zh{口} (lui-même). \cle{3 traits :} 1. Vertical gauche (丨) 2. Horizontal + crochet bas (㇆) 3. Horizontal de fermeture (一). \emph{Ne pas faire un carré en 4 traits !}

\textbf{6. 幸 (xìng) / 福 (fú) — Bonheur} \\
\zh{幸} : clé \zh{干} (gān). 8 traits. \zh{福} : clé \zh{示} (shì, « autel/rite »). 13 traits. \zh{福} = \zh{示} + \zh{畐} (tonneau plein) = bénédiction divine/abondance.
\end{definition}

\begin{experience}[Atelier écriture : « Mon arbre généalogique »]
Chaque élève dessine son arbre généalogique simplifié (3 générations) sur une feuille A4. Il écrit les caractères chinois pour chaque membre (ex. \zh{爷爷}, \zh{爸爸}, \zh{我}, \zh{弟弟}) en respectant l'ordre des traits. En binôme, ils se présentent mutuellement : \zh{这是我的爸爸，他叫……} (Zhè shì wǒ de bàba, tā jiào…). L'enseignant circule et corrige l'ordre des traits (vertical avant horizontal si croisé, haut → bas, gauche → droite).
\end{experience}

\courssub{Compréhension du texte}

\begin{exoresolu}[Lecture guidée du texte support]
\textbf{Énoncé.} Lis le texte de la \textbf{textebox} ci-dessus. Réponds en français aux questions :
1. Combien de personnes composent la famille ? Cite le classificateur utilisé.
2. Qui fait la cuisine chaque jour ? Qui l'aide ?
3. Que fait la famille le week-end ? (Deux activités).
4. Qui aide qui pour les devoirs ?
5. Quels deux adjectifs décrivent le foyer à la fin ?

\tcblower
\textbf{Solution détaillée.}
\textbf{1.} Cinq personnes (\zh{五口人}). Classificateur : \zh{口} (kǒu).
\textbf{2.} Maman (\zh{妈妈}) fait la cuisine (\zh{做饭}) ; « je » (\zh{我}) l'aide (\zh{帮助她}).
\textbf{3.} Manger ensemble (\zh{一起吃饭}) et regarder la télévision ensemble (\zh{一起看电视}).
\textbf{4.} Le frère aîné (\zh{哥哥}) aide la sœur cadette (\zh{妹妹}) à faire ses devoirs (\zh{做作业}).
\textbf{5.} \zh{温暖} (wēnnuǎn, chaleureux) et \zh{幸福} (xìngfú, heureux).
\end{exoresolu}

\begin{exercice}[Compréhension : nouveau texte court]
\textbf{Énoncé.} Lis le texte suivant et réponds en chinois par des phrases complètes.

\zh{大卫家在杜阿拉。他家有六口人：爸爸、妈妈、两个姐姐、一个哥哥和他。爸爸是工程师，妈妈在医院工作。大卫今年十六岁，上高中。他的姐姐们都在大学读书。每天晚上，全家人一起吃饭，聊天。大卫觉得很幸福。}

\emph{Dàwèi jiā zài Dùālā. Tā jiā yǒu liù kǒu rén : bàba, māma, liǎng ge jiějie, yí ge gēge hé tā. Bàba shì gōngchéngshī, māma zài yīyuàn gōngzuò. Dàwèi jīnnián shíliù suì, shàng gāozhōng. Tā de jiějie men dōu zài dàxué dúshū. Měitiān wǎnshàng, quán jiā rén yìqǐ chī fàn, liáo tiān. Dàwèi juéde hěn xìngfú.}

\textbf{Vocabulaire nouveau :} \zh{杜阿拉} (Dùālā, Douala), \zh{工程师} (gōngchéngshī, ingénieur), \zh{医院} (yīyuàn, hôpital), \zh{工作} (gōngzuò, travailler), \zh{十六} (shíliù, 16), \zh{高中} (gāozhōng, lycée), \zh{大学} (dàxué, université), \zh{读书} (dúshū, étudier/lire), \zh{晚上} (wǎnshàng, soir), \zh{聊天} (liáo tiān, discuter).

\textbf{Questions :}
1. 大卫家在哪儿？ 2. 他家有几口人？ 3. 爸爸是做什么的？ 4. 姐姐们在哪儿读书？ 5. 每天晚上他们做什么？ 6. 大卫觉得怎么样？
\end{exercice}

\courssub{Collocations et production guidée}

\begin{propriete}[Collocations indispensables du thème (HSK 1-3)]
\begin{center}
\begin{tabularx}{\linewidth}{|X|X|X|}
\hline
\rowcolor{popblueL} \textbf{Structure / Collocation} & \textbf{Exemple (Caractères + Pinyin)} & \textbf{Traduction} \\
\hline
\zh{有} + Nombre + \zh{口} + \zh{人} & \zh{我家有四口人。} (Wǒ jiā yǒu sì kǒu rén.) & Nous sommes quatre dans ma famille. \\
\hline
\zh{爱} + \zh{的} + Nom & \zh{我爱我的家。} (Wǒ ài wǒ de jiā.) & J'aime ma famille. \\
\hline
\zh{帮助} + Personne + Verbe & \zh{我帮助妈妈做饭。} (Wǒ bāngzhù māma zuò fàn.) & J'aide maman à cuisiner. \\
\hline
\zh{照顾} + Personne & \zh{哥哥照顾妹妹。} (Gēge zhàogù mèimei.) & Le grand frère s'occupe de la petite sœur. \\
\hline
Sujet + \zh{一起} + Verbe & \zh{我们一起吃饭。} (Wǒmen yìqǐ chī fàn.) & Nous mangeons ensemble. \\
\hline
Sujet + Temps + Verbe & \zh{每天我做作业。} (Měitiān wǒ zuò zuòyè.) & Chaque jour je fais mes devoirs. \\
\hline
Sujet + \zh{很} + Adjectif & \zh{我很幸福。} (Wǒ hěn xìngfú.) & Je suis très heureux. \\
\hline
Sujet + \zh{觉得} + Adjectif/Clause & \zh{我觉得很温暖。} (Wǒ juéde hěn wēnnuǎn.) & Je trouve cela très chaleureux. \\
\hline
A + \zh{比} + B + Adj (+ Quantité) & \zh{哥哥比我大两岁。} (Gēge bǐ wǒ dà liǎng suì.) & Mon frère a 2 ans de plus que moi. \\
\hline
\end{tabularx}
\end{center}
\end{propriete}

\begin{methode}[Rédiger un paragraphe « Ma famille » (5-6 phrases) — Modèle d'examen]
\begin{enumerate}
\item \textbf{Phrase 1 — Présentation :} \zh{我家在[Ville], 有[Nombre]口人：[Liste membres].} (Utiliser \zh{两} devant \zh{个/口} pour « deux »).
\item \textbf{Phrase 2 — Profession/Activité :} \zh{爸爸是[Métier]，妈妈是[Métier].} (Pas d'article devant le métier).
\item \textbf{Phrase 3 — Solidarité quotidienne :} \zh{每天[Qui]帮助[Qui][Action].} (Ex. \zh{每天我帮助妈妈做饭。})
\item \textbf{Phrase 4 — Moment collectif :} \zh{周末/晚上，我们全家一起[Activité 1]、[Activité 2].}
\item \textbf{Phrase 5 — Sentiment final :} \zh{我觉得我们的家很[Adjectif 1]，很[Adjectif 2].} (\zh{幸福}、\zh{温暖}、\zh{和睦} hé mù « harmonieux »).
\item \textbf{Relecture :} Vérifier l'ordre Sujet-Temps-Verbe, les classificateurs (\zh{口}/\zh{个}), \zh{两} vs \zh{二}, l'absence de \zh{是} devant adjectif.
\end{enumerate}
\end{methode}

\begin{exoresolu}[Production guidée : mon paragraphe]
\textbf{Énoncé.} Écris un paragraphe de 5 phrases pour présenter ta famille et la solidarité qui y règne. Mots imposés : \zh{有}、\zh{爱}、\zh{帮助}、\zh{一起}、\zh{幸福}.

\tcblower
\textbf{Modèle de rédaction (niveau HSK 2-3) :}

\zh{我家在雅温得，有五口人：爸爸、妈妈、哥哥、妹妹和我。} \\
\emph{Wǒ jiā zài Yǎwēndé, yǒu wǔ kǒu rén : bàba, māma, gēge, mèimei hé wǒ.}

\zh{爸爸是老师，妈妈是护士，我们都很爱我们的家。} \\
\emph{Bàba shì lǎoshī, māma shì hùshì, wǒmen dōu hěn ài wǒmen de jiā.}

\zh{每天我帮助妈妈做饭，哥哥帮助妹妹做作业。} \\
\emph{Měitiān wǒ bāngzhù māma zuò fàn, gēge bāngzhù mèimei zuò zuòyè.}

\zh{周末我们全家一起吃饭、看电视、聊天。} \\
\emph{Zhōumò wǒmen quán jiā yìqǐ chī fàn, kàn diànshì, liáo tiān.}

\zh{我觉得我们的家很温暖，很幸福。} \\
\emph{Wǒ juéde wǒmen de jiā hěn wēnnuǎn, hěn xìngfú.}

\textbf{Analyse grammaticale :} Phrase 1 : structure d'existence \zh{有} + \zh{口人}. Phrase 2 : métier sans article, \zh{都很爱} (adverbe \zh{都} + \zh{很} + verbe). Phrase 3 : deux clauses coordonnées, temps \zh{每天} en tête, structure \zh{帮助} + O + V. Phrase 4 : \zh{全家一起} + série de VO séparés par \zh{、}. Phrase 5 : \zh{觉得} + double adjectif prédicatif.
\end{exoresolu}

\begin{exercice}[Production libre : ma famille camerounaise]
\textbf{Consigne.} En t'inspirant du modèle ci-dessus, écris ton propre paragraphe (5-6 phrases) en adaptant la réalité de ta famille (nombre, métiers, activités du week-end au Cameroun : marché, église, football, repas de fête...). Utilise au moins 8 mots du glossaire. Écris en caractères, pinyin et français.
\end{exercice}

\courssub{Points de vigilance et erreurs types (Bac Blanc)}

\begin{attention}[Les 7 erreurs qui coûtent des points au Bac]
\begin{enumerate}
\item \textbf{\zh{两} (liǎng) vs \zh{二} (èr).} \cle{Règle :} Devant un classificateur (\zh{个}、\zh{口}、\zh{岁}...), on utilise \zh{两} : \zh{两个姐姐} (liǎng ge jiějie), \zh{两口人} (liǎng kǒu rén). \zh{二} s'emploie dans les nombres purs (\zh{十二} shíèr), les dates (\zh{二月} èryuè), les étages (\zh{二楼} èrlóu). \textbf{Piège :} \zh{二岁} se dit parfois à l'oral rapide, mais \zh{两岁} est la norme écrite standard.
\item \textbf{Classificateur \zh{口} (kǒu) obligatoire pour \zh{人} (membres famille).} \zh{我家有五口人} (Correct) vs \zh{我家有五个} (Incorrect pour le nombre total de la famille). \zh{个} s'utilise pour compter des individus isolés : \zh{三个哥哥} (trois frères aînés).
\item \textbf{Position du complément de temps (AdvT).} \cle{Ordre immuable :} Sujet + \textbf{Temps} + Verbe + Complément. \\ \textbf{Faux :} \zh{我帮助妈妈每天。} \textbf{Juste :} \zh{我每天帮助妈妈。} / \zh{每天我帮助妈妈。} (Temps en tête = insistance).
\item \textbf{Adjectif prédicatif sans \zh{是} (shì).} \textbf{Faux :} \zh{我是很幸福。} \zh{是} ne lie pas sujet et adjectif. \textbf{Juste :} \zh{我很幸福。} (\zh{很} = lien copulatif + « très » affaibli). \textbf{Négation :} \zh{我不幸福。} (Pas \zh{不很}).
\item \textbf{Confusion caractères / tons : paires minimales.} \zh{妈} (mā, mère, clé \zh{女}) ≠ \zh{马} (mǎ, cheval) ≠ \zh{骂} (mà, gronder) ≠ \zh{吗} (ma, particule interrogative). \zh{爱} (ài, aimer, clé \zh{爫}/\zh{心}) ≠ \zh{受} (shòu, recevoir, clé \zh{又}). \zh{好} (hǎo, bon) ≠ \zh{号} (hào, numéro). \textbf{Action :} Apprendre le caractère \textbf{avec} son radical et son ton.
\item \textbf{\zh{和} (hé) ne lie que des noms/pronoms, pas des phrases.} \textbf{Faux :} \zh{我喜欢吃饭和我看电视。} \textbf{Juste :} \zh{我喜欢吃饭，也喜欢看电视。} ou \zh{我喜欢吃饭和看电视。} (liaison de deux VO).
\item \textbf{Oubli de \zh{的} (de) pour la possession/attribut.} \zh{我妈妈} (oral, familier) → \zh{我的妈妈} (écrit, standard). \zh{我们家} (oral) → \zh{我们的家} (écrit). À l'examen écrit, \textbf{garde \zh{的}}.
\end{enumerate}
\end{attention}

\begin{corrige}[Exercice : Compréhension (Douala)]
\textbf{1.} \zh{大卫家在杜阿拉。} (Dàwèi jiā zài Dùālā.) — Sa famille est à Douala.
\textbf{2.} \zh{他家有六口人。} (Tā jiā yǒu liù kǒu rén.) — Sa famille compte six personnes.
\textbf{3.} \zh{爸爸是工程师。} (Bàba shì gōngchéngshī.) — Papa est ingénieur.
\textbf{4.} \zh{姐姐们在大学读书。} (Jiějie men zài dàxué dúshū.) — Les sœurs aînées étudient à l'université.
\textbf{5.} \zh{每天晚上他们一起吃饭、聊天。} (Měitiān wǎnshàng tāmen yìqǐ chī fàn, liáo tiān.) — Chaque soir ils mangent et discutent ensemble.
\textbf{6.} \zh{大卫觉得很幸福。} (Dàwèi juéde hěn xìngfú.) — David se sent très heureux.
\end{corrige}

\begin{aretenir}[Synthèse de la leçon 1 — Check-list compétences]
\begin{itemize}
\item \textbf{Vocabulaire :} 30+ mots (membres famille, verbes vie quotidienne, adjectifs sentiment, mots-outils \zh{的、和、都、很、一起、比、两}).
\item \textbf{Grammaire :} Structure \zh{有} + Nb + \zh{口人} ; Ordre S + Temps + V ; Adj prédicatif avec \zh{很} (pas \zh{是}) ; Comparatif \zh{比} ; \zh{帮助/照顾} + O + V ; \zh{一起} + V ; \zh{觉得} + Adj.
\item \textbf{Écriture :} 10 caractères clés analysés (radical, ordre traits) : \zh{家、爱、爸、妈、哥、姐、弟、妹、口、福}.
\item \textbf{Compréhension :} Répondre par phrase complète en reprenant la structure de la question (\zh{几}→Nb, \zh{谁}→Pers, \zh{什么时候}→Temps, \zh{为什么}→Raison \zh{因为}...\zh{所以}...).
\item \textbf{Expression écrite :} Produire un paragraphe cohérent de 5 phrases (Présentation → Métiers → Entraide quotidienne → Moment collectif → Sentiment).
\item \textbf{Socioculturel :} Précision du vocabulaire de parenté chinois (âge, lignée) vs terme générique « frère/sœur » français ; valeur partagée de solidarité intergénérationnelle Cameroun/Chine.
\end{itemize}
\end{aretenir}

\newpage

\coursec{Exercices}

\courssub{Je vérifie mes connaissances}

\begin{exercice}[QCM : le bon sens]
Pour chaque phrase, choisis la réponse correcte (A, B ou C).

\begin{enumerate}
\item \zh{我家有四口人。} signifie :
\begin{enumerate}
\item[A.] Ma maison a quatre portes.
\item[B.] Ma famille compte quatre personnes.
\item[C.] J'ai quatre maisons.
\end{enumerate}
\item \zh{我们互相帮助。} signifie :
\begin{enumerate}
\item[A.] Nous nous aidons mutuellement.
\item[B.] Nous nous disputons souvent.
\item[C.] Nous nous ignorons.
\end{enumerate}
\item Dans \zh{家人}, le caractère \zh{人} signifie :
\begin{enumerate}
\item[A.] maison
\item[B.] cœur
\item[C.] personne
\end{enumerate}
\item \zh{家是最温暖的地方。} signifie :
\begin{enumerate}
\item[A.] La maison est l'endroit le plus froid.
\item[B.] La famille est l'endroit le plus chaleureux.
\item[C.] La famille est le plus grand endroit.
\end{enumerate}
\end{enumerate}
\end{exercice}

\begin{exercice}[Vrai ou faux justifié]
D'après le texte support de la leçon, dis si chaque affirmation est vraie (对) ou fausse (不对), puis justifie ta réponse en français par une citation du texte en chinois.

\begin{enumerate}
\item \zh{全家人周末一起吃饭。} (Quánjiārén zhōumò yìqǐ chīfàn.)
\item \zh{我不爱我的家。} (Wǒ bú ài wǒ de jiā.)
\item \zh{弟弟帮助我学习汉语。} (Dìdi bāngzhù wǒ xuéxí Hànyǔ.)
\item \zh{我们家有五口人。} (Wǒmen jiā yǒu wǔ kǒu rén.)
\end{enumerate}
\end{exercice}

\begin{exercice}[Vocabulaire : le mot qui convient]
Complète chaque phrase avec le mot qui convient, choisi entre parenthèses.

\begin{enumerate}
\item \zh{我们互相\_\_\_。} (帮助 / 关心)
\item \zh{我觉得家是最\_\_\_的地方。} (温暖 / 冷)
\item \zh{爸爸每天\_\_\_我们。} (爱 / 走)
\item \zh{一家人要\_\_\_在一起。} (团结 / 分开)
\end{enumerate}
\end{exercice}

\begin{exercice}[Pinyin et caractères]
Associe chaque caractère ou mot à son pinyin, puis donne sa traduction française.

\begin{center}
\begin{tabularx}{\linewidth}{|X|X|X|}\hline
\rowcolor{popblueL} Caractères & Pinyin à associer & Traduction \\ \hline
爱 & a) jiā & ........................ \\ \hline
家 & b) ài & ........................ \\ \hline
帮助 & c) tuánjié & ........................ \\ \hline
团结 & d) bāngzhù & ........................ \\ \hline
温暖 & e) wēnnuǎn & ........................ \\ \hline
\end{tabularx}
\end{center}
\end{exercice}

\courssub{Je m'entraîne}

\begin{exercice}[Traduction français-chinois]
Traduis en chinois simplifié (avec le pinyin) les phrases suivantes :

\begin{enumerate}
\item Ma famille compte quatre personnes.
\item Nous nous entraidons chaque jour.
\item J'aime beaucoup ma famille.
\item Le week-end, toute la famille mange ensemble.
\end{enumerate}
\end{exercice}

\begin{exercice}[Texte à trous : mots-outils et vocabulaire]
Complète le texte suivant avec les mots de la liste : \zh{的}、\zh{口}、\zh{互相}、\zh{最}、\zh{爱}。

\begin{textebox}[我家]
我\_\_\_家有四\_\_\_人：爸爸、妈妈、弟弟和我。\\
Wǒ \_\_\_ jiā yǒu sì \_\_\_ rén : bàba, māma, dìdi hé wǒ.\\
我们\_\_\_帮助，我\_\_\_我的家。\\
Wǒmen \_\_\_ bāngzhù, wǒ \_\_\_ wǒ de jiā.\\
家是我\_\_\_喜欢的地方。\\
Jiā shì wǒ \_\_\_ xǐhuan de dìfang.
\end{textebox}
\end{exercice}

\begin{exercice}[Remise en ordre]
Remets les mots dans l'ordre pour former une phrase correcte, puis traduis-la en français.

\begin{enumerate}
\item \zh{人 / 我家 / 五 / 有 / 口}
\item \zh{帮助 / 互相 / 我们 / 每天}
\item \zh{最 / 家 / 的 / 地方 / 是 / 温暖}
\item \zh{一起 / 周末 / 吃饭 / 全家人}
\end{enumerate}
\end{exercice}

\begin{exercice}[Caractères, clés et traits]
\begin{enumerate}
\item Associe chaque caractère à sa clé (radical) et à sa traduction : \zh{爱}、\zh{家}、\zh{帮}、\zh{团}、\zh{暖}. Radicaux proposés : \zh{宀} (toit), \zh{爫} (griffe), \zh{日} (soleil), \zh{囗} (enceinte), \zh{巾} (tissu).
\item Écris les caractères \zh{爱} et \zh{家} en respectant l'ordre des traits (haut $\rightarrow$ bas, gauche $\rightarrow$ droite). Le respect de l'ordre des traits sera noté.
\end{enumerate}
\end{exercice}

\courssub{Je me prépare au Bac}

\begin{exercice}[Type Bac : compréhension de texte]
Lis le texte suivant, puis réponds aux questions en chinois, par des phrases complètes.

\begin{textebox}[我爱我的家]
我爱我的家。我家有四口人：爸爸、妈妈、妹妹和我。\\
Wǒ ài wǒ de jiā. Wǒ jiā yǒu sì kǒu rén : bàba, māma, mèimei hé wǒ.\\
爸爸是老师，妈妈是医生。他们每天工作很忙，但是晚上都回家吃饭。\\
Bàba shì lǎoshī, māma shì yīshēng. Tāmen měitiān gōngzuò hěn máng, dànshì wǎnshang dōu huí jiā chīfàn.\\
周末，全家人一起吃饭、看电视。妹妹帮助我学习汉语，我帮助她学习法语。\\
Zhōumò, quánjiārén yìqǐ chīfàn, kàn diànshì. Mèimei bāngzhù wǒ xuéxí Hànyǔ, wǒ bāngzhù tā xuéxí Fǎyǔ.\\
我们互相关心，互相帮助。我觉得家是最温暖的地方。\\
Wǒmen hùxiāng guānxīn, hùxiāng bāngzhù. Wǒ juéde jiā shì zuì wēnnuǎn de dìfang.
\end{textebox}

Questions de compréhension (réponds en chinois) :
\begin{enumerate}
\item \zh{他家有几口人？都是谁？} (Qui ? Combien de personnes ?)
\item \zh{爸爸和妈妈做什么工作？} (Quoi ? Quels métiers ?)
\item \zh{全家人什么时候一起吃饭、看电视？} (Quand ?)
\item \zh{为什么作者觉得家是最温暖的地方？} (Pourquoi ?)
\item \zh{你喜欢你的家吗？为什么？} (Ton opinion personnelle.)
\end{enumerate}

Questions de langue :
\begin{enumerate}
\item Relève dans le texte deux mots construits avec la clé \zh{宀} (toit).
\item Explique en français le rôle de \zh{互相} dans la dernière phrase.
\end{enumerate}
\end{exercice}

\begin{exercice}[Type Bac : thème et version]
\textbf{Thème.} Traduis en chinois simplifié (avec le pinyin) :

\begin{enumerate}
\item Ma famille compte cinq personnes : mon père, ma mère, mes deux frères et moi.
\item Nous nous entraidons chaque jour et nous nous aimons beaucoup.
\item Le dimanche, toute la famille mange ensemble à la maison.
\end{enumerate}

\textbf{Version.} Traduis en français :

\begin{textebox}[Version]
我的家在雅温得。我们家人很多，但是大家很团结。\\
Wǒ de jiā zài Yǎwēndé. Wǒmen jiā rén hěn duō, dànshì dàjiā hěn tuánjié.\\
有困难的时候，我们互相帮助。我爱我的家人。\\
Yǒu kùnnan de shíhou, wǒmen hùxiāng bāngzhù. Wǒ ài wǒ de jiārén.
\end{textebox}
\end{exercice}

\begin{exercice}[Type Bac : expression écrite]
\textbf{Sujet :} \zh{我的家} (Wǒ de jiā) — « Ma famille ».

Rédige un court texte en chinois simplifié (entre 60 et 80 caractères) dans lequel tu :
\begin{itemize}
\item présentes les membres de ta famille (combien de personnes, qui) ;
\item décris une activité que vous faites ensemble (repas, week-end, fête) ;
\item exprimes ton amour pour ta famille et l'importance de la solidarité familiale.
\end{itemize}

\textbf{Contraintes :} utilise au moins une fois les mots \zh{互相}、\zh{帮助}、\zh{温暖} et la structure \zh{我家有……口人}. Écris le pinyin sous chaque phrase. Soigne l'ordre des traits des caractères appris dans la leçon (\zh{爱}、\zh{家}).
\end{exercice}

\newpage

\coursec{Corrigés des exercices}

\coursec{Leçon 1 — Vocabulaire et compréhension de texte : Amour de la famille et solidarité familiale}

\courssub{Mise en route et découverte du thème}

\begin{textebox}[Situation de communication : Une famille à Yaoundé]
Regardez l'image et lisez le dialogue court. Essayez de deviner le sens des mots en gras.
\begin{center}
\begin{tabularx}{\linewidth}{|X|X|}\hline
\zh{图片：一家四口围坐在餐桌旁，父亲夹菜给母亲，妹妹帮哥哥写作业。} & \zh{— 爸爸，妈妈，妹妹，我。我们四口人。}\\\zh{— 周末一起吃饭、看电视。}\\\shortstack[l]{\zh{— 妹妹帮助我学习汉语。}\\ \zh{— 我们互相关心，互相帮助。}\\ \zh{— 家，真温暖！}} \\ \hline
\end{tabularx}
\end{center}
\textbf{Question orale :} Combien de personnes y a-t-il dans cette famille ? Quelles activités font-elles ensemble le week-end ?
\end{textebox}

\begin{experience}[Brainstorming lexical : « La famille »]
En binôme, listez en 2 minutes tous les mots chinois (caractères + pinyin) que vous connaissez sur le thème de la famille (membres, verbes d'action, adjectifs de sentiment). Mettez en commun au tableau. Le professeur classe les mots par nature (noms, verbes, adjectifs, classificateurs).
\end{experience}

\courssub{Texte support : Lecture et traduction}

\begin{textebox}[Texte original : 我的家]
\zh{我家有四口人：爸爸、妈妈、妹妹和我。} \\
\zh{Wǒ jiā yǒu sì kǒu rén : bàba, māma, mèimei hé wǒ.} \\
« Ma famille compte quatre personnes : papa, maman, ma petite sœur et moi. » \\[0.5em]

\zh{爸爸是老师，妈妈是医生，妹妹是学生。} \\
\zh{Bàba shì lǎoshī, māma shì yīshēng, mèimei shì xuésheng.} \\
« Papa est professeur, maman est médecin, ma petite sœur est élève. » \\[0.5em]

\zh{周末，全家人一起吃饭、看电视，家里很安静，也很温暖。} \\
\zh{Zhōumò, quánjiārén yìqǐ chīfàn, kàn diànshì, jiāli hěn ānjìng, yě hěn wēnnuǎn.} \\
« Le week-end, toute la famille mange ensemble et regarde la télévision ; à la maison c'est très calme et aussi très chaleureux. » \\[0.5em]

\zh{妹妹帮助我学习汉语，我们互相关心，互相帮助。} \\
\zh{Mèimei bāngzhù wǒ xuéxí Hànyǔ, wǒmen hùxiāng guānxīn, hùxiāng bāngzhù.} \\
« Ma petite sœur m'aide à apprendre le chinois ; nous nous soucions les uns des autres, nous nous aidons mutuellement. » \\[0.5em]

\zh{因为家人很团结，所以我觉得家是最温暖的地方。我爱我的家。} \\
\zh{Yīnwèi jiārén hěn tuánjié, suǒyǐ wǒ juéde jiā shì zuì wēnnuǎn de dìfang. Wǒ ài wǒ de jiā.} \\
« Parce que les membres de la famille sont très unis, je pense que la famille est l'endroit le plus chaleureux. J'aime ma famille. »
\end{textebox}

\begin{methode}[Comment aborder un texte chinois en classe de Première]
\begin{enumerate}
\item \textbf{Lecture globale :} Lisez le texte à haute voix (professeur puis élèves) en vous concentrant sur les tons et le rythme (groupes de sens). Ne cherchez pas à tout comprendre tout de suite.
\item \textbf{Repérage des connus :} Surlignez les caractères et mots déjà appris (HSK 1-2 : \zh{我, 家, 有, 人, 爸爸, 妈妈, 妹妹, 是, 老师, 学生, 周末, 吃饭, 看, 电视, 很, 也, 帮助, 学习, 汉语, 我们, 因为, 所以, 觉得, 地方, 爱}).
\item \textbf{Décodage des nouveaux :} Identifiez les 5-6 mots nouveaux ciblés (\zh{口, 四, 安静, 温暖, 互相, 关心, 团结, 最}). Utilisez le contexte, la structure du caractère (clé phonétique/sémantique) ou le dictionnaire.
\item \textbf{Traduction segmentée :} Traduisez phrase par phrase en respectant la syntaxe française, puis proposez une version française naturelle.
\item \textbf{Analyse structurelle :} Repérez les patrons grammaticaux de la leçon (\zh{有}+classif, \zh{互相}+V, \zh{一起}+V, \zh{最}+Adj, \zh{因为}...\zh{所以}...).
\end{enumerate}
\end{methode}

\courssub{Glossaire du thème : Vocabulaire de base (HSK 1-3)}

\begin{center}
\begin{tabularx}{\linewidth}{|X|X|X|X|}\hline
\rowcolor{popblueL} \textbf{Caractères} & \textbf{Pinyin} & \textbf{Nature} & \textbf{Traduction} \\ \hline
家 & jiā & nom & famille, maison, foyer \\ \hline
口 & kǒu & classif. & classificateur pour les membres de la famille \\ \hline
人 & rén & nom & personne, être humain \\ \hline
爸爸 & bàba & nom & papa, père \\ \hline
妈妈 & māma & nom & maman, mère \\ \hline
妹妹 & mèimei & nom & petite sœur \\ \hline
弟弟 & dìdi & nom & petit frère \\ \hline
哥哥 & gēge & nom & grand frère \\ \hline
姐姐 & jiějie & nom & grande sœur \\ \hline
老师 & lǎoshī & nom & professeur, enseignant \\ \hline
医生 & yīshēng & nom & médecin \\ \hline
学生 & xuésheng & nom & élève, étudiant \\ \hline
周末 & zhōumò & nom/temps & week-end \\ \hline
全家 & quánjiā & nom & toute la famille (ensemble) \\ \hline
全家人 & quánjiārén & nom & tous les membres de la famille \\ \hline
一起 & yìqǐ & adv. & ensemble (action commune) \\ \hline
吃饭 & chīfàn & verbe & manger (un repas) \\ \hline
看电视 & kàn diànshì & verbe & regarder la télévision \\ \hline
安静 & ānjìng & adj. & calme, tranquille, silencieux \\ \hline
温暖 & wēnnuǎn & adj. & chaleureux, doux (température ou sentiment) \\ \hline
帮助 & bāngzhù & verbe/nom & aider, aide \\ \hline
学习 & xuéxí & verbe/nom & étudier, apprendre, étude \\ \hline
汉语 & Hànyǔ & nom & langue chinoise (han) \\ \hline
互相 & hùxiāng & adv. & mutuellement, réciproquement \\ \hline
关心 & guānxīn & verbe & se soucier de, prendre soin de \\ \hline
团结 & tuánjié & verbe/adj. & s'unir, unir ; uni, solidaire \\ \hline
因为 & yīnwèi & conj. & parce que (cause) \\ \hline
所以 & suǒyǐ & conj. & donc, par conséquent (conséquence) \\ \hline
觉得 & juéde & verbe & penser que, trouver que (impression) \\ \hline
最 & zuì & adv. & le plus (superlatif) \\ \hline
地方 & dìfang & nom & endroit, lieu, place \\ \hline
爱 & ài & verbe & aimer, amour \\ \hline
\end{tabularx}
\end{center}

\begin{attention}[Faux-amis et pièges fréquents pour francophones]
\begin{itemize}
\item \zh{人} vs \zh{口} : \zh{人} = « personne » (contenu) ; \zh{口} = « bouche » (classificateur pour compter les gens de la famille). On dit \zh{三口人} (trois personnes), jamais \zh{三个人} dans le cadre familial strict.
\item \zh{温暖} vs \zh{暖和} : \zh{温暖} = chaleureux (sentiment, atmosphère, climat doux) ; \zh{暖和} = il fait doux (météo, ressenti physique immédiat). Le texte dit \zh{家里很温暖} (ambiance affective).
\item \zh{互相} (adv) vs \zh{一起} (adv) : \zh{互相} + V simple (action réciproque A$\leftrightarrow$B) ; \zh{一起} + V (action conjointe A+B$\rightarrow$C). Ex: \zh{我们互相帮助} (je t'aide, tu m'aides) vs \zh{我们一起吃饭} (nous mangeons au même moment).
\item \zh{帮助} vs \zh{帮} : \zh{帮助} est plus formel, transitif direct (\zh{帮助我}) ; \zh{帮} est plus oral, souvent \zh{帮 + 的忙} ou \zh{帮 + VV} (\zh{帮我做}).
\item Ordre des mots \zh{因为}...\zh{所以}... : En chinois, les deux conjonctions sont souvent présentes (corrélatives). En français, on choisit « Parce que... » OU « ...donc ». Erreur : \zh{因为...所以...} traduit par « Parce que... donc... » (redondant en français standard).
\end{itemize}
\end{attention}

\courssub{Structures grammaticales clés}

\begin{propriete}[Structure 1 : Existence / Possession avec \zh{有} et le classificateur \zh{口}]
\textbf{Schéma :} Sujet (Possesseur) + \zh{有} + Nombre + \zh{口} + \zh{人}
\end{propriete}

\begin{exemplebox}[Exemples]
\zh{我家有四口人。} (Wǒ jiā yǒu sì kǒu rén.) — Ma famille compte 4 personnes. \\
\zh{他家有五口人。} (Tā jiā yǒu wǔ kǒu rén.) — Sa famille compte 5 personnes. \\
\zh{你家有几口人？} (Nǐ jiā yǒu jǐ kǒu rén ?) — Combien de personnes dans ta famille ?
\end{exemplebox}
\textbf{Règle :} \zh{有} exprime la possession ou l'existence. Pour les membres de la famille, le classificateur obligatoire est \zh{口}. \zh{几} (jǐ) s'utilise pour « combien » (nombre < 10).


\begin{propriete}[Structure 2 : Réciprocité avec \zh{互相} / \zh{彼此}]
\textbf{Schéma :} Sujet (Pluriel) + \zh{互相} / \zh{彼此} + Verbe (souvent psychologique ou action sociale)
\end{propriete}

\begin{exemplebox}[Exemples]
\zh{我们互相帮助。} (Wǒmen hùxiāng bāngzhù.) — Nous nous aidons mutuellement. \\
\zh{他们互相关心。} (Tāmen hùxiāng guānxīn.) — Ils se soucient l'un de l'autre. \\
\zh{同学之间要互相尊重。} (Tóngxué zhījiān yào hùxiāng zūnzhòng.) — Les camarades doivent se respecter mutuellement.
\end{exemplebox}
\textbf{Note :} \zh{互相} et \zh{彼此} sont interchangeables. \zh{互相} est plus courant à l'oral. Le sujet doit être pluriel (nous, vous, ils, ou noms collectifs).


\begin{propriete}[Structure 3 : Action conjointe avec \zh{一起}]
\textbf{Schéma :} Sujet + \zh{一起} + Verbe (+ Complément)
\end{propriete}

\begin{exemplebox}[Exemples]
\zh{全家人一起吃饭。} (Quánjiārén yìqǐ chīfàn.) — Toute la famille mange ensemble. \\
\zh{我们周末一起去公园。} (Wǒmen zhōumò yìqǐ qù gōngyuán.) — Nous allons au parc ensemble le week-end.
\end{exemplebox}
\textbf{Position :} \zh{一起} se place \textbf{avant} le verbe. Ne pas dire \zh{吃饭一起}.


\begin{propriete}[Structure 4 : Superlatif avec \zh{最}]
\textbf{Schéma :} Sujet + \zh{最} + Adjectif (verbe d'état)
\end{propriete}

\begin{exemplebox}[Exemples]
\zh{家是最温暖的地方。} (Jiā shì zuì wēnnuǎn de dìfang.) — La famille est l'endroit le plus chaleureux. \\
\zh{妈妈做的菜最好吃。} (Māma zuò de cài zuì hǎochī.) — La cuisine de maman est la meilleure.
\end{exemplebox}
\textbf{Règle :} Pas de \zh{比} (bǐ) ni de \zh{更} (gèng). \zh{最} marque le degré suprême au sein d'un groupe implicite ou explicite.


\begin{propriete}[Structure 5 : Causalité \zh{因为} ... \zh{所以} ...]
\textbf{Schéma :} \zh{因为} + Cause (phrase), \zh{所以} + Conséquence (phrase)
\end{propriete}

\begin{exemplebox}[Exemples]
\zh{因为下雨，所以我不去了。} (Yīnwèi xiàyǔ, suǒyǐ wǒ bù qù le.) — Parce qu'il pleut, je n'y vais pas. \\
\zh{因为家人很团结，所以我觉得家很温暖。} (Yīnwèi jiārén hěn tuánjié, suǒyǐ wǒ juéde jiā hěn wēnnuǎn.)
\end{exemplebox}
\textbf{Attention :} En chinois écrit formel, les deux sont souvent utilisés ensemble. À l'oral, on peut n'utiliser que \zh{因为} (en tête) ou que \zh{所以} (en seconde position).


\courssub{Clés (radicaux) et ordre des traits}

\begin{definition}[Rappel : Les clés (radicaux) et l'ordre des traits]
Une clé (部首, bùshǒu) est un composant graphique qui permet de classer les caractères dans le dictionnaire et donne souvent une indication sémantique. L'ordre des traits (笔顺, bǐshùn) suit des règles fixes : de haut en bas (上到下), de gauche à droite (左到右), horizontal avant vertical (横竖), extérieur avant intérieur (外到里), fermer en dernier (最后封口).
\end{definition}

\begin{aretenir}[Analyse de 5 caractères clés de la leçon]
\begin{center}
\begin{tabularx}{\linewidth}{|X|X|X|X|}\hline
\rowcolor{popblueL} \textbf{Caractère} & \textbf{Clé (Radical)} & \textbf{Sens de la clé} & \textbf{Nb traits} \\ \hline
爱 & 爫 & griffe / main qui couvre & 10 \\ \hline
家 & 宀 & toit, maison (bǎogàitóu) & 10 \\ \hline
帮 & 巾 & tissu, étoffe (jīnzìpáng) & 10 \\ \hline
团 & 囗 & enceinte, clôture (wéizìkuàng) & 6 \\ \hline
暖 & 日 & soleil (rìzìpáng) & 13 \\ \hline
\end{tabularx}
\end{center}

\textbf{Détail de l'ordre des traits (à pratiquer sur cahier d'écriture) :}
\begin{itemize}
\item \zh{爱} (10 traits) : 1. \zh{爫} (point, horizontal, crochet) — 2. \zh{冖} (horizontal, crochet vertical) — 3. \zh{心} (point, point, crochet vertical, point) — 4. \zh{友} (horizontal, vertical, horizontal, horizontal, horizontal, vertical, horizontal). \textit{Règle : haut $\rightarrow$ bas.}
\item \zh{家} (10 traits) : 1. \zh{宀} (point, point, horizontal+crochet) — 2. \zh{豕} (horizontal, vertical, horizontal, horizontal, horizontal, vertical, horizontal, horizontal). \textit{Règle : toit $\rightarrow$ intérieur.}
\item \zh{帮} (10 traits) : 1. \zh{巾} (vertical, horizontal+crochet, vertical, horizontal) — 2. \zh{邦} (gauche \zh{土} puis droite \zh{攵}). \textit{Règle : gauche $\rightarrow$ droite (clé à gauche).}
\item \zh{团} (6 traits) : 1. \zh{囗} (vertical, horizontal+crochet, horizontal) — 2. \zh{专} (point, horizontal, vertical, horizontal). \textit{Règle : extérieur $\rightarrow$ intérieur $\rightarrow$ fermer.}
\item \zh{暖} (13 traits) : 1. \zh{日} (vertical, horizontal+crochet, horizontal, horizontal) — 2. \zh{爱} sans le \zh{爫} (partie droite complexe : \zh{冖}, \zh{心}, \zh{友}). \textit{Règle : gauche $\rightarrow$ droite.}
\end{itemize}

\end{aretenir}

\begin{savaistu}[Le caractère \zh{爱} : simplification et cœur]
En caractères traditionnels (繁體字), « aimer » s'écrit \zh{愛} : clé \zh{心} (cœur) au milieu, \zh{爪} (griffe) en haut, \zh{友} (ami) en bas $\rightarrow$ « Le cœur va vers l'ami ». En simplifié \zh{爱}, le cœur \zh{心} est gardé, mais le haut est simplifié en \zh{爫} et le bas \zh{友} reste. \textbf{Leçon :} En chinois, aimer implique le cœur (\zh{心}). Sans cœur, ce n'est pas de l'amour.
\end{savaistu}

\courssub{Compréhension et exercices d'application}

\courssub{Exercice 1 — QCM : Le bon sens (Vocabulaire \& Grammaire)}
\begin{exercice}[QCM]
Choisissez la bonne réponse (une seule par question).
\begin{enumerate}
\item Que signifie \zh{我家有四口人} ?
\begin{itemize}
\item[A.] Ma famille a quatre bouches.
\item[B.] Ma famille compte quatre personnes.
\item[C.] Ma famille a quatre portes.
\end{itemize}
\item Complétez : \zh{我们 \_\_\_\_ 帮助。} (Nous nous aidons mutuellement.)
\begin{itemize}
\item[A.] 互相
\item[B.] 一起
\item[C.] 因为
\end{itemize}
\item Dans le mot \zh{家人} (jiārén), que signifie \zh{人} ?
\begin{itemize}
\item[A.] Maison
\item[B.] Famille
\item[C.] Personne
\end{itemize}
\item Traduisez : \zh{家是最温暖的地方。}
\begin{itemize}
\item[A.] La maison est l'endroit le plus froid.
\item[B.] La famille est l'endroit le plus chaleureux.
\item[C.] La famille est un endroit chaud.
\end{itemize}
\end{enumerate}
\end{exercice}

\courssub{Exercice 2 — Vrai ou Faux justifié (Compréhension de texte)}
\begin{exercice}[V/F]
D'après le texte « 我的家 », dites si les affirmations sont \zh{对} (Vrai) ou \zh{不对} (Faux). Justifiez en citant la phrase du texte en chinois et sa traduction.
\begin{enumerate}
\item Le week-end, toute la famille mange ensemble et regarde la télévision.
\item L'auteur n'aime pas sa famille.
\item Le petit frère de l'auteur l'aide à apprendre le chinois.
\item La famille de l'auteur compte cinq personnes.
\end{enumerate}
\end{exercice}

\courssub{Exercice 3 — Vocabulaire : Le mot qui convient (Collocations)}
\begin{exercice}[Choix lexical]
Complétez les phrases avec le mot approprié parmi la liste : \zh{互相, 冷, 走, 团结, 分开, 帮助, 爱, 一起}.
\begin{enumerate}
\item 我们 \_\_\_\_\_\_\_\_\_\_\_\_The user wants me to continue exactly where I stopped in the previous response. Looking at the previous response, it ended mid-sentence in the middle of Exercise 3 (Vocabulaire : Le mot qui convient). The last visible text was:

"1. 我们 \_\_\_\_\_\_\_\_\_\_\_\_
\end{enumerate}
\end{exercice}


\newpage

\coursec{Fiche bilan}

\begin{popfigure}
\begin{tikzpicture}[mindmap, grow cyclic, every node/.style={concept, align=center, text width=2.6cm, font=\small},
root concept/.append style={concept color=popblue, text=white, text width=3.2cm, font=\bfseries},
level 1/.append style={level distance=4.2cm, sibling angle=51},
level 1 concept/.append style={text width=2.6cm}]
\node[root concept]{L'amour de la famille et la solidarité familiale}
child[concept color=poporange]{node[align=center]{Lexique famille\\爸爸 妈妈 哥哥 姐姐 弟弟 妹妹}}
child[concept color=popgreen]{node[align=center]{Sentiments\\爱 喜欢 关心 想念}}
child[concept color=poppurple]{node[align=center]{Solidarité\\帮助 照顾 支持 团结}}
child[concept color=popgold]{node[align=center]{Possession\\的\\我的家 妈妈的爱}}
child[concept color=poppink]{node[align=center]{Fréquence et degré\\常常 很 非常 都}}
child[concept color=popblue!60]{node[align=center]{Compréhension\\lire, repérer, répondre}}
child[concept color=popgreen!60]{node[align=center]{Culture\\famille Cameroun / Chine}};
\end{tikzpicture}
\legende{Carte mentale de la leçon : la famille et la solidarité familiale}
\end{popfigure}

\begin{bilan}
\textbf{1. Lexique clé à connaître par cœur}
\begin{itemize}
\item \zh{家} (jiā) : famille, maison ; \zh{家人} (jiārén) : membres de la famille
\item \zh{爸爸} (bàba) : papa ; \zh{妈妈} (māma) : maman ; \zh{父母} (fùmǔ) : les parents
\item \zh{哥哥} (gēge) : frère aîné ; \zh{姐姐} (jiějie) : sœur aînée ; \zh{弟弟} (dìdi) : frère cadet ; \zh{妹妹} (mèimei) : sœur cadette
\item \zh{爱} (ài) : aimer, amour ; \zh{喜欢} (xǐhuan) : aimer bien ; \zh{想念} (xiǎngniàn) : manquer (qqn)
\item \zh{关心} (guānxīn) : se soucier de ; \zh{帮助} (bāngzhù) : aider ; \zh{照顾} (zhàogù) : prendre soin de
\item \zh{支持} (zhīchí) : soutenir ; \zh{团结} (tuánjié) : solidarité, unité ; \zh{一起} (yìqǐ) : ensemble
\end{itemize}

\textbf{2. Règles de grammaire à connaître par cœur}
\begin{center}
\begin{tabularx}{\linewidth}{|X|X|X|}
\hline
\rowcolor{popblueL} Structure & Exemple (caractères + pinyin) & Traduction \\
\hline
A + 的 + B (possession) & \zh{我妈妈的爱} Wǒ māma de ài & l'amour de ma maman \\
\hline
Sujet + 很 / 非常 + adjectif & \zh{我们家很团结。} Wǒmen jiā hěn tuánjié. & Notre famille est très unie. \\
\hline
Sujet + 常常 + verbe & \zh{哥哥常常帮助我。} Gēge chángcháng bāngzhù wǒ. & Mon frère aîné m'aide souvent. \\
\hline
Sujet + 都 + verbe & \zh{我们都爱我们的家。} Wǒmen dōu ài wǒmen de jiā. & Nous aimons tous notre famille. \\
\hline
Sujet + 跟 / 和 + qqn + 一起 + verbe & \zh{我跟家人一起吃饭。} Wǒ gēn jiārén yìqǐ chīfàn. & Je mange avec ma famille. \\
\hline
\end{tabularx}
\end{center}

\textbf{3. Écriture des caractères (clés et ordre des traits)}
\begin{itemize}
\item \zh{家} : clé du toit 宀 en haut, puis 豕 en bas ; on écrit toujours de haut en bas.
\item \zh{爱} : clé de la main 又 ; on écrit le haut avant le bas, la gauche avant la droite.
\item \zh{妈} : clé de la femme 女 à gauche, puis 马 à droite (gauche avant droite).
\end{itemize}

\textbf{4. Méthode express : comprendre un texte chinois}
\begin{enumerate}
\item Lire le titre et repérer les mots connus (noms de personnes, 家, 爱…).
\item Lire le texte deux fois : d'abord le sens général, ensuite les détails.
\item Souligner les mots-clés de la question, puis retrouver la phrase correspondante dans le texte.
\item Rédiger la réponse avec une structure simple : Sujet + verbe + complément.
\end{enumerate}
\end{bilan}

\begin{aretenir}[Checklist avant le Bac]
\begin{itemize}
\item[$\square$] Je sais nommer tous les membres de la famille en chinois avec le bon pinyin et les bons tons.
\item[$\square$] Je distingue frère aîné \zh{哥哥} et frère cadet \zh{弟弟}, sœur aînée \zh{姐姐} et sœur cadette \zh{妹妹}.
\item[$\square$] J'emploie correctement \zh{的} pour exprimer la possession (\zh{我的家}).
\item[$\square$] Je place l'adverbe de fréquence \zh{常常} avant le verbe, jamais après.
\item[$\square$] Je sais que \zh{很} est presque obligatoire devant un adjectif attribut.
\item[$\square$] J'utilise \zh{都} après le sujet pour dire « tous ».
\item[$\square$] Je sais construire une phrase avec \zh{和 / 跟 … 一起} (faire quelque chose avec quelqu'un).
\item[$\square$] Je connais les verbes de la solidarité : \zh{帮助}, \zh{照顾}, \zh{关心}, \zh{支持}.
\item[$\square$] Je sais écrire \zh{家}, \zh{爱} et \zh{妈} dans le bon ordre des traits.
\item[$\square$] Je sais répondre à une question de compréhension par une phrase complète et correcte.
\item[$\square$] Je peux présenter ma famille en 4 ou 5 phrases chinoises simples.
\item[$\square$] Je sais comparer la famille au Cameroun et en Chine en deux idées claires.
\end{itemize}
\end{aretenir}

\findecours

\end{document}
